% !TEX program = pdflatex
\documentclass[12pt]{article}

\usepackage[utf8]{inputenc}
\usepackage{CJKutf8}
\usepackage{amsmath,amssymb}
\usepackage{mathtools}
\usepackage{amsthm}
\usepackage{geometry}
\geometry{
  a4paper,
  top=25mm,
  bottom=25mm,
  left=20mm,
  right=20mm
}
\usepackage{xcolor}
\usepackage{url}
\usepackage[unicode,colorlinks=true,linkcolor=blue,urlcolor=blue]{hyperref}
\usepackage{xurl}
\Urlmuskip=0mu plus 4mu
\sloppy
\usepackage{array}
\usepackage{tocloft}

\renewcommand{\cftsecleader}{\cftdotfill{\cftdotsep}}
\renewcommand{\refname}{参考文献}

\newtheorem{definition}{定义}[section]
\newtheorem{axiom}{公理}[section]
\newtheorem{assumption}{假设}[section]
\newtheorem{theorem}{定理}[section]
\newtheorem{proposition}{命题}[section]
\newtheorem{corollary}{推论}[section]
\newtheorem{lemma}{引理}[section]
\newtheorem{remark}{备注}[section]
\newtheorem{Certificate}{证书}[section]

\newcommand{\NN}{\mathbb{N}}
\newcommand{\RR}{\mathbb{R}}
\newcommand{\GFramework}{\textsf{G-Framework}}
\newcommand{\MicroTunnel}{\mathsf{MicroTunnel}}
\newcommand{\CrossWall}{\mathsf{CrossWall}}
\newcommand{\Path}{\mathcal P}
\newcommand{\DefEq}{\coloneqq}
\newcommand{\TV}{\mathrm{TV}}
\newcommand{\Unif}{\mathrm{unif}}
\newcommand{\Normalize}{\operatorname{Normalize}}
\newcommand{\id}{\mathrm{id}}
\newcommand{\incl}{\iota}
\newcommand{\code}[1]{\path{#1}}

\setlength{\emergencystretch}{8em}
\setlength{\parindent}{0pt}
\setlength{\parskip}{0.6em}

\begin{document}
\begin{CJK*}{UTF8}{gkai}

\title{基于高阶正交确定性收敛的统计力学同伦命中：PDE 求解范式替代\\（工作笔记：形式系统版）}
\author{Gao Zheng（高政）}
\date{2026-05-13}
\maketitle

\section*{前言}
\addcontentsline{toc}{section}{前言}

本文的目标不是另写一套孤立的 PDE 求解术语，而是在 \GFramework{} 的既有谱系中，
把障碍、修复、同伦、命中、统计力学、正交残差、雅可比间隙、确定性收敛与
PDE 残差解组织成一条可追踪的形式链条。若这些概念分散存在，它们容易只表现为
互相呼应的解释性母题；本文将它们压合为一个可定义、可优化、可判定、可收敛的
统计力学同伦命中范式。

本文的第一层立场是求解对象的迁移。传统 PDE 求解通常以方程本身为中心：
给定微分算子，进行离散化，求解代数系统，再恢复函数近似。本文则把求解中心从
“显式求解方程”迁移到“命中低残差集合”。这并不是削弱 PDE 的数学地位，而是将
PDE 在求解层中的直接反演角色降级为残差能量的语义来源。由此，核心问题变为：
如何构造统计力学分布律、同伦扭曲轨迹与命中泛函，使系统以可控方式命中低残差集合。
高阶 PDE 中常见的算子反演、边界耦合、高阶导数病态与局部线性化敏感性，因而被
转写为残差能量、命中集合、轨迹控制与雅可比间隙下降的问题。

本文的第二层立场是随机过程的确定化。本文严格区分分布律、随机变量、命中事件、
命中指示变量、首次命中时间、命中泛函、生存概率与命中指数。命中不是抽象地
“趋向某个分布”，而是由分布律索引的随机变量取值落入指定命中集合。单个样本仍然
保持随机性；趋于确定的是经验测度、经验命中泛函、经验命中指数、经验雅可比间隙
以及由它们诱导的轨迹函数。由此，统计力学采样不再停留在启发式层面，而被纳入
有限样本、有限截断下的工程确定性下降系统，并在一致收敛、投影稳定与雅可比稳定
条件下进入数学严格的确定性极限。

本文的第三层立场是微观隧穿的可控化。第~\ref{sec:tex55-main} 节以墙、孔洞、
正厚度隧道与穿墙芯线为几何原型，证明正厚度孔洞与穿墙芯线在同一端口粘合下
具有同伦等价关系。它的作用并不止于一个拓扑事实，而是为后文提供“障碍穿透可被
收缩为可控通路”的几何锚点。后续章节将这一思想从拓扑空间推广到算法空间：
复杂的随机搜索过程被表达为命中轨迹、正交方向与雅可比间隙的控制问题。于是，
障碍修复不再只是直观比喻，而成为贯穿拓扑、统计、优化与 PDE 残差求解的统一操作。

本文的第四层立场是变分算子的扩展。传统变分框架通常围绕能量泛函、一阶变分、
Euler--Lagrange 方程与二阶变分结构展开；本文把变分对象从函数空间上的能量极值
提升为轨迹函数空间上的命中指数极值。换言之，变分不再只问“哪个函数使能量下降”，
而进一步问“哪条同伦轨迹使未命中生存概率下降”。因此，变分算子在本文中被扩展为
残差能量变分、命中指数变分、正交残差变分与雅可比间隙变分的组合链条。

本文的第五层立场是高阶正交机制的系统化。所谓“高阶”并不只是增加复杂度，而是
表示剩余命中增益的递归提取。第 \(k\) 阶方向不是重复前序方向，而是在目标扭曲与
已用方向张成空间的正交补中提取命中梯度残差。每一阶都回答同一个问题：当前系统
是否仍存在未被利用的命中方向；若存在，则通过正交投影取出；若不存在，则说明该
正交补中的一阶命中增益已经耗尽。

本文的第六层立场是雅可比间隙的工程实现作用。理论上知道理想命中方向，并不表示
参数化同伦系统能够实际产生该方向。本文以雅可比间隙衡量实际同伦扭曲雅可比方向
与理想高阶正交命中方向之间的差距。这个间隙不是附属量，而是每阶扭曲是否真正
有效的判据：只有当实际可实现方向足够逼近理想方向，并且不泄漏回已用方向空间时，
同伦扭曲才获得明确的收敛指引。

本文还刻意区分问题转化与极限保真。问题转化本身是构造层命题：给定 PDE 后，只要
能够定义离散状态、残差能量、命中集合、统计力学分布律与命中泛函，便已经完成从
显式 PDE 求解到低残差命中的范式转化。至于命中所得对象能否在 \(N\to\infty\)、
\(L\to\infty\)、\(\varepsilon\to0\) 的极限中恢复原 PDE 解，则属于保真性定理，
需要残差相容、紧性、闭性、一致收敛与雅可比稳定性等证书。本文通过这种层级区分，
避免把“可转化”误写成“无条件等价”，也避免把“需要保真条件”误写成“转化本身未成立”。

因此，本文在整体 \GFramework{} 中承担的是范式锚定作用。它向前连接障碍穿透、
空洞修复与穿墙芯线的同伦几何，向中间连接统计力学分布律、随机变量簇与命中泛函，
向后连接高阶正交、雅可比间隙、工程截断确定性与 PDE 解恢复。本文所说的“替代”
并不是删除 PDE 理论，而是将 PDE 从求解层中心移动到残差语义层与保真性证书层；
也不是放弃确定性，而是把随机命中通过经验泛函和雅可比间隙下降转化为可检验的
确定性极限系统。

概括地说，本文把 \GFramework{} 的核心精神从“发现结构”推进到“控制结构”，从
“描述障碍”推进到“穿透障碍”，从“形式化问题”推进到“构造可收敛的求解范式”。
它提供的不是单个算法，而是一条由低残差命中、同伦轨迹、高阶正交残差、雅可比间隙
与确定性极限共同构成的方法论主干。

\setcounter{tocdepth}{3}
\setcounter{secnumdepth}{3}
\tableofcontents
\clearpage

%%%%%%%%%%%%%%%%%%%%%%%%%%%%%%%%%%%%%%%%%%%%%%%%%%%%%%%%%%%%
\section{形式逻辑：正厚度孔洞与穿墙芯线的同伦等价}
\label{sec:tex55-main}
%%%%%%%%%%%%%%%%%%%%%%%%%%%%%%%%%%%%%%%%%%%%%%%%%%%%%%%%%%%%

\begin{remark}[核心陈述]
本文把“墙中打通正厚度孔洞”和“保留穿墙芯线通路”写成同一个障碍修复
\(\rho\) 下的两个可达空间：
\[
  \MicroTunnel_{\rho}(W),
  \qquad
  \CrossWall_{\rho}(W).
\]
要证明的形式化命题是
\[
  \boxed{
  \MicroTunnel_{\rho}(W)
  \simeq
  \CrossWall_{\rho}(W)
  }.
\]
其中 \(\MicroTunnel_{\rho}(W)\) 表示墙 \(W\) 被修复后得到的正厚度隧道可达空间；
\(\CrossWall_{\rho}(W)\) 表示在同一修复端口下，仅保留穿墙芯线与两端端口的可达空间。
结论是：只要修复确实打通拓扑通道，正厚度孔洞与其穿墙芯线在同伦意义下等价。
\end{remark}

\begin{remark}[阅读导航与引用口径]
本文的写作体例、证书格式与“公理降级为可证条目”的组织方式，对齐
\cite{RepoTex37SemanticGauge,RepoTex51TruthAnchor} 的形式系统稿。
拓扑语言只使用代数拓扑中的标准概念：强变形收缩、同伦等价、商空间粘合与路径空间，
基础背景可参考 \cite{Hatcher2002AT}。
仓内概念谱系方面，雅可比间隙与同伦扭曲势能路径参考
\cite{RepoTex28JacGap}；
微观隧穿命中、PDE 表象升维与障碍修复判据参考
\cite{RepoTex45ObstructionGlobalSection}；
统计格、同伦隧穿与 PDE 策略簇参考
\cite{RepoTex53StrategyCluster}；
变分同伦修复塔与 \(L_\infty\) 无故障修复链条参考
\cite{RepoTex17VariationalTower,RepoTex50LinftyRepair}。
本文不声称度量代价、概率权重或能量函数的改变自动等价于拓扑隧道；
它只证明在\textbf{同一拓扑端口粘合}下，厚通道与芯线通道的同伦型一致。
\end{remark}

\subsection{对象域与目标命题}
\label{subsec:tex55-objects}

\begin{definition}[墙与未修复可达空间]
\label{def:tex55-wall}
令 \(M\) 为环境拓扑空间，令 \(W\subset M\) 为障碍。
若 \(W\) 在局部把可达区域分成两侧，并且存在指定的穿越方向，则称 \(W\) 为墙。
未修复前的可达空间记为
\[
  X_0\DefEq M\setminus W.
\]
墙的两侧可达端口区域记为
\[
  X_-,X_+\subset X_0.
\]
这里的 \(X_-\) 与 \(X_+\) 只记录待粘合端口所在的两侧区域，不要求它们穷尽 \(X_0\)。
\end{definition}

\begin{definition}[障碍修复与正厚度通道]
\label{def:tex55-repair}
障碍修复 \(\rho\) 指在墙 \(W\) 中打通一个正厚度通道，并指定两个端口粘合映射。
对某个 \(k\in\NN\cup\{0\}\)，令
\[
  D^k\DefEq\{x\in\RR^k:\|x\|\le 1\},
  \qquad
  I\DefEq[0,1],
\]
并记正厚度通道为
\[
  T_k\DefEq D^k\times I.
\]
修复数据包含两个端口嵌入
\[
  \alpha_-:D^k\times\{0\}\longrightarrow X_-,
  \qquad
  \alpha_+:D^k\times\{1\}\longrightarrow X_+.
\]
它们把入口端口与出口端口分别粘到墙的两侧。
\end{definition}

\begin{definition}[微观隧穿空间]
\label{def:tex55-microtunnel}
在修复 \(\rho=(T_k,\alpha_-,\alpha_+)\) 下，完整的正厚度隧道可达空间定义为商空间
\[
  X_{\mu}
  \DefEq
  (X_0\sqcup T_k)/{\sim_{\rho}},
\]
其中等价关系 \(\sim_{\rho}\) 由
\[
  (x,0)\sim_{\rho}\alpha_-(x,0),
  \qquad
  (x,1)\sim_{\rho}\alpha_+(x,1)
  \quad (x\in D^k)
\]
生成。
定义
\[
  \MicroTunnel_{\rho}(W)\DefEq X_{\mu}.
\]
\end{definition}

\begin{definition}[穿墙芯线空间]
\label{def:tex55-crosswall}
定义隧道的穿墙芯线加端口为
\[
  C_k
  \DefEq
  (D^k\times\{0\})
  \cup
  (\{0\}\times I)
  \cup
  (D^k\times\{1\})
  \subset T_k.
\]
在同一端口粘合数据 \(\rho\) 下，穿墙芯线可达空间定义为
\[
  X_{\times}
  \DefEq
  (X_0\sqcup C_k)/{\sim_{\rho}},
\]
其中 \(\sim_{\rho}\) 是定义~\ref{def:tex55-microtunnel} 的端口粘合关系在
\(X_0\sqcup C_k\) 上的限制。
定义
\[
  \CrossWall_{\rho}(W)\DefEq X_{\times}.
\]
由于 \(C_k\subset T_k\)，自然包含映射诱导
\[
  \incl:X_{\times}\hookrightarrow X_{\mu}.
\]
\end{definition}

\subsection{公理降级：强变形收缩判据}
\label{subsec:tex55-deformation-criterion}

\begin{axiom}[公理（降级为定理）：强变形收缩推出同伦等价]
\label{ax:tex55-sdr-homotopy-equivalence}
设 \(A\subset X\)。
若存在连续映射
\[
  r:X\to A
\]
和同伦
\[
  H:X\times I\to X
\]
满足
\[
  H(x,0)=x,
  \qquad
  H(x,1)=\incl r(x),
\]
并且
\[
  \forall a\in A,\quad H(a,s)=a,
\]
则 \(A\) 是 \(X\) 的强变形收缩核。
此时
\[
  X\simeq A.
\]
\end{axiom}

\begin{Certificate}[证书：同伦逆证书]
\label{cert:tex55-sdr-homotopy-equivalence}
证书由包含映射、收缩映射与固定同伦组成：
\[
  \pi_{\mathrm{sdr}}
  =
  \bigl(\incl:A\hookrightarrow X,\ r:X\to A,\ H:X\times I\to X\bigr).
\]
其验证条件为
\[
  r\circ\incl=\id_A,
  \qquad
  \incl\circ r\simeq\id_X.
\]
\end{Certificate}

\begin{proof}[证明（基于数学符号推演逻辑的谓词因果链）]
由强变形收缩条件，
\[
  H(x,0)=x
  \Longrightarrow
  H_0=\id_X.
\]
又由
\[
  H(x,1)=\incl r(x)
  \Longrightarrow
  H_1=\incl\circ r.
\]
因此
\[
  \id_X\simeq \incl\circ r.
\]
对任意 \(a\in A\)，固定条件给出
\[
  H(a,1)=a.
\]
另一方面
\[
  H(a,1)=\incl r(a),
\]
于是
\[
  r(\incl(a))=a,
  \qquad
  r\circ\incl=\id_A.
\]
所以 \(\incl\) 与 \(r\) 互为同伦逆，得到
\[
  X\simeq A.
\]
证毕。
\end{proof}

\subsection{引理：正厚度隧道收缩到穿墙芯线加端口}
\label{subsec:tex55-tunnel-spine-lemma}

\begin{lemma}[正厚度隧道的芯线强变形收缩]
\label{lem:tex55-tunnel-spine-sdr}
对任意 \(k\in\NN\cup\{0\}\)，有
\[
  D^k\times I\searrow C_k.
\]
也就是说，\(C_k\) 是 \(T_k=D^k\times I\) 的强变形收缩核。
\end{lemma}

\begin{Certificate}[数学归纳法证书：维度无关的径向收缩模板]
\label{cert:tex55-tunnel-spine-induction}
定义谓词
\[
  P(k):\quad D^k\times I\text{ 强变形收缩到 }C_k.
\]
基步为
\[
  k=0.
\]
此时 \(D^0=\{0\}\)，从而
\[
  D^0\times I=\{0\}\times I=C_0,
\]
故 \(P(0)\) 成立。

归纳步采用同一个维度无关模板。
对 \(x\in D^{k+1}\)，令
\[
  r=\|x\|\in[0,1].
\]
下文构造的收缩公式只依赖 \(r\)、端点 \(0,1\) 与线性凸组合，不依赖坐标个数。
因此该公式在 \(D^{k+1}\times I\) 上仍给出合法强变形收缩，故
\[
  P(k)\Longrightarrow P(k+1).
\]
于是
\[
  \forall k\in\NN\cup\{0\},\quad P(k).
\]
实际证明中将直接给出统一公式；归纳证书用于记录该构造随维度升高保持不变。
\end{Certificate}

\begin{proof}[证明（基于数学符号推演逻辑的谓词因果链）]
令
\[
  T_k=D^k\times I.
\]
取任意点
\[
  (x,t)\in D^k\times I,
\]
并记
\[
  r=\|x\|.
\]
定义收缩映射
\[
  q:T_k\to C_k.
\]
若 \(r=0\)，令
\[
  q(0,t)=(0,t).
\]
若 \(r>0\)，按三段定义：
\[
  q(x,t)=
  \begin{cases}
    \left(\left(1-\dfrac{3t}{r}\right)x,\,0\right),
    & 0\le t\le \dfrac r3,\\[1.2em]
    \left(0,\,\dfrac{t-\frac r3}{1-\frac{2r}{3}}\right),
    & \dfrac r3\le t\le 1-\dfrac r3,\\[1.2em]
    \left(\left(1-\dfrac{3(1-t)}{r}\right)x,\,1\right),
    & 1-\dfrac r3\le t\le 1.
  \end{cases}
\]
由于 \(0<r\le 1\)，中间段分母满足
\[
  1-\frac{2r}{3}\ge \frac13>0.
\]
第一段落在下端口：
\[
  q(x,t)\in D^k\times\{0\}\subset C_k.
\]
第二段落在芯线：
\[
  q(x,t)\in\{0\}\times I\subset C_k.
\]
第三段落在上端口：
\[
  q(x,t)\in D^k\times\{1\}\subset C_k.
\]
故
\[
  q(T_k)\subset C_k.
\]

现在验证 \(q\) 固定 \(C_k\)。
若 \(t=0\)，则
\[
  q(x,0)=(x,0).
\]
若 \(t=1\)，则
\[
  q(x,1)=(x,1).
\]
若 \(x=0\)，则按定义
\[
  q(0,t)=(0,t).
\]
因此
\[
  \forall z\in C_k,\quad q(z)=z.
\]

再验证连续性。
在分段边界 \(t=r/3\) 处，第一段与第二段均给出
\[
  q(x,t)=(0,0).
\]
在分段边界 \(t=1-r/3\) 处，第二段与第三段均给出
\[
  q(x,t)=(0,1).
\]
当 \(r\to0\) 时，第一段只能靠近 \(t=0\)，第三段只能靠近 \(t=1\)，中间段趋向
\[
  (0,t).
\]
这与 \(q(0,t)=(0,t)\) 一致。
故 \(q\) 连续。

定义同伦
\[
  H:T_k\times I\to T_k
\]
为
\[
  H((x,t),s)
  \DefEq
  (1-s)(x,t)+s\,q(x,t).
\]
由于 \(D^k\) 与 \(I\) 均为凸集，
\[
  H((x,t),s)\in D^k\times I.
\]
并且
\[
  H((x,t),0)=(x,t),
  \qquad
  H((x,t),1)=q(x,t).
\]
若 \((x,t)\in C_k\)，则 \(q(x,t)=(x,t)\)，因此
\[
  H((x,t),s)=(x,t).
\]
于是 \(H\) 是 \(T_k\) 到 \(C_k\) 的强变形收缩。
故
\[
  D^k\times I\searrow C_k.
\]
引理得证。
\end{proof}

\subsection{命题：修复后空间的相对强变形收缩}
\label{subsec:tex55-repaired-space-sdr}

\begin{proposition}[微观隧穿空间强变形收缩到穿墙芯线空间]
\label{prop:tex55-micro-sdr-cross}
在定义~\ref{def:tex55-wall}--\ref{def:tex55-crosswall} 下，
\[
  X_{\mu}\searrow X_{\times}.
\]
\end{proposition}

\begin{Certificate}[证书：端口固定的商空间下降]
\label{cert:tex55-quotient-descent}
证书由三部分组成：
\[
\begin{aligned}
  Q_1&:\quad T_k\searrow C_k;\\
  Q_2&:\quad H_s|_{D^k\times\{0\}}=\id,\quad
  H_s|_{D^k\times\{1\}}=\id;\\
  Q_3&:\quad R|_{X_0}=\id_{X_0},\quad
  \overline H_s|_{X_0}=\id_{X_0}.
\end{aligned}
\]
其中 \(Q_2\) 保证收缩同伦尊重端口粘合关系，因此可下降到商空间。
\end{Certificate}

\begin{proof}[证明（基于数学符号推演逻辑的谓词因果链）]
由定义，
\[
  X_{\mu}=(X_0\sqcup T_k)/{\sim_{\rho}},
  \qquad
  X_{\times}=(X_0\sqcup C_k)/{\sim_{\rho}}.
\]
由引理~\ref{lem:tex55-tunnel-spine-sdr}，存在收缩映射
\[
  q:T_k\to C_k
\]
与同伦
\[
  H:T_k\times I\to T_k
\]
满足
\[
  H_0=\id_{T_k},
  \qquad
  H_1=\incl q,
  \qquad
  H_s|_{C_k}=\id_{C_k}.
\]
因为
\[
  D^k\times\{0\}\subset C_k,
  \qquad
  D^k\times\{1\}\subset C_k,
\]
所以
\[
  q|_{D^k\times\{0\}}=\id,
  \qquad
  q|_{D^k\times\{1\}}=\id,
\]
并且
\[
  H_s|_{D^k\times\{0\}}=\id,
  \qquad
  H_s|_{D^k\times\{1\}}=\id.
\]
这说明 \(q\) 与 \(H_s\) 固定所有被 \(\alpha_-\) 与 \(\alpha_+\) 粘合的端口点。

在 \(X_0\) 上定义恒等映射，并在 \(T_k\) 上使用 \(q\)，得到商空间上的连续映射
\[
  R:X_{\mu}\to X_{\times}.
\]
同理，在 \(X_0\) 上取恒等同伦，在 \(T_k\) 上取 \(H\)，得到商空间上的同伦
\[
  \overline H:X_{\mu}\times I\to X_{\mu}.
\]
它们满足
\[
  \overline H_0=\id_{X_{\mu}},
  \qquad
  \overline H_1=\incl\circ R,
\]
并且
\[
  \overline H_s|_{X_{\times}}=\id_{X_{\times}}.
\]
同时
\[
  R\circ\incl=\id_{X_{\times}}.
\]
因此 \(X_{\times}\) 是 \(X_{\mu}\) 的强变形收缩核，即
\[
  X_{\mu}\searrow X_{\times}.
\]
命题得证。
\end{proof}

\subsection{定理：微观隧穿等价于穿墙而过}
\label{subsec:tex55-main-theorem}

\begin{theorem}[微观隧穿空间与穿墙芯线空间同伦等价]
\label{thm:tex55-micro-cross-homotopy-equivalence}
若障碍 \(W\) 被视为墙，且障碍修复 \(\rho\) 在墙中打通一个拓扑通道，则
\[
  \MicroTunnel_{\rho}(W)
  \simeq
  \CrossWall_{\rho}(W).
\]
\end{theorem}

\begin{Certificate}[证书：强变形收缩到同伦等价]
\label{cert:tex55-main-theorem}
证书为
\[
  \pi_{\mathrm{main}}
  =
  \bigl(
    X_{\mu}=\MicroTunnel_{\rho}(W),\
    X_{\times}=\CrossWall_{\rho}(W),\
    X_{\mu}\searrow X_{\times}
  \bigr).
\]
其中第三项由命题~\ref{prop:tex55-micro-sdr-cross} 签发。
\end{Certificate}

\begin{proof}[证明（基于数学符号推演逻辑的谓词因果链）]
由定义~\ref{def:tex55-microtunnel} 与定义~\ref{def:tex55-crosswall}，
\[
  \MicroTunnel_{\rho}(W)=X_{\mu},
  \qquad
  \CrossWall_{\rho}(W)=X_{\times}.
\]
由命题~\ref{prop:tex55-micro-sdr-cross}，
\[
  X_{\mu}\searrow X_{\times}.
\]
由公理（降级为定理）~\ref{ax:tex55-sdr-homotopy-equivalence}，
\[
  X_{\mu}\simeq X_{\times}.
\]
代回记号，得到
\[
  \MicroTunnel_{\rho}(W)
  \simeq
  \CrossWall_{\rho}(W).
\]
证毕。
\end{proof}

\subsection{推论：路径空间版本}
\label{subsec:tex55-path-space-corollary}

\begin{corollary}[穿墙路径空间的同伦等价]
\label{cor:tex55-path-space}
取两点
\[
  a\in X_-,
  \qquad
  b\in X_+,
\]
并假设 \(a,b\) 不在端口粘合的内部歧义点上，或已在商空间中选定对应点。
定义端点固定路径空间
\[
  \Path_{\mu}(a,b)
  \DefEq
  \{\gamma:I\to X_{\mu}\mid \gamma(0)=a,\ \gamma(1)=b\},
\]
\[
  \Path_{\times}(a,b)
  \DefEq
  \{\gamma:I\to X_{\times}\mid \gamma(0)=a,\ \gamma(1)=b\}.
\]
则
\[
  \Path_{\mu}(a,b)\simeq \Path_{\times}(a,b).
\]
\end{corollary}

\begin{Certificate}[证书：端点固定的路径空间诱导]
\label{cert:tex55-path-space}
证书由收缩映射与包含映射诱导：
\[
  R_{\#}:\Path_{\mu}(a,b)\to\Path_{\times}(a,b),
  \qquad
  R_{\#}(\gamma)=R\circ\gamma,
\]
\[
  \incl_{\#}:\Path_{\times}(a,b)\to\Path_{\mu}(a,b),
  \qquad
  \incl_{\#}(\eta)=\incl\circ\eta.
\]
由于命题~\ref{prop:tex55-micro-sdr-cross} 的收缩在 \(X_0\) 上为恒等，
端点 \(a,b\) 被固定。
\end{Certificate}

\begin{proof}[证明（基于数学符号推演逻辑的谓词因果链）]
由命题~\ref{prop:tex55-micro-sdr-cross}，存在
\[
  R:X_{\mu}\to X_{\times}
\]
和相对 \(X_{\times}\) 固定的同伦
\[
  \overline H:X_{\mu}\times I\to X_{\mu}.
\]
并且
\[
  R|_{X_0}=\id_{X_0},
  \qquad
  \overline H_s|_{X_0}=\id_{X_0}.
\]
因此 \(a,b\in X_0\) 的端点条件被 \(R\) 与 \(\overline H\) 保持。

对任意 \(\eta\in\Path_{\times}(a,b)\)，有
\[
  (R_{\#}\circ\incl_{\#})(\eta)
  =
  R\circ\incl\circ\eta
  =
  \eta,
\]
故
\[
  R_{\#}\circ\incl_{\#}
  =
  \id_{\Path_{\times}(a,b)}.
\]
对任意 \(\gamma\in\Path_{\mu}(a,b)\)，定义路径空间同伦
\[
  \mathcal H(\gamma,s)(t)
  \DefEq
  \overline H(\gamma(t),s).
\]
端点被固定，所以
\[
  \mathcal H(\gamma,s)(0)=a,
  \qquad
  \mathcal H(\gamma,s)(1)=b.
\]
并且
\[
  \mathcal H(\gamma,0)=\gamma,
  \qquad
  \mathcal H(\gamma,1)=\incl_{\#}R_{\#}(\gamma).
\]
于是
\[
  \incl_{\#}\circ R_{\#}
  \simeq
  \id_{\Path_{\mu}(a,b)}.
\]
所以
\[
  \Path_{\mu}(a,b)\simeq \Path_{\times}(a,b).
\]
推论得证。
\end{proof}

\subsection{谓词因果链与备注}
\label{subsec:tex55-predicate-chain}

核心谓词因果链为：
\[
\begin{aligned}
  &W\text{ 是墙}\\
  &\Longrightarrow
  W\text{ 有两侧端口区域 }X_-,X_+\\
  &\Longrightarrow
  \rho\text{ 在 }X_-\text{ 与 }X_+\text{ 之间加入 }T_k=D^k\times I\\
  &\Longrightarrow
  \MicroTunnel_{\rho}(W)=X_0\cup_{\rho}T_k\\
  &\Longrightarrow
  T_k\searrow C_k\\
  &\Longrightarrow
  X_0\cup_{\rho}T_k\searrow X_0\cup_{\rho}C_k\\
  &\Longrightarrow
  \MicroTunnel_{\rho}(W)\simeq\CrossWall_{\rho}(W).
\end{aligned}
\]
一句话压缩为
\[
  \boxed{
  \text{孔洞打通的正厚度隧道}
  \simeq
  \text{穿墙芯线通路}
  }.
\]

\begin{remark}[拓扑修复的边界条件]
上述证明成立的关键条件是：所谓障碍修复 \(\rho\) 确实改变了可达空间的拓扑结构，
也就是把原来被墙阻断的两侧通过端口粘合与通道 \(T_k\) 连接起来。
若修复只改变代价、概率、度量、能量函数或局部权重，而没有生成同一端口粘合下的拓扑通道，
则不能推出
\[
  \MicroTunnel_{\rho}(W)\simeq\CrossWall_{\rho}(W).
\]
此时命题需要先重新定义为度量、概率或动力系统问题，不能直接套用本文的同伦等价结论。
\end{remark}

\clearpage

%%%%%%%%%%%%%%%%%%%%%%%%%%%%%%%%%%%%%%%%%%%%%%%%%%%%%%%%%%%%
\section{形式逻辑：时域局部遍历、频域共振随机变量簇与近似全局最优命中}
\label{sec:tex55-time-frequency-global-hit}
%%%%%%%%%%%%%%%%%%%%%%%%%%%%%%%%%%%%%%%%%%%%%%%%%%%%%%%%%%%%

\begin{remark}[核心陈述]
本章把“时域局部遍历算法经谱展开与截断转化为频域共振分布律，并由该分布律生成每次命中判定的随机变量簇”
严格化为一条近似桥接命题：
\[
  \boxed{
  \text{时域局部遍历算法}
  \xrightarrow{\text{谱展开/截断}}
  \text{频域共振随机变量簇}
  \xrightarrow{\text{统计力学命中}}
  \text{近似全局最优}
  }.
\]
必须先区分两个命题：
\[
  \text{全局可达性}\neq \text{必然命中全局最优}.
\]
全局可达性只表示
\[
  \forall x,y\in\Omega,\quad \exists t\ge0,\quad P^t(x,y)>0,
\]
而全局最优命中要求某个采样或退火过程满足
\[
  \Pr(X_t\in S_*)\to 1.
\]
因此，本文可证明的严格版本不是“任何局部最优算法都会自动变成全局最优算法”，而是：
\[
  \boxed{
  \begin{gathered}
  \text{在可达性、谱可表达、低温集中、收敛近似与采样充分条件下，}\\
  \text{频域共振随机变量簇给出全局最优命中的可控近似。}
  \end{gathered}
  }
\]
\end{remark}

\begin{remark}[阅读导航与引用口径]
本章的马尔可夫链、可逆核、谱展开与全变差距离使用有限状态空间上的标准理论，
参考 \cite{LevinPeresWilmer2017Markov,Norris1998MarkovChains}；
Gibbs/Boltzmann 权重与退火式命中直觉参考统计力学与模拟退火经典文献
\cite{Kirkpatrick1983Optimization,Hajek1988Cooling}。
与前一章的拓扑修复类似，本章只签发带条件的桥接结论：
频域共振在严格对象层级上是“谱模态重加权的分布律 + 由该分布律索引的命中随机变量簇”，不是单个确定性分布函数，更不是无条件的全局最优保证。
\end{remark}

\subsection{对象域：状态空间、局部遍历与全局命中}
\label{subsec:tex55-bridge-objects}

\begin{definition}[状态空间、能量函数与全局最优集]
\label{def:tex55-state-energy}
令
\[
  \Omega=\{x_1,x_2,\ldots,x_N\}
\]
为有限状态空间。
令
\[
  E:\Omega\to\RR
\]
为目标函数，也称能量函数。
定义全局最优值与全局最优集合为
\[
  E_*\DefEq\min_{x\in\Omega}E(x),
  \qquad
  S_*\DefEq\{x\in\Omega:E(x)=E_*\}.
\]
优化目标不是抽象地“变好”，而是以高概率命中集合 \(S_*\)。
\end{definition}

\begin{definition}[时域局部遍历算法]
\label{def:tex55-local-markov}
时域局部算法表示为马尔可夫分布迭代
\[
  \mu_{t+1}=\mu_tP,
\]
其中
\[
  P(x,y)=\Pr(X_{t+1}=y\mid X_t=x)
\]
是从状态 \(x\) 跳到状态 \(y\) 的转移概率。
若 \(P(x,y)>0\) 主要发生在局部邻域
\[
  y\in N(x),
\]
则称 \(P\) 为局部转移核。
对应的样本路径可写为
\[
  x_{t+1}=A(x_t,\xi_t),
\]
其中 \(\xi_t\) 是随机扰动。
\end{definition}

\begin{definition}[全局可达性与全局命中]
\label{def:tex55-reachability-hit}
称转移核 \(P\) 具有全局可达性，若
\[
  \forall x,y\in\Omega,\quad \exists t\ge0,\quad P^t(x,y)>0.
\]
等价地说，状态图是不可约的。
称一族分布或算法实现全局最优命中，若
\[
  \Pr(X_t\in S_*)\to1.
\]
两者的逻辑强度不同：
\[
  \text{可达性只给出存在正概率路径，命中性要求概率质量集中到 }S_*.
\]
\end{definition}

\begin{definition}[局部最优陷阱]
\label{def:tex55-local-trap}
若存在
\[
  x_{\mathrm{loc}}\notin S_*
\]
满足
\[
  E(x_{\mathrm{loc}})\le E(y),
  \qquad
  \forall y\in N(x_{\mathrm{loc}}),
\]
则称 \(x_{\mathrm{loc}}\) 为局部最优点。
若算法只接受下降方向
\[
  E(y)<E(x),
\]
则一旦进入 \(x_{\mathrm{loc}}\)，就可能无法离开。
因此
\[
  \text{纯局部下降}
  \not\Rightarrow
  \text{全局命中}.
\]
\end{definition}

\begin{lemma}[全局可达性不推出必然命中全局最优]
\label{lem:tex55-reachable-not-hit}
存在有限状态空间、能量函数与不可约非周期马尔可夫核 \(P\)，使得
\[
  \forall x,y\in\Omega,\quad \exists t,\quad P^t(x,y)>0,
\]
但
\[
  \Pr(X_t\in S_*)\not\to1.
\]
\end{lemma}

\begin{Certificate}[证书：不可约链的平稳质量反例]
\label{cert:tex55-reachable-not-hit}
取任意有限集合 \(\Omega\) 与非空真子集 \(S_*\subsetneq\Omega\)。
令 \(P\) 为懒惰完全图随机游走：
\[
  P(x,y)=
  \begin{cases}
    \dfrac12, & y=x,\\[0.6em]
    \dfrac{1}{2(N-1)}, & y\ne x.
  \end{cases}
\]
该链不可约、非周期，平稳分布为均匀分布
\[
  \pi(x)=\frac1N.
\]
\end{Certificate}

\begin{proof}[证明（基于数学符号推演逻辑的谓词因果链）]
由证书构造，任意 \(x,y\in\Omega\) 都有
\[
  P(x,y)>0
\]
或在一步内保持原点，因此全局可达性成立。
该链有限、不可约、非周期，所以
\[
  \mu_t\to\pi.
\]
于是
\[
  \Pr(X_t\in S_*)=\mu_t(S_*)\to\pi(S_*)=\frac{|S_*|}{|\Omega|}.
\]
由于 \(S_*\subsetneq\Omega\)，
\[
  \frac{|S_*|}{|\Omega|}<1.
\]
故
\[
  \Pr(X_t\in S_*)\not\to1.
\]
这证明全局可达性不能替代全局最优命中。证毕。
\end{proof}

\subsection{统计力学目标分布与频域谱表示}
\label{subsec:tex55-gibbs-spectrum}

\begin{definition}[统计力学目标分布]
\label{def:tex55-gibbs}
给定权重函数
\[
  w:\Omega\to(0,\infty).
\]
均匀先验对应
\[
  w(x)=\frac1N,
\]
非均匀先验对应重要性权重、经验分布或结构偏置。
定义 Gibbs/Boltzmann 分布
\[
  \pi_{\beta,w}(x)
  \DefEq
  \frac{w(x)e^{-\beta E(x)}}{Z_{\beta,w}},
  \qquad
  Z_{\beta,w}\DefEq\sum_{z\in\Omega}w(z)e^{-\beta E(z)}.
\]
其中 \(\beta>0\) 为逆温度。
当 \(\beta\to\infty\) 时，若存在正能隙，则分布质量向低能量状态集中。
\end{definition}

\begin{definition}[可逆转移核的频域谱展开]
\label{def:tex55-spectral}
令 \(P\) 是有限状态空间上的可逆马尔可夫转移核，平稳分布为 \(\pi\)。
在内积
\[
  \langle f,g\rangle_{\pi}
  \DefEq
  \sum_{x\in\Omega}f(x)g(x)\pi(x)
\]
下，若
\[
  P\phi_j=\lambda_j\phi_j,
\]
并取
\[
  1=\lambda_0,
  \qquad
  \phi_0=1,
\]
则任意密度比可展开为
\[
  \frac{\mu_t}{\pi}
  =
  1+\sum_{j=1}^{N-1}c_j(t)\phi_j.
\]
频域 \(K\)-模近似定义为
\[
  \frac{\mu_t^{(K)}}{\pi}
  =
  1+\sum_{j=1}^{K}c_j(t)\phi_j.
\]
这表示把时域遍历压缩为前 \(K\) 个主谱模态的演化。
\end{definition}

\begin{theorem}[局部时域遍历的频域截断近似]
\label{thm:tex55-spectral-truncation}
若 \(P\) 是有限状态空间上的不可约、非周期、\(\pi\)-可逆马尔可夫核，且
\[
  \mu_t=\mu_0P^t,
\]
则可用前 \(K\) 个谱模态近似 \(\mu_t\)，并且
\[
  \left\|
  \mu_t-\mu_t^{(K)}
  \right\|_{\TV}
  \le
  \frac12
  \left(
  \sum_{j>K}c_j(0)^2|\lambda_j|^{2t}
  \right)^{1/2}.
\]
\end{theorem}

\begin{Certificate}[证书：正交谱展开与全变差控制]
\label{cert:tex55-spectral-truncation}
证书包含四项：
\[
\begin{aligned}
  S_1&:\quad P\text{ 在 }L^2(\pi)\text{ 中自伴};\\
  S_2&:\quad \{\phi_0,\phi_1,\ldots,\phi_{N-1}\}\text{ 为正交特征基};\\
  S_3&:\quad c_j(t)=c_j(0)\lambda_j^t;\\
  S_4&:\quad
  \|\mu-\nu\|_{\TV}
  \le \frac12
  \left\|\frac{\mu}{\pi}-\frac{\nu}{\pi}\right\|_{L^2(\pi)}.
\end{aligned}
\]
\end{Certificate}

\begin{proof}[证明（基于数学符号推演逻辑的谓词因果链）]
令
\[
  h_t=\frac{\mu_t}{\pi}.
\]
由于 \(P\) 对 \(\pi\) 可逆，\(P\) 在 \(L^2(\pi)\) 中自伴，因此存在正交特征基
\[
  \{\phi_0,\phi_1,\ldots,\phi_{N-1}\}.
\]
展开初始密度：
\[
  h_0
  =
  1+\sum_{j=1}^{N-1}c_j(0)\phi_j.
\]
经过 \(t\) 步转移：
\[
  h_t
  =
  P^th_0
  =
  1+\sum_{j=1}^{N-1}c_j(0)\lambda_j^t\phi_j.
\]
取前 \(K\) 个频域模态：
\[
  h_t^{(K)}
  =
  1+\sum_{j=1}^{K}c_j(0)\lambda_j^t\phi_j.
\]
于是
\[
  h_t-h_t^{(K)}
  =
  \sum_{j>K}c_j(0)\lambda_j^t\phi_j.
\]
由正交性，
\[
  \|h_t-h_t^{(K)}\|_{L^2(\pi)}
  =
  \left(
  \sum_{j>K}c_j(0)^2|\lambda_j|^{2t}
  \right)^{1/2}.
\]
又由 Cauchy 不等式，
\[
  \|\mu_t-\mu_t^{(K)}\|_{\TV}
  =
  \frac12\sum_{x\in\Omega}\pi(x)|h_t(x)-h_t^{(K)}(x)|
  \le
  \frac12\|h_t-h_t^{(K)}\|_{L^2(\pi)}.
\]
合并可得
\[
  \left\|
  \mu_t-\mu_t^{(K)}
  \right\|_{\TV}
  \le
  \frac12
  \left(
  \sum_{j>K}c_j(0)^2|\lambda_j|^{2t}
  \right)^{1/2}.
\]
证毕。
\end{proof}

\begin{corollary}[长时间局部遍历等价于主频模态演化]
\label{cor:tex55-time-to-frequency}
在定理~\ref{thm:tex55-spectral-truncation} 的条件下，若
\[
  \sum_{j>K}c_j(0)^2|\lambda_j|^{2t}
\]
足够小，则
\[
  \mu_0\to\mu_1\to\mu_2\to\cdots
\]
可近似表示为
\[
  c_j(t)=c_j(0)\lambda_j^t,
  \qquad
  1\le j\le K.
\]
即
\[
  \boxed{
  \text{长时间局部遍历}
  \approx
  \text{少数主频模态的演化}
  }.
\]
\end{corollary}

\begin{Certificate}[证书：高频残差小]
\label{cert:tex55-time-to-frequency}
取
\[
  \varepsilon_{\mathrm{spec}}
  \DefEq
  \frac12
  \left(
  \sum_{j>K}c_j(0)^2|\lambda_j|^{2t}
  \right)^{1/2}.
\]
若 \(\varepsilon_{\mathrm{spec}}\le\delta\)，则
\[
  \|\mu_t-\mu_t^{(K)}\|_{\TV}\le\delta.
\]
\end{Certificate}

\begin{proof}[证明（基于数学符号推演逻辑的谓词因果链）]
由定理~\ref{thm:tex55-spectral-truncation}，
\[
  \|\mu_t-\mu_t^{(K)}\|_{\TV}
  \le
  \varepsilon_{\mathrm{spec}}.
\]
若高频残差不超过 \(\delta\)，则时域分布与前 \(K\) 个谱模态表示之间的全变差误差不超过 \(\delta\)。
因此在该误差容忍度下，可把局部遍历视为主模态系数
\[
  c_j(t)=c_j(0)\lambda_j^t
\]
的演化。证毕。
\end{proof}

\subsection{频域共振分布律、命中随机变量与统计力学命中}
\label{subsec:tex55-resonance-hit}

\begin{definition}[频域共振分布律与样本随机变量]
\label{def:tex55-resonance-distribution}
令
\[
  \phi_1,\ldots,\phi_K
\]
为选取的频域模态。
定义频域共振分布律
\[
  \nu_{\beta,\theta,K}(x)
  \DefEq
  \frac{
  w(x)e^{-\beta E(x)}
  \exp\left(\sum_{j=1}^{K}\theta_j\phi_j(x)\right)
  }{
  Z_{\beta,\theta,K}
  },
\]
其中
\[
  \theta=(\theta_1,\ldots,\theta_K)
\]
为共振调制参数，归一化常数为
\[
  Z_{\beta,\theta,K}
  \DefEq
  \sum_{z\in\Omega}
  w(z)e^{-\beta E(z)}
  \exp\left(\sum_{j=1}^{K}\theta_j\phi_j(z)\right).
\]
当 \(\theta=0\) 时，
\[
  \nu_{\beta,0,K}=\pi_{\beta,w}.
\]
当某些 \(\theta_j\ne0\) 时，表示在频域中增强或抑制对应谱模态。
严格地说，\(\nu_{\beta,\theta,K}\) 只是样本随机变量的分布律；若第 \(i\) 次命中判定使用该分布律，
真正的随机对象应写为
\[
  X_i^{(\beta,\theta,K)}\sim \nu_{\beta,\theta,K},
  \qquad
  Y_i^{(\beta,\theta,K)}
  \DefEq
  \mathbf 1_{S_*}\bigl(X_i^{(\beta,\theta,K)}\bigr).
\]
因此本节中的“频域共振采样”均应理解为由分布律 \(\nu_{\beta,\theta,K}\) 索引的随机变量簇采样机制。
\end{definition}

\begin{definition}[低能量相关模态]
\label{def:tex55-low-energy-mode}
给定低能量阈值 \(q\)，若
\[
  \left|
  \left\langle
  \mathbf 1_{\{E\le q\}},
  \phi_j
  \right\rangle_{\pi}
  \right|
\]
较大，则称 \(\phi_j\) 与低能量盆地相关。
在分布律层面，频域共振的数学含义是：
\[
  \boxed{
  \text{选择并调制与低能量盆地相关的谱模态}
  }.
\]
它不是未定义地宣称“共振会自动找到全局最优”。
\end{definition}

\begin{theorem}[独立采样的全局最优命中概率]
\label{thm:tex55-hit-probability}
令
\[
  p_{\nu}\DefEq\nu(S_*).
\]
若从分布 \(\nu\) 中独立采样 \(n\) 次，则至少一次命中全局最优集合的概率为
\[
  \Pr(\exists i\le n:X_i\in S_*)
  =
  1-(1-p_{\nu})^n.
\]
并且
\[
  \Pr(\exists i\le n:X_i\in S_*)
  \ge
  1-e^{-np_{\nu}}.
\]
\end{theorem}

\begin{Certificate}[证书：独立未命中事件乘法]
\label{cert:tex55-hit-probability}
证书为
\[
\begin{aligned}
  H_1&:\quad \Pr(X_i\notin S_*)=1-p_{\nu};\\
  H_2&:\quad X_1,\ldots,X_n\text{ 相互独立};\\
  H_3&:\quad 1-p_{\nu}\le e^{-p_{\nu}}.
\end{aligned}
\]
\end{Certificate}

\begin{proof}[证明（基于数学符号推演逻辑的谓词因果链）]
单次没有命中 \(S_*\) 的概率为
\[
  1-p_{\nu}.
\]
独立采样 \(n\) 次均未命中的概率为
\[
  (1-p_{\nu})^n.
\]
因此至少一次命中的概率为
\[
  1-(1-p_{\nu})^n.
\]
又因为
\[
  1-p_{\nu}\le e^{-p_{\nu}},
\]
所以
\[
  (1-p_{\nu})^n\le e^{-np_{\nu}}.
\]
于是
\[
  1-(1-p_{\nu})^n
  \ge
  1-e^{-np_{\nu}}.
\]
证毕。
\end{proof}

\begin{corollary}[均匀采样的命中率]
\label{cor:tex55-uniform-hit}
若直接均匀采样，则
\[
  p_{\Unif}
  =
  \frac{|S_*|}{|\Omega|}.
\]
因此
\[
  \Pr(\text{命中})
  =
  1-
  \left(
  1-\frac{|S_*|}{|\Omega|}
  \right)^n.
\]
若 \(|S_*|\ll|\Omega|\)，则有限 \(n\) 下的均匀命中概率可能很低。
\end{corollary}

\begin{Certificate}[证书：均匀质量计算]
\label{cert:tex55-uniform-hit}
证书为
\[
  \nu_{\Unif}(x)=\frac1{|\Omega|},
  \qquad
  \nu_{\Unif}(S_*)=\sum_{x\in S_*}\frac1{|\Omega|}.
\]
\end{Certificate}

\begin{proof}[证明（基于数学符号推演逻辑的谓词因果链）]
由均匀分布定义，
\[
  \nu_{\Unif}(S_*)
  =
  \sum_{x\in S_*}\frac1{|\Omega|}
  =
  \frac{|S_*|}{|\Omega|}.
\]
将该值代入定理~\ref{thm:tex55-hit-probability}，得到
\[
  \Pr(\text{命中})
  =
  1-
  \left(
  1-\frac{|S_*|}{|\Omega|}
  \right)^n.
\]
证毕。
\end{proof}

\begin{corollary}[非均匀 Gibbs 采样的最优集质量下界]
\label{cor:tex55-gibbs-hit}
令
\[
  W_*\DefEq\sum_{x\in S_*}w(x),
  \qquad
  W_{\neg *}\DefEq\sum_{x\notin S_*}w(x).
\]
设次优能隙为
\[
  \Delta
  \DefEq
  \min_{x\notin S_*}E(x)-E_*>0.
\]
则
\[
  \pi_{\beta,w}(S_*)
  \ge
  \frac{W_*}{W_*+W_{\neg *}e^{-\beta\Delta}}.
\]
因此
\[
  1-\pi_{\beta,w}(S_*)
  \le
  \frac{W_{\neg *}}{W_*}e^{-\beta\Delta}.
\]
\end{corollary}

\begin{Certificate}[证书：能隙压制非最优质量]
\label{cert:tex55-gibbs-hit}
证书包含
\[
\begin{aligned}
  G_1&:\quad E(x)=E_*,\quad x\in S_*;\\
  G_2&:\quad E(x)\ge E_*+\Delta,\quad x\notin S_*;\\
  G_3&:\quad W_*>0,\quad W_{\neg *}\ge0.
\end{aligned}
\]
\end{Certificate}

\begin{proof}[证明（基于数学符号推演逻辑的谓词因果链）]
由 Gibbs 分布定义，
\[
  \pi_{\beta,w}(S_*)
  =
  \frac{
  \sum_{x\in S_*}w(x)e^{-\beta E(x)}
  }{
  \sum_{z\in\Omega}w(z)e^{-\beta E(z)}
  }.
\]
当 \(x\in S_*\) 时，\(E(x)=E_*\)，所以分子为
\[
  W_*e^{-\beta E_*}.
\]
当 \(x\notin S_*\) 时，由能隙条件
\[
  E(x)\ge E_*+\Delta,
\]
得到
\[
  \sum_{x\notin S_*}w(x)e^{-\beta E(x)}
  \le
  W_{\neg *}e^{-\beta(E_*+\Delta)}.
\]
因此
\[
  \pi_{\beta,w}(S_*)
  \ge
  \frac{
  W_*e^{-\beta E_*}
  }{
  W_*e^{-\beta E_*}
  +
  W_{\neg *}e^{-\beta(E_*+\Delta)}
  }.
\]
约去 \(e^{-\beta E_*}\)，得到
\[
  \pi_{\beta,w}(S_*)
  \ge
  \frac{W_*}{W_*+W_{\neg *}e^{-\beta\Delta}}.
\]
进一步，
\[
  1-\pi_{\beta,w}(S_*)
  \le
  \frac{W_{\neg *}e^{-\beta\Delta}}{W_*+W_{\neg *}e^{-\beta\Delta}}
  \le
  \frac{W_{\neg *}}{W_*}e^{-\beta\Delta}.
\]
证毕。
\end{proof}

\subsection{公理降级：非线性遍历的误差递推闭合}
\label{subsec:tex55-nonlinear-induction}

\begin{definition}[非线性遍历更新算子]
\label{def:tex55-nonlinear-update}
非线性遍历可形式化为分布更新算子
\[
  \mu_{t+1}=\mathcal F(\mu_t),
\]
其中 \(\mathcal F\) 可依赖当前经验分布、能量、频域模态与采样反馈。
例如
\[
  \mathcal F(\mu)
  =
  \Normalize\left[
  \mu P+\eta R(\mu)
  \right],
\]
其中 \(P\) 是局部转移，\(R(\mu)\) 是频域共振分布律重加权项，\(\eta>0\) 是调制强度。
\end{definition}

\begin{axiom}[公理（降级为定理）：非线性遍历的收敛近似]
\label{ax:tex55-nonlinear-contraction}
设目标频域统计力学分布为
\[
  \nu=\nu_{\beta,\theta,K}.
\]
若存在距离 \(d\)、常数 \(0\le\kappa<1\) 与误差 \(\varepsilon\ge0\)，使得
\[
  d(\mathcal F(\mu),\nu)
  \le
  \kappa d(\mu,\nu)+\varepsilon,
\]
则对 \(\mu_{t+1}=\mathcal F(\mu_t)\) 有
\[
  d(\mu_t,\nu)
  \le
  \kappa^t d(\mu_0,\nu)
  +
  \frac{1-\kappa^t}{1-\kappa}\varepsilon.
\]
特别地，
\[
  \limsup_{t\to\infty}d(\mu_t,\nu)
  \le
  \frac{\varepsilon}{1-\kappa}.
\]
\end{axiom}

\begin{Certificate}[数学归纳法证书：近似收缩递推]
\label{cert:tex55-nonlinear-contraction}
定义目标谓词
\[
  P(t):\quad
  d(\mu_t,\nu)
  \le
  \kappa^t d(\mu_0,\nu)
  +
  \frac{1-\kappa^t}{1-\kappa}\varepsilon.
\]
证书为
\[
  \pi_{\mathrm{nonlin}}
  =
  \bigl(P(0),\ \forall t,\ P(t)\Rightarrow P(t+1)\bigr).
\]
\end{Certificate}

\begin{proof}[证明（基于数学符号推演逻辑的谓词因果链）]
基步 \(t=0\)：
\[
  \kappa^0d(\mu_0,\nu)
  +
  \frac{1-\kappa^0}{1-\kappa}\varepsilon
  =
  d(\mu_0,\nu).
\]
故 \(P(0)\) 成立。

归纳步：假设 \(P(t)\) 成立，即
\[
  d(\mu_t,\nu)
  \le
  \kappa^t d(\mu_0,\nu)
  +
  \frac{1-\kappa^t}{1-\kappa}\varepsilon.
\]
由收缩条件，
\[
  d(\mu_{t+1},\nu)
  =
  d(\mathcal F(\mu_t),\nu)
  \le
  \kappa d(\mu_t,\nu)+\varepsilon.
\]
代入归纳假设：
\[
\begin{aligned}
  d(\mu_{t+1},\nu)
  &\le
  \kappa
  \left[
  \kappa^t d(\mu_0,\nu)
  +
  \frac{1-\kappa^t}{1-\kappa}\varepsilon
  \right]
  +
  \varepsilon\\
  &=
  \kappa^{t+1}d(\mu_0,\nu)
  +
  \frac{\kappa-\kappa^{t+1}}{1-\kappa}\varepsilon
  +
  \varepsilon\\
  &=
  \kappa^{t+1}d(\mu_0,\nu)
  +
  \frac{1-\kappa^{t+1}}{1-\kappa}\varepsilon.
\end{aligned}
\]
故 \(P(t+1)\) 成立。
由数学归纳法，结论对所有 \(t\ge0\) 成立。
令 \(t\to\infty\)，由于 \(0\le\kappa<1\)，得到
\[
  \limsup_{t\to\infty}d(\mu_t,\nu)
  \le
  \frac{\varepsilon}{1-\kappa}.
\]
证毕。
\end{proof}

\subsection{总近似公式与主定理}
\label{subsec:tex55-total-bridge}

\begin{definition}[四类误差通道]
\label{def:tex55-error-channels}
定义总近似中的四类误差：
\[
  \varepsilon_{\mathrm{spec}}
  \DefEq
  \frac12
  \left(
  \sum_{j>K}c_j(0)^2|\lambda_j|^{2t}
  \right)^{1/2},
\]
\[
  \varepsilon_{\mathrm{nonlin}}
  \DefEq
  \kappa^t d(\mu_0,\nu)
  +
  \frac{1-\kappa^t}{1-\kappa}\varepsilon,
\]
\[
  \varepsilon_{\mathrm{temp}}
  \DefEq
  \frac{W_{\neg *}}{W_*}e^{-\beta\Delta},
\]
\[
  \varepsilon_{\mathrm{miss}}
  \DefEq
  (1-p_{\nu})^n
  \le
  e^{-np_{\nu}}.
\]
它们分别对应频域截断、非线性遍历收敛、有限温度非最优质量、有限采样未命中。
\end{definition}

\begin{theorem}[时域局部遍历到频域共振随机变量簇全局命中的近似桥接定理]
\label{thm:tex55-approximate-bridge}
设以下条件成立：
\[
\begin{aligned}
  B_1&:\quad P\text{ 具有所需可达性，并可构造可逆谱近似};\\
  B_2&:\quad \mathbf 1_{S_*}\text{ 可由前 }K\text{ 个主模态近似表达};\\
  B_3&:\quad \nu_{\beta,\theta,K}\text{ 不削弱 }S_*\text{ 的 Gibbs 质量下界};\\
  B_4&:\quad \mathcal F\text{ 满足公理~\ref{ax:tex55-nonlinear-contraction} 的近似收缩条件};\\
  B_5&:\quad \text{构造 }X_i^{(\beta,\theta,K)}\sim\nu_{\beta,\theta,K}\text{ 或其足够近似分布并判定 }n\text{ 次}.
\end{aligned}
\]
则有近似桥接
\[
  \mu_0P^t
  \approx
  \mu_t^{(K)}
  \approx
  \nu_{\beta,\theta,K}
  \Longrightarrow
  \Pr(X\in S_*)\text{ 较大}.
\]
更精确地，误差可按
\[
  \varepsilon_{\mathrm{total}}
  \le
  \varepsilon_{\mathrm{spec}}
  +
  \varepsilon_{\mathrm{nonlin}}
  +
  \varepsilon_{\mathrm{temp}}
  +
  \varepsilon_{\mathrm{miss}}
\]
控制。
\end{theorem}

\begin{Certificate}[证书：四误差桥接证书]
\label{cert:tex55-approximate-bridge}
证书由四个已签发条目拼接：
\[
\begin{aligned}
  C_1&:\quad
  \|\mu_t-\mu_t^{(K)}\|_{\TV}
  \le
  \varepsilon_{\mathrm{spec}}
  \quad\text{由定理~\ref{thm:tex55-spectral-truncation}};\\
  C_2&:\quad
  d(\mu_t,\nu)
  \le
  \varepsilon_{\mathrm{nonlin}}
  \quad\text{由公理~\ref{ax:tex55-nonlinear-contraction}};\\
  C_3&:\quad
  1-\pi_{\beta,w}(S_*)
  \le
  \varepsilon_{\mathrm{temp}}
  \quad\text{由推论~\ref{cor:tex55-gibbs-hit}};\\
  C_4&:\quad
  \Pr(\text{未命中 }S_*)
  =
  \varepsilon_{\mathrm{miss}}
  \le
  e^{-np_{\nu}}
  \quad\text{由定理~\ref{thm:tex55-hit-probability}}.
\end{aligned}
\]
\end{Certificate}

\begin{proof}[证明（基于数学符号推演逻辑的谓词因果链）]
由 \(C_1\)，时域局部遍历分布与频域前 \(K\) 模态近似之间的误差不超过
\[
  \varepsilon_{\mathrm{spec}}.
\]
由 \(C_2\)，非线性频域共振机制更新到目标分布律 \(\nu=\nu_{\beta,\theta,K}\) 的误差不超过
\[
  \varepsilon_{\mathrm{nonlin}}.
\]
由 \(C_3\)，在能隙 \(\Delta>0\) 与温度 \(\beta\) 足够大时，普通 Gibbs 分布留在非最优集合的质量满足
\[
  1-\pi_{\beta,w}(S_*)
  \le
  \varepsilon_{\mathrm{temp}}.
\]
若 \(B_3\) 成立，即共振重加权不削弱 \(S_*\) 的 Gibbs 质量下界，则该温度误差仍可作为频域共振随机变量簇命中机制的保守控制项。
由 \(C_4\)，即使采样分布已给定，有限 \(n\) 次采样仍有未命中误差
\[
  \varepsilon_{\mathrm{miss}}
  =
  (1-p_{\nu})^n
  \le
  e^{-np_{\nu}}.
\]
将四类独立证明义务作为串联近似链的误差通道相加，得到
\[
  \varepsilon_{\mathrm{total}}
  \le
  \varepsilon_{\mathrm{spec}}
  +
  \varepsilon_{\mathrm{nonlin}}
  +
  \varepsilon_{\mathrm{temp}}
  +
  \varepsilon_{\mathrm{miss}}.
\]
因此
\[
  \boxed{
  A_{\mathrm{time-local}}
  \approx
  A_{\mathrm{freq-resonant}}
  \approx
  A_{\mathrm{global-hit}}
  }
\]
只能在上述误差均可控时签发。
证毕。
\end{proof}

\begin{corollary}[高概率近似全局最优命中]
\label{cor:tex55-high-probability-hit}
若
\[
  \nu_{\beta,\theta,K}(S_*)\ge p_0>0,
\]
则构造独立随机变量 \(X_1,\ldots,X_n\sim\nu_{\beta,\theta,K}\) 并进行 \(n\) 次命中判定，有
\[
  \Pr(\text{至少一次命中 }S_*)
  \ge
  1-(1-p_0)^n
  \ge
  1-e^{-np_0}.
\]
\end{corollary}

\begin{Certificate}[证书：最优集质量下界]
\label{cert:tex55-high-probability-hit}
证书为
\[
  p_{\nu}=\nu_{\beta,\theta,K}(S_*)\ge p_0.
\]
\end{Certificate}

\begin{proof}[证明（基于数学符号推演逻辑的谓词因果链）]
由定理~\ref{thm:tex55-hit-probability}，
\[
  \Pr(\text{至少一次命中 }S_*)
  =
  1-(1-p_{\nu})^n.
\]
由证书 \(p_{\nu}\ge p_0\)，得到
\[
  1-(1-p_{\nu})^n
  \ge
  1-(1-p_0)^n.
\]
又由
\[
  (1-p_0)^n\le e^{-np_0},
\]
得到
\[
  \Pr(\text{至少一次命中 }S_*)
  \ge
  1-e^{-np_0}.
\]
证毕。
\end{proof}

\subsection{算法结构、命题与边界}
\label{subsec:tex55-algorithm-boundaries}

\begin{proposition}[局部最优困难来自谱慢模态]
\label{prop:tex55-slow-mode-trap}
若某些非平凡谱值满足
\[
  |\lambda_j|\approx1,
\]
则对应模态衰减很慢：
\[
  c_j(t)=c_j(0)\lambda_j^t.
\]
这些慢模态可表示能垒、局部盆地、慢混合方向或局部最优陷阱。
\end{proposition}

\begin{Certificate}[证书：谱系数慢衰减]
\label{cert:tex55-slow-mode-trap}
证书为
\[
  |\lambda_j|\approx1
  \Longrightarrow
  |\lambda_j|^t\text{ 在有限 }t\text{ 内不快速趋零}.
\]
\end{Certificate}

\begin{proof}[证明（基于数学符号推演逻辑的谓词因果链）]
由谱展开，
\[
  c_j(t)=c_j(0)\lambda_j^t.
\]
若 \(|\lambda_j|\approx1\)，则有限时间内
\[
  |c_j(t)|=|c_j(0)|\,|\lambda_j|^t
\]
仍可能保持显著大小。
因此该模态所表达的分布不均匀性、盆地隔离或能垒方向不会迅速消失。
这就是时域局部算法中慢混合与局部陷阱在频域中的表型。证毕。
\end{proof}

\begin{proposition}[频域共振随机变量簇增强命中质量而非替代可达性]
\label{prop:tex55-resonance-not-reachability}
频域共振机制的目标不是替代
\[
  \forall x,y,\quad \exists t,\quad P^t(x,y)>0,
\]
而是在可达或可采样的前提下提高
\[
  p_{\mathrm{res}}
  \DefEq
  \nu_{\beta,\theta,K}(S_*).
\]
理想情形为
\[
  p_{\mathrm{res}}
  \gg
  p_{\Unif}
  =
  \frac{|S_*|}{|\Omega|}.
\]
\end{proposition}

\begin{Certificate}[证书：可达性与命中质量分离]
\label{cert:tex55-resonance-not-reachability}
证书为
\[
  \text{可达性控制路径存在，}
  \qquad
  \nu(S_*)\text{ 控制采样命中概率}.
\]
\end{Certificate}

\begin{proof}[证明（基于数学符号推演逻辑的谓词因果链）]
由定义~\ref{def:tex55-reachability-hit}，可达性只说明存在 \(t\) 使 \(P^t(x,y)>0\)。
由定理~\ref{thm:tex55-hit-probability}，有限采样命中概率由
\[
  p_{\nu}=\nu(S_*)
\]
控制。
因此二者属于不同谓词：
\[
  \operatorname{Reach}(P)
  \not\equiv
  \operatorname{HitProb}(\nu,S_*).
\]
频域共振只能通过调制分布律 \(\nu_{\beta,\theta,K}\)，并由该分布律生成命中随机变量，来提高 \(S_*\) 的概率质量，
不能在缺乏可达性或采样机制时凭空签发全局命中。证毕。
\end{proof}

\begin{proposition}[均匀、Gibbs 与频域共振分布律采样机制的统一公式]
\label{prop:tex55-unified-sampling}
令
\[
  R_K(x)
  \DefEq
  \exp\left(\sum_{j=1}^{K}\theta_j\phi_j(x)\right).
\]
统一采样分布可写为
\[
  \nu(x)
  =
  \frac{
  w(x)e^{-\beta E(x)}R_K(x)
  }{
  Z
  }.
\]
其中：
\[
\begin{aligned}
  &\text{均匀采样：}
  &&w(x)=\frac1N,\quad \beta=0,\quad R_K(x)=1;\\
  &\text{普通统计力学采样：}
  &&w(x)>0,\quad \beta>0,\quad R_K(x)=1;\\
  &\text{频域共振统计力学分布律采样机制：}
  &&w(x)>0,\quad \beta>0,\quad R_K(x)\ne1.
\end{aligned}
\]
因此
\[
  \boxed{
  \text{频域共振分布律}
  =
  \text{Gibbs 采样}
  +
  \text{谱模态重加权}
  }.
\]
\end{proposition}

\begin{Certificate}[证书：参数退化关系]
\label{cert:tex55-unified-sampling}
证书为三组参数代入：
\[
  (\beta,\theta,w)
  =
  (0,0,1/N),
  \qquad
  (\beta,0,w),
  \qquad
  (\beta,\theta,w).
\]
\end{Certificate}

\begin{proof}[证明（基于数学符号推演逻辑的谓词因果链）]
将
\[
  R_K(x)=1
\]
且 \(\beta=0,w(x)=1/N\) 代入统一公式，得到均匀分布。
将 \(R_K(x)=1\) 且 \(\beta>0\) 代入，得到 Gibbs 分布
\[
  \pi_{\beta,w}.
\]
保留 \(R_K(x)\ne1\)，则得到定义~\ref{def:tex55-resonance-distribution} 中的
\[
  \nu_{\beta,\theta,K}.
\]
故三类采样是同一公式在不同参数下的特例。证毕。
\end{proof}

\begin{remark}[算法结构]
对应算法可写为五步：
\begin{enumerate}
  \item 建立局部转移图
  \[
    G=(\Omega,\mathcal E),
    \qquad
    (x,y)\in\mathcal E\Longleftrightarrow y\in N(x).
  \]
  \item 构造局部转移核，并在需要时加入少量全局跳跃：
  \[
    P_{\alpha}
    =
    (1-\alpha)P_{\mathrm{local}}
    +
    \alpha Q_{\mathrm{global}},
    \qquad
    \alpha>0.
  \]
  若
  \[
    Q_{\mathrm{global}}(x,y)>0,\quad \forall x,y,
  \]
  则可修复纯局部核的不可达问题。
  \item 求解或近似求解
  \[
    P_{\alpha}\phi_j=\lambda_j\phi_j
  \]
  得到前 \(K\) 个主模态 \(\phi_1,\ldots,\phi_K\)。
  \item 按低能量相关性选择共振参数，例如
  \[
    \theta_j
    =
    \eta
    \left\langle
    \mathbf 1_{\{E\le q\}},
    \phi_j
    \right\rangle_{\pi}.
  \]
  \item 从
  \[
    \nu_{\beta,\theta,K}
  \]
  中采样 \(n\) 次，或用 Metropolis--Hastings 校正采样，并用
  \[
    1-(1-\nu_{\beta,\theta,K}(S_*))^n
  \]
  评估命中概率。
\end{enumerate}
\end{remark}

\begin{remark}[必要边界]
本章结论有四个必要边界。

第一，没有全局可达性，就不能保证全局命中。
若存在 \(x,y\) 使
\[
  P^t(x,y)=0,\qquad \forall t,
\]
且 \(y\in S_*\)，则从 \(x\) 出发全局最优不可达。

第二，没有正能隙，低温集中结论会变弱。
若
\[
  \Delta=\min_{x\notin S_*}E(x)-E_*
\]
很小，则 \(e^{-\beta\Delta}\) 下降很慢，非最优状态仍可能占有明显概率质量。

第三，前 \(K\) 个频域模态必须表达全局盆地。
若
\[
  \left\langle \mathbf 1_{S_*},\phi_j\right\rangle_{\pi}\approx0,
  \qquad
  j=1,\ldots,K,
\]
则前 \(K\) 个模态无法有效增强全局最优区域。
形式上需要
\[
  \Pi_K\mathbf 1_{S_*}\approx \mathbf 1_{S_*}.
\]

第四，有限计算预算只能得到概率保证。
有限 \(n\) 次采样只给出
\[
  \Pr(\text{命中})=1-(1-p)^n,
\]
不能推出 \(\Pr(\text{命中})=1\)，除非 \(p=1\) 或 \(n\to\infty\)。
所以有限预算下的严格表述应是
\[
  \boxed{
  \text{以高概率命中全局最优的近似算法}
  }.
\]
\end{remark}

\begin{remark}[本章最终口径]
本章最终签发的结论是
\[
  \boxed{
  \text{时域局部最优遍历}
  \xrightarrow{\text{谱分解}}
  \text{频域模态表示}
  \xrightarrow{\text{共振重加权}}
  \text{统计力学采样}
  \xrightarrow{\text{概率命中}}
  \text{近似全局最优}.
  }
\]
其核心概率公式为
\[
  \boxed{
  \Pr(\text{命中 }S_*)
  \ge
  1-
  \exp\left(
  -n\nu_{\beta,\theta,K}(S_*)
  \right)
  }.
\]
其核心误差公式为
\[
  \boxed{
  \varepsilon_{\mathrm{total}}
  \le
  \varepsilon_{\mathrm{spec}}
  +
  \varepsilon_{\mathrm{nonlin}}
  +
  \varepsilon_{\mathrm{temp}}
  +
  \varepsilon_{\mathrm{miss}}.
  }
\]
因此，该思想的严格版本是：
\[
  \boxed{
  \text{用频域谱模态压缩时域局部遍历，用统计力学分布提高全局最优集合的命中概率。}
  }
\]
它不是无条件“局部最优变全局最优”的定理；它是在全局可达、谱可表达、低温集中、收敛近似与采样充分条件下的可控近似命中定理。
\end{remark}

\clearpage

%%%%%%%%%%%%%%%%%%%%%%%%%%%%%%%%%%%%%%%%%%%%%%%%%%%%%%%%%%%%
\section{形式逻辑：时域局部算法优选到频域统计力学分布律优化}
\label{sec:tex55-local-algorithm-to-law-selection}
%%%%%%%%%%%%%%%%%%%%%%%%%%%%%%%%%%%%%%%%%%%%%%%%%%%%%%%%%%%%

\begin{remark}[核心陈述]
本章把“优选时域局部最优算法”严格化为“优选频域统计力学分布律的类型与参数”。
形式上，要刻画的桥接关系是
\[
  \boxed{
  \text{从优选时域局部算法}
  \Longrightarrow
  \text{转化为优选频域统计力学分布律的类型与参数}
  }.
\]
更精确地说，有限时间局部算法优选问题
\[
  \max_{a\in\mathcal A_{\mathrm{local}}}J_T(a)
\]
可近似改写为
\[
  \max_{\tau,\vartheta}
  \left[
    1-\left(1-\pi_{\tau,\vartheta}(S_\varepsilon)\right)^m
    -m\delta_{\mathrm{mix}}(\tau,\vartheta)
    -\delta_{\mathrm{est}}(\tau,\vartheta)
  \right].
\]
当混合误差与估计误差足够小时，该问题进一步退化为
\[
  \boxed{
  \max_{\tau,\vartheta}\pi_{\tau,\vartheta}(S_\varepsilon)
  }.
\]
核心判断是：
\[
  \boxed{
  \begin{gathered}
  \text{局部算法优选问题可以转化为目标分布律优化问题，}\\
  \text{但必须同时控制采样该分布的谱混合性质。}
  \end{gathered}
  }
\]
只优化分布律而不控制混合速度，并不足以推出有限时间命中。
\end{remark}

\begin{remark}[阅读导航与引用口径]
本章承接上一章的马尔可夫链谱表示、Gibbs 采样与频域共振随机变量簇口径，
并把焦点从“给定分布后的命中概率”推进到“如何选择分布律类型与参数”。
马尔可夫链混合与谱隙估计参考 \cite{LevinPeresWilmer2017Markov,Norris1998MarkovChains}；
自由能与 Gibbs 分布的变分形式参考统计力学标准口径；
精英采样、重要性采样和交叉熵式参数更新可参考 \cite{Rubinstein1999CrossEntropy}。
\end{remark}

\subsection{对象域：近似最优集、局部算法与分布律}
\label{subsec:tex55-law-selection-objects}

\begin{definition}[优化状态空间与近似全局最优集合]
\label{def:tex55-law-state-space}
令
\[
  \Omega=\{x_1,x_2,\ldots,x_N\}
\]
为有限状态空间，令
\[
  E:\Omega\to\RR
\]
为待最小化的目标函数，也称能量函数。
定义全局最优值、全局最优集合与近似全局最优集合为
\[
  E_*\DefEq\min_{x\in\Omega}E(x),
  \qquad
  S_*\DefEq\{x\in\Omega:E(x)=E_*\},
\]
\[
  S_\varepsilon
  \DefEq
  \{x\in\Omega:E(x)\le E_*+\varepsilon\},
  \qquad
  \varepsilon\ge0.
\]
当 \(\varepsilon=0\) 时，
\[
  S_\varepsilon=S_*.
\]
\end{definition}

\begin{definition}[时域局部算法与有限时间命中目标]
\label{def:tex55-local-algorithm-objective}
一个时域局部算法 \(A_a\) 表示为马尔可夫过程
\[
  X_{t+1}\sim P_a(X_t,\cdot),
  \qquad
  a\in\mathcal A_{\mathrm{local}},
\]
其中 \(a\) 是算法类型或算法参数。
若
\[
  P_a(x,y)>0
\]
主要发生在局部邻域
\[
  y\in N_a(x),
\]
则称 \(P_a\) 是局部转移核。

局部算法的目标不是单步最优，而是有限时间内的整体命中质量：
\[
  J_T(a)
  \DefEq
  \Pr_a(\exists t\le T:X_t\in S_\varepsilon).
\]
因此原始优选问题为
\[
  a^\star
  \in
  \arg\max_{a\in\mathcal A_{\mathrm{local}}}J_T(a).
\]
\end{definition}

\begin{definition}[频域统计力学分布律]
\label{def:tex55-statistical-law}
定义一族频域统计力学分布律
\[
  \pi_{\tau,\vartheta}(x)
  \DefEq
  \frac{
  W_\vartheta(x)
  \Psi_{\tau}\!\left(
  -\beta E(x)+\sum_{j=1}^{K}\theta_j\phi_j(x)
  \right)
  }{
  Z_{\tau,\vartheta}
  }.
\]
其中：
\[
\begin{aligned}
  \tau&:\quad \text{分布类型};\\
  \vartheta&=(\beta,W_\vartheta,\theta_1,\ldots,\theta_K,K):\quad \text{分布参数};\\
  W_\vartheta(x)&>0:\quad \text{均匀或非均匀先验权重};\\
  \beta&>0:\quad \text{逆温度};\\
  \phi_j&:\quad \text{频域模态，例如转移核或图拉普拉斯算子的特征函数};\\
  \theta_j&:\quad \text{频域共振分布律参数};\\
  Z_{\tau,\vartheta}&:\quad \text{归一化常数}.
\end{aligned}
\]
若
\[
  W_\vartheta(x)=\frac1{|\Omega|},
\]
则为均匀先验；若
\[
  W_\vartheta(x)\ne\frac1{|\Omega|},
\]
则为非均匀先验。
\end{definition}

\begin{definition}[频域分布律优化目标]
\label{def:tex55-law-objective}
频域分布律优化问题定义为
\[
  (\tau^\star,\vartheta^\star)
  \in
  \arg\max_{\tau,\vartheta}
  \mathcal H_T(\tau,\vartheta),
\]
其中 \(\mathcal H_T\) 同时刻画分布命中概率、混合误差与估计误差。
一个典型形式为
\[
  \mathcal H_T(\tau,\vartheta)
  \DefEq
  1-\exp\left(
    -m\left[
      \pi_{\tau,\vartheta}(S_\varepsilon)
      -\delta_{\mathrm{mix}}(\tau,\vartheta)
    \right]_+
  \right)
  -
  \delta_{\mathrm{est}}(\tau,\vartheta).
\]
其中 \(m\) 是有效采样次数；
\(\delta_{\mathrm{mix}}\) 表示采样分布尚未接近目标分布导致的混合误差；
\(\delta_{\mathrm{est}}\) 表示有限样本估计频域模态、能量分布或参数梯度造成的误差。
\end{definition}

\subsection{公理降级：独立采样命中概率的数学归纳证书}
\label{subsec:tex55-hit-induction-law}

\begin{axiom}[公理（降级为定理）：独立采样命中概率归纳闭合]
\label{ax:tex55-iid-hit-induction}
设
\[
  X_1,\ldots,X_m\sim\pi
\]
独立同分布，并令
\[
  p\DefEq\pi(S_\varepsilon).
\]
则至少一次命中 \(S_\varepsilon\) 的概率为
\[
  \Pr(\exists i\le m:X_i\in S_\varepsilon)
  =
  1-(1-p)^m.
\]
并且
\[
  1-(1-p)^m\ge 1-e^{-mp}.
\]
\end{axiom}

\begin{Certificate}[数学归纳法证书：未命中事件乘法链]
\label{cert:tex55-iid-hit-induction}
定义谓词
\[
  P(m):\quad
  \Pr(\exists i\le m:X_i\in S_\varepsilon)
  =
  1-(1-p)^m.
\]
证书为
\[
  \pi_{\mathrm{hit}}
  =
  \bigl(P(1),\ \forall m\ge1,\ P(m)\Rightarrow P(m+1)\bigr).
\]
\end{Certificate}

\begin{proof}[证明（基于数学符号推演逻辑的谓词因果链）]
基步 \(m=1\)：
\[
  \Pr(X_1\in S_\varepsilon)=p
  =
  1-(1-p)^1.
\]
故 \(P(1)\) 成立。

归纳步：假设 \(P(m)\) 成立。
对 \(m+1\) 次采样，有
\[
\begin{aligned}
  &\Pr(\exists i\le m+1:X_i\in S_\varepsilon)\\
  &=
  1-
  \Pr(X_1\notin S_\varepsilon,\ldots,X_{m+1}\notin S_\varepsilon).
\end{aligned}
\]
由独立性，
\[
  \Pr(X_1\notin S_\varepsilon,\ldots,X_{m+1}\notin S_\varepsilon)
  =
  (1-p)^{m+1}.
\]
因此
\[
  \Pr(\exists i\le m+1:X_i\in S_\varepsilon)
  =
  1-(1-p)^{m+1},
\]
故 \(P(m+1)\) 成立。
由数学归纳法，
\[
  \forall m\ge1,\quad P(m).
\]
又因为
\[
  1-p\le e^{-p},
\]
所以
\[
  (1-p)^m\le e^{-mp},
\]
从而
\[
  1-(1-p)^m\ge1-e^{-mp}.
\]
证毕。
\end{proof}

\subsection{引理与命题：从算法优选到分布律优选}
\label{subsec:tex55-selection-proposition}

\begin{lemma}[分布命中质量与有限重启命中目标的单调关系]
\label{lem:tex55-hit-quality-monotone}
固定有效采样次数 \(m\ge1\)。
理想独立采样命中函数
\[
  G_m(p)\DefEq1-(1-p)^m
\]
在 \(p\in[0,1]\) 上单调递增。
\end{lemma}

\begin{Certificate}[证书：一阶导数非负]
\label{cert:tex55-hit-quality-monotone}
证书为
\[
  G_m'(p)=m(1-p)^{m-1}\ge0.
\]
\end{Certificate}

\begin{proof}[证明（基于数学符号推演逻辑的谓词因果链）]
对
\[
  G_m(p)=1-(1-p)^m
\]
求导，得
\[
  G_m'(p)=m(1-p)^{m-1}.
\]
当 \(p\in[0,1]\) 且 \(m\ge1\) 时，
\[
  m(1-p)^{m-1}\ge0.
\]
因此 \(G_m\) 关于 \(p\) 单调递增。
证毕。
\end{proof}

\begin{proposition}[时域局部算法优选可转化为频域分布律优化]
\label{prop:tex55-local-to-law}
若每个局部算法 \(A_a\) 都对应一个平稳分布 \(\pi_a\) 与一个采样核 \(P_a\)，
并且存在参数化表示
\[
  \pi_a=\pi_{\tau(a),\vartheta(a)},
\]
则时域局部算法优选问题
\[
  \max_{a\in\mathcal A_{\mathrm{local}}}J_T(a)
\]
可近似转化为频域分布律优化问题
\[
  \max_{\tau,\vartheta}
  \left[
    \pi_{\tau,\vartheta}(S_\varepsilon)
    -
    \delta_{\mathrm{mix}}(\tau,\vartheta)
    -
    \delta_{\mathrm{est}}(\tau,\vartheta)
  \right],
\]
或更精确地转化为定义~\ref{def:tex55-law-objective} 中的 \(\mathcal H_T\) 优化。
\end{proposition}

\begin{Certificate}[证书：平稳分布参数化与误差扣除]
\label{cert:tex55-local-to-law}
证书由四项组成：
\[
\begin{aligned}
  L_1&:\quad A_a\mapsto(P_a,\pi_a);\\
  L_2&:\quad \pi_a=\pi_{\tau(a),\vartheta(a)};\\
  L_3&:\quad \text{混合充分时单次命中概率近似为 }\pi_a(S_\varepsilon);\\
  L_4&:\quad \text{有限时间与有限估计引入 }\delta_{\mathrm{mix}},\delta_{\mathrm{est}}.
\end{aligned}
\]
\end{Certificate}

\begin{proof}[证明（基于数学符号推演逻辑的谓词因果链）]
时域局部算法的有限时间命中概率为
\[
  J_T(a)=\Pr_a(\exists t\le T:X_t\in S_\varepsilon).
\]
若算法 \(A_a\) 的长期分布为 \(\pi_a\)，则在混合充分后，单次采样命中 \(S_\varepsilon\) 的概率近似为
\[
  p_a=\pi_a(S_\varepsilon).
\]
若
\[
  \pi_a=\pi_{\tau(a),\vartheta(a)},
\]
则
\[
  p_a=\pi_{\tau(a),\vartheta(a)}(S_\varepsilon).
\]
由引理~\ref{lem:tex55-hit-quality-monotone}，在固定有效采样次数下，理想命中概率关于 \(p_a\) 单调递增。
因此优化算法 \(a\) 的核心目标可转化为选择使
\[
  \pi_{\tau,\vartheta}(S_\varepsilon)
\]
尽量大的分布律类型与参数。

但是有限时间内算法未必已达到目标分布，因此必须扣除混合误差
\[
  \delta_{\mathrm{mix}}(\tau,\vartheta).
\]
此外，频域模态 \(\phi_j\)、能量分布与参数梯度常由样本估计，因此还需扣除估计误差
\[
  \delta_{\mathrm{est}}(\tau,\vartheta).
\]
综上，
\[
  \max_a J_T(a)
\]
可近似转化为
\[
  \max_{\tau,\vartheta}
  \left[
    \pi_{\tau,\vartheta}(S_\varepsilon)
    -
    \delta_{\mathrm{mix}}(\tau,\vartheta)
    -
    \delta_{\mathrm{est}}(\tau,\vartheta)
  \right],
\]
或采用 \(\mathcal H_T\) 的重启命中形式。
命题得证。
\end{proof}

\subsection{定理：有限时间混合误差与谱控制}
\label{subsec:tex55-mixing-theorems-law}

\begin{theorem}[分布律命中概率与时域算法命中概率的误差界]
\label{thm:tex55-finite-time-hit-error}
设局部算法每次从同一初始分布出发，经过 \(b\) 步后的分布为 \(\mu_b\)。
若
\[
  \|\mu_b-\pi_{\tau,\vartheta}\|_{\TV}\le\delta_b,
\]
并且进行 \(m\) 次独立重启采样，每次运行 \(b\) 步，则实际命中概率 \(J_{m,b}\) 满足
\[
  \left|
  J_{m,b}
  -
  \left[
  1-
  \left(
  1-\pi_{\tau,\vartheta}(S_\varepsilon)
  \right)^m
  \right]
  \right|
  \le
  m\delta_b.
\]
\end{theorem}

\begin{Certificate}[证书：单次全变差误差到乘积事件误差]
\label{cert:tex55-finite-time-hit-error}
证书包含
\[
\begin{aligned}
  M_1&:\quad
  \|\mu_b-\pi_{\tau,\vartheta}\|_{\TV}\le\delta_b;\\
  M_2&:\quad
  |\mu_b(B)-\pi_{\tau,\vartheta}(B)|\le\delta_b,\quad \forall B\subseteq\Omega;\\
  M_3&:\quad
  \|\mu_b^{\otimes m}-\pi_{\tau,\vartheta}^{\otimes m}\|_{\TV}\le m\delta_b.
\end{aligned}
\]
\end{Certificate}

\begin{proof}[证明（基于数学符号推演逻辑的谓词因果链）]
令
\[
  p=\pi_{\tau,\vartheta}(S_\varepsilon).
\]
如果每次采样都严格来自 \(\pi_{\tau,\vartheta}\)，则由公理~\ref{ax:tex55-iid-hit-induction}，
\[
  J_{\mathrm{ideal}}
  =
  1-(1-p)^m.
\]
实际每次采样来自 \(\mu_b\)。
由全变差条件，对任意事件 \(B\subseteq\Omega\)，有
\[
  |\mu_b(B)-\pi_{\tau,\vartheta}(B)|\le\delta_b.
\]
对 \(m\) 次独立重启的乘积分布，利用乘积测度的标准耦合或逐坐标替换界，有
\[
  \|\mu_b^{\otimes m}-\pi_{\tau,\vartheta}^{\otimes m}\|_{\TV}\le m\delta_b.
\]
令
\[
  \mathcal E_m
  \DefEq
  \{(x_1,\ldots,x_m):\exists i\le m,\ x_i\in S_\varepsilon\}
\]
为至少一次命中事件。
于是
\[
  |J_{m,b}-J_{\mathrm{ideal}}|
  =
  |\mu_b^{\otimes m}(\mathcal E_m)-\pi_{\tau,\vartheta}^{\otimes m}(\mathcal E_m)|
  \le
  m\delta_b.
\]
即
\[
  \left|
  J_{m,b}
  -
  \left[
  1-
  \left(
  1-\pi_{\tau,\vartheta}(S_\varepsilon)
  \right)^m
  \right]
  \right|
  \le
  m\delta_b.
\]
证毕。
\end{proof}

\begin{theorem}[混合误差由频域谱决定]
\label{thm:tex55-spectral-mixing-law}
若采样核 \(P_{\tau,\vartheta}\) 对目标分布 \(\pi_{\tau,\vartheta}\) 可逆，且特征值满足
\[
  1=\lambda_0>
  |\lambda_1|\ge|\lambda_2|\ge\cdots,
\]
则从初始分布 \(\mu_0\) 出发，经过 \(b\) 步后的混合误差满足
\[
  \|\mu_0P_{\tau,\vartheta}^b-\pi_{\tau,\vartheta}\|_{\TV}
  \le
  \frac12
  \sqrt{\chi^2(\mu_0\mid\pi_{\tau,\vartheta})}\,
  |\lambda_1|^b,
\]
其中
\[
  \chi^2(\mu_0\mid\pi)
  \DefEq
  \sum_{x\in\Omega}
  \frac{(\mu_0(x)-\pi(x))^2}{\pi(x)}.
\]
\end{theorem}

\begin{Certificate}[证书：谱半径控制 \(L^2\) 残差]
\label{cert:tex55-spectral-mixing-law}
证书由三项构成：
\[
\begin{aligned}
  R_1&:\quad h_0=\mu_0/\pi_{\tau,\vartheta}
  =1+\sum_{j\ge1}c_j\phi_j;\\
  R_2&:\quad h_b=1+\sum_{j\ge1}c_j\lambda_j^b\phi_j;\\
  R_3&:\quad
  \|\mu_b-\pi\|_{\TV}\le
  \frac12\|h_b-1\|_{L^2(\pi)}.
\end{aligned}
\]
\end{Certificate}

\begin{proof}[证明（基于数学符号推演逻辑的谓词因果链）]
令
\[
  \pi=\pi_{\tau,\vartheta},
  \qquad
  h_0=\frac{\mu_0}{\pi}.
\]
在 \(L^2(\pi)\) 中展开
\[
  h_0=1+\sum_{j\ge1}c_j\phi_j.
\]
由于 \(P_{\tau,\vartheta}\) 对 \(\pi\) 可逆，其特征函数构成正交谱基，经过 \(b\) 步后
\[
  h_b=1+\sum_{j\ge1}c_j\lambda_j^b\phi_j.
\]
于是
\[
  \|h_b-1\|_{L^2(\pi)}
  \le
  |\lambda_1|^b
  \left(
  \sum_{j\ge1}c_j^2
  \right)^{1/2}.
\]
由 \(\chi^2\) 展开，
\[
  \sum_{j\ge1}c_j^2
  =
  \chi^2(\mu_0\mid\pi).
\]
又由 Cauchy 不等式，
\[
  \|\mu_0P_{\tau,\vartheta}^b-\pi\|_{\TV}
  \le
  \frac12\|h_b-1\|_{L^2(\pi)}.
\]
合并得到
\[
  \|\mu_0P_{\tau,\vartheta}^b-\pi_{\tau,\vartheta}\|_{\TV}
  \le
  \frac12
  \sqrt{\chi^2(\mu_0\mid\pi_{\tau,\vartheta})}\,
  |\lambda_1|^b.
\]
证毕。
\end{proof}

\subsection{总转化定理}
\label{subsec:tex55-total-selection-theorem}

\begin{theorem}[时域局部算法优选到频域统计力学分布律优化的近似等价]
\label{thm:tex55-local-selection-equivalence}
设 \(A_a\) 是时域局部算法，\(P_a\) 是其转移核，\(\pi_a\) 是其平稳分布。
若存在参数化映射
\[
  a\mapsto(\tau,\vartheta)
\]
使得
\[
  \pi_a=\pi_{\tau,\vartheta},
\]
并且 \(P_a\) 的混合误差为
\[
  \delta_b(a)
  \DefEq
  \|\mu_0P_a^b-\pi_a\|_{\TV},
\]
则有限时间重启命中优化问题
\[
  \max_a J_{m,b}(a)
\]
可近似为
\[
  \max_{\tau,\vartheta}
  \left\{
  1-
  \left[
  1-\pi_{\tau,\vartheta}(S_\varepsilon)
  \right]^m
  -
  m\delta_b(\tau,\vartheta)
  \right\}.
\]
当
\[
  m\delta_b(\tau,\vartheta)\ll1
\]
且估计误差可忽略时，该问题进一步近似为
\[
  \boxed{
  \max_{\tau,\vartheta}\pi_{\tau,\vartheta}(S_\varepsilon)
  }.
\]
\end{theorem}

\begin{Certificate}[证书：算法参数到分布参数的桥接]
\label{cert:tex55-local-selection-equivalence}
证书为
\[
\begin{aligned}
  T_1&:\quad \pi_a=\pi_{\tau,\vartheta};\\
  T_2&:\quad
  J_{m,b}(a)
  =
  1-\left[1-\pi_a(S_\varepsilon)\right]^m
  +O(m\delta_b(a));\\
  T_3&:\quad
  G_m(p)=1-(1-p)^m\text{ 对 }p\text{ 单调递增};\\
  T_4&:\quad
  m\delta_b(\tau,\vartheta)\text{ 与 }\delta_{\mathrm{est}}\text{ 可控}.
\end{aligned}
\]
\end{Certificate}

\begin{proof}[证明（基于数学符号推演逻辑的谓词因果链）]
由定理~\ref{thm:tex55-finite-time-hit-error}，
\[
  J_{m,b}(a)
  =
  1-
  \left[
  1-\pi_a(S_\varepsilon)
  \right]^m
  +
  O(m\delta_b(a)).
\]
由参数化条件
\[
  \pi_a=\pi_{\tau,\vartheta},
\]
得到
\[
  J_{m,b}(a)
  =
  1-
  \left[
  1-\pi_{\tau,\vartheta}(S_\varepsilon)
  \right]^m
  +
  O(m\delta_b(\tau,\vartheta)).
\]
因此优化 \(a\) 等价于优化
\[
  1-
  \left[
  1-\pi_{\tau,\vartheta}(S_\varepsilon)
  \right]^m
  -
  m\delta_b(\tau,\vartheta)
\]
的保守下界。
当
\[
  m\delta_b(\tau,\vartheta)\ll1
\]
且估计误差可控时，主要目标由
\[
  1-
  \left[
  1-\pi_{\tau,\vartheta}(S_\varepsilon)
  \right]^m
\]
决定。
由引理~\ref{lem:tex55-hit-quality-monotone}，该项关于
\[
  \pi_{\tau,\vartheta}(S_\varepsilon)
\]
单调递增。
故问题进一步近似为
\[
  \max_{\tau,\vartheta}\pi_{\tau,\vartheta}(S_\varepsilon).
\]
证毕。
\end{proof}

\subsection{分布律类型、参数梯度与自由能}
\label{subsec:tex55-law-types-gradient-free-energy}

\begin{definition}[统计力学分布律类型库]
\label{def:tex55-law-type-library}
本章考虑以下五类常用分布律。

第一，Gibbs/Boltzmann 分布：
\[
  \pi_{\beta,w}(x)
  =
  \frac{w(x)e^{-\beta E(x)}}{Z_{\beta,w}},
  \qquad
  \vartheta=(\beta,w).
\]
其中 \(\beta\) 控制开发强度，\(w\) 控制先验偏置。

第二，退火分布：
\[
  \pi_t(x)
  =
  \frac{w(x)e^{-\beta_tE(x)}}{Z_t},
  \qquad
  \vartheta=(\beta_t)_{t\ge0}.
\]
早期 \(\beta_t\) 较小有利于遍历，后期 \(\beta_t\) 较大有利于低能量集中。

第三，频域共振 Gibbs 分布律：
\[
  \pi_{\beta,\theta,K}(x)
  =
  \frac{
  w(x)e^{-\beta E(x)+\sum_{j=1}^{K}\theta_j\phi_j(x)}
  }{
  Z_{\beta,\theta,K}
  },
  \qquad
  \vartheta=(\beta,w,\theta_1,\ldots,\theta_K,K).
\]
其中 \(\phi_j\) 表达全局几何、慢混合方向、局部盆地连接方向，\(\theta_j\) 调节这些模态的影响。

第四，混合分布：
\[
  \pi(x)=\sum_{\ell=1}^{L}\alpha_\ell\pi_{\beta_\ell,w_\ell}(x),
  \qquad
  \alpha_\ell\ge0,\quad \sum_{\ell=1}^{L}\alpha_\ell=1.
\]
它保留多个温度、多个先验或多个盆地，以降低单一分布陷入局部盆地的风险。

第五，Tsallis 型重尾分布：
\[
  \pi_{q,\beta}(x)
  \propto
  w(x)
  \left[
  1-(1-q)\beta E(x)
  \right]_+^{\frac1{1-q}}.
\]
重尾分布可保留远距离跳跃与跨能垒探索能力。
\end{definition}

\begin{proposition}[频域共振 Gibbs 分布律的参数梯度公式]
\label{prop:tex55-gradient-law}
考虑
\[
  \pi_{\beta,\theta,K}(x)
  =
  \frac{
  w(x)e^{-\beta E(x)+\sum_{j=1}^{K}\theta_j\phi_j(x)}
  }{
  Z_{\beta,\theta,K}
  }.
\]
令
\[
  p(\beta,\theta)
  \DefEq
  \pi_{\beta,\theta,K}(S_\varepsilon).
\]
则
\[
  \frac{\partial}{\partial\theta_j}\log p(\beta,\theta)
  =
  \mathbb E_{\pi}[\phi_j\mid S_\varepsilon]
  -
  \mathbb E_{\pi}[\phi_j],
\]
并且
\[
  \frac{\partial}{\partial\beta}\log p(\beta,\theta)
  =
  \mathbb E_{\pi}[E]
  -
  \mathbb E_{\pi}[E\mid S_\varepsilon].
\]
\end{proposition}

\begin{Certificate}[证书：指数族得分函数]
\label{cert:tex55-gradient-law}
证书为
\[
\begin{aligned}
  D_{\theta_j}\log\pi(x)
  &=
  \phi_j(x)-\mathbb E_\pi[\phi_j];\\
  D_{\beta}\log\pi(x)
  &=
  -E(x)+\mathbb E_\pi[E];\\
  p(\beta,\theta)&=\sum_{x\in S_\varepsilon}\pi(x).
\end{aligned}
\]
\end{Certificate}

\begin{proof}[证明（基于数学符号推演逻辑的谓词因果链）]
由定义，
\[
  p(\beta,\theta)=\sum_{x\in S_\varepsilon}\pi_{\beta,\theta,K}(x).
\]
对 \(\theta_j\) 求导：
\[
  \frac{\partial}{\partial\theta_j}\log\pi(x)
  =
  \phi_j(x)-\mathbb E_\pi[\phi_j].
\]
因此
\[
  \frac{\partial p}{\partial\theta_j}
  =
  \sum_{x\in S_\varepsilon}
  \pi(x)
  \left[
  \phi_j(x)-\mathbb E_\pi[\phi_j]
  \right].
\]
两边除以 \(p\)，得到
\[
  \frac{\partial}{\partial\theta_j}\log p
  =
  \mathbb E_{\pi}[\phi_j\mid S_\varepsilon]
  -
  \mathbb E_{\pi}[\phi_j].
\]
同理，
\[
  \frac{\partial}{\partial\beta}\log\pi(x)
  =
  -E(x)+\mathbb E_\pi[E].
\]
故
\[
  \frac{\partial}{\partial\beta}\log p
  =
  -\mathbb E_{\pi}[E\mid S_\varepsilon]
  +
  \mathbb E_{\pi}[E].
\]
即
\[
  \frac{\partial}{\partial\beta}\log p
  =
  \mathbb E_{\pi}[E]
  -
  \mathbb E_{\pi}[E\mid S_\varepsilon].
\]
证毕。
\end{proof}

\begin{theorem}[Gibbs 分布是自由能最优分布]
\label{thm:tex55-free-energy-gibbs}
设 \(w\) 是先验分布，定义自由能泛函
\[
  \mathcal F_\beta(q)
  \DefEq
  \mathbb E_q[E]
  +
  \frac1\beta
  \operatorname{KL}(q\mid w),
  \qquad
  q\in\Delta(\Omega).
\]
则
\[
  \pi_{\beta,w}(x)
  =
  \frac{w(x)e^{-\beta E(x)}}{Z_{\beta,w}}
\]
是 \(\mathcal F_\beta(q)\) 的唯一最小点。
\end{theorem}

\begin{Certificate}[证书：KL 非负性]
\label{cert:tex55-free-energy-gibbs}
证书为
\[
  \operatorname{KL}(q\mid\pi_{\beta,w})
  =
  \beta\mathcal F_\beta(q)+\log Z_{\beta,w},
\]
以及
\[
  \operatorname{KL}(q\mid\pi_{\beta,w})\ge0,
  \quad
  \text{等号当且仅当 }q=\pi_{\beta,w}.
\]
\end{Certificate}

\begin{proof}[证明（基于数学符号推演逻辑的谓词因果链）]
由 KL 散度定义，
\[
  \operatorname{KL}(q\mid\pi_{\beta,w})
  =
  \sum_x q(x)\log\frac{q(x)}{\pi_{\beta,w}(x)}.
\]
代入
\[
  \pi_{\beta,w}(x)=\frac{w(x)e^{-\beta E(x)}}{Z_{\beta,w}},
\]
得到
\[
\begin{aligned}
  \operatorname{KL}(q\mid\pi_{\beta,w})
  &=
  \sum_x q(x)\log\frac{q(x)}{w(x)}
  +
  \beta\sum_x q(x)E(x)
  +
  \log Z_{\beta,w}\\
  &=
  \beta\mathcal F_\beta(q)+\log Z_{\beta,w}.
\end{aligned}
\]
由于
\[
  \operatorname{KL}(q\mid\pi_{\beta,w})\ge0
\]
且等号成立当且仅当
\[
  q=\pi_{\beta,w},
\]
所以 \(\mathcal F_\beta(q)\) 的唯一最小点为
\[
  \pi_{\beta,w}.
\]
证毕。
\end{proof}

\subsection{算法结构、意义与边界}
\label{subsec:tex55-law-algorithm-boundaries}

\begin{remark}[算法结构]
从局部算法选择到分布律选择可整理为六步。
\begin{enumerate}
  \item 对每个局部算法 \(A_a\) 写出转移核
  \[
    P_a(x,y),
    \qquad
    P_a(x,y)>0\Longleftrightarrow y\in N_a(x)
  \]
  或其带重启、退火、跃迁修复后的版本。
  \item 对 \(P_a\) 或状态图拉普拉斯算子 \(L_a\) 做谱分解：
  \[
    P_a\phi_j=\lambda_j\phi_j,
    \qquad
    L_a\phi_j=\omega_j\phi_j.
  \]
  保留前 \(K\) 个模态 \(\phi_1,\ldots,\phi_K\)，作为全局盆地、瓶颈、能垒与慢混合方向的频域特征。
  \item 选择分布类型
  \[
    \tau\in
    \{\text{Gibbs},\text{退火},\text{频域共振 Gibbs 分布律},\text{混合分布},\text{重尾分布}\}.
  \]
  \item 求解参数优化：
  \[
    (\tau^\star,\vartheta^\star)
    \in
    \arg\max_{\tau,\vartheta}
    \left[
    1-\left(1-\pi_{\tau,\vartheta}(S_\varepsilon)\right)^m
    -m\delta_b(\tau,\vartheta)
    -\delta_{\mathrm{est}}(\tau,\vartheta)
    \right].
  \]
  若 \(S_\varepsilon\) 未知，可用经验低能量精英集合
  \[
    \widehat S_q
    \DefEq
    \{x:E(x)\le q\text{-分位能量}\}
  \]
  替代，得到经验目标 \(\max_{\tau,\vartheta}\pi_{\tau,\vartheta}(\widehat S_q)\)。
  \item 为目标分布 \(\pi_{\tau^\star,\vartheta^\star}\) 构造采样核
  \[
    P_{\tau^\star,\vartheta^\star},
    \qquad
    \pi_{\tau^\star,\vartheta^\star}P_{\tau^\star,\vartheta^\star}
    =
    \pi_{\tau^\star,\vartheta^\star},
  \]
  并要求 \(|\lambda_1|\) 尽量小，或谱隙 \(1-|\lambda_1|\) 尽量大。
  \item 运行采样器得到 \(X_1,\ldots,X_m\)，输出
  \[
    \hat x=\arg\min_{1\le i\le m}E(X_i).
  \]
  若
  \[
    \pi_{\tau^\star,\vartheta^\star}(S_\varepsilon)-\delta_b\ge p_0,
  \]
  则命中概率具有保守下界
  \[
    \Pr(\hat x\in S_\varepsilon)\ge1-e^{-mp_0}.
  \]
\end{enumerate}
\end{remark}

\begin{remark}[转化意义]
第一，该转化把“哪个局部搜索算法更好”改写为
\[
  \text{哪个分布更容易把概率质量压到 }S_\varepsilon\text{ 上}.
\]
即
\[
  \boxed{
  \text{优选局部算法}
  \Longrightarrow
  \text{优选平稳分布}
  }.
\]

第二，整体遍历质量可拆成两个可计算量：
\[
  \text{命中质量}+\text{混合质量}.
\]
命中质量是
\[
  \pi_{\tau,\vartheta}(S_\varepsilon),
\]
混合质量由谱控制：
\[
  |\lambda_1|^b.
\]
所以好算法的结构是
\[
  \boxed{
  \text{高最优质量分布}
  +
  \text{快混合采样动力学}
  }.
\]

第三，频域共振分布律及其命中随机变量簇给出局部盆地之间的全局调制：
\[
  \pi_{\beta,\theta,K}(x)
  \propto
  w(x)e^{-\beta E(x)+\sum_{j=1}^{K}\theta_j\phi_j(x)}.
\]
这里局部邻域 \(N(x)\) 给出时域移动约束，而频域模态 \(\phi_j\) 表达瓶颈、能垒、盆地连接方向与慢混合方向。
\end{remark}

\begin{remark}[必要边界]
本章转化有四个边界。

第一，只优化分布律不等于自动全局最优。
即使
\[
  \pi_{\tau,\vartheta}(S_\varepsilon)
\]
很大，若采样核混合极慢
\[
  |\lambda_1|\approx1,
\]
则有限时间内仍可能无法命中。
因此
\[
  \text{分布律优}
  \not\Rightarrow
  \text{有限时间命中优}.
\]

第二，先验权重不能把最优区排除。
若
\[
  W_\vartheta(x)=0,
  \qquad
  \forall x\in S_\varepsilon,
\]
则
\[
  \pi_{\tau,\vartheta}(S_\varepsilon)=0.
\]
因此非均匀分布必须满足
\[
  W_\vartheta(S_\varepsilon)>0.
\]

第三，频域模态必须能表达低能量区域。
若
\[
  \langle\mathbf 1_{S_\varepsilon},\phi_j\rangle\approx0,
  \qquad
  j=1,\ldots,K,
\]
则频域共振参数项
\[
  \sum_{j=1}^{K}\theta_j\phi_j(x)
\]
难以提高 \(\pi_{\beta,\theta,K}(S_\varepsilon)\)。
所以需要
\[
  \Pi_K\mathbf 1_{S_\varepsilon}\approx\mathbf 1_{S_\varepsilon}.
\]

第四，过低温会提高集中性但降低遍历性。
增大 \(\beta\) 通常提高低能量区域概率质量，但过大的 \(\beta\) 也会导致采样停滞并增大
\[
  \delta_{\mathrm{mix}}.
\]
所以温度优化不是 \(\beta\to\infty\) 的单向结论，而是
\[
  \boxed{
  \beta^\star
  \in
  \arg\max_{\beta}
  \left[
  \pi_{\beta,w}(S_\varepsilon)-\delta_{\mathrm{mix}}(\beta)
  \right].
  }
\]
\end{remark}

\begin{remark}[本章最终口径]
本章最终签发的严格表达是：
\[
  \boxed{
  \max_{a\in\mathcal A_{\mathrm{local}}}
  \Pr_a(\exists t\le T:X_t\in S_\varepsilon)
  }
\]
可转化为
\[
  \boxed{
  \max_{\tau,\vartheta}
  \left[
  1-
  \left(
  1-\pi_{\tau,\vartheta}(S_\varepsilon)
  \right)^m
  -
  m\delta_{\mathrm{mix}}(\tau,\vartheta)
  -
  \delta_{\mathrm{est}}(\tau,\vartheta)
  \right].
  }
\]
当混合误差与估计误差可控时，进一步简化为
\[
  \boxed{
  \max_{\tau,\vartheta}
  \pi_{\tau,\vartheta}(S_\varepsilon).
  }
\]
因此，该思想可以概括为：
\[
  \boxed{
  \begin{gathered}
  \text{把时域中选择哪个局部算法的问题，}\\
  \text{转化为频域中选择哪类统计力学分布及其参数的问题。}
  \end{gathered}
  }
\]
但严格版本必须保留
\[
  \boxed{
  \text{分布律优化}
  +
  \text{谱混合优化}
  +
  \text{有限样本误差控制}.
  }
\]
因此，“频域分布律质量高”与“有限时间实际命中高”之间必须保留混合误差与估计误差通道。
\end{remark}

\clearpage

%%%%%%%%%%%%%%%%%%%%%%%%%%%%%%%%%%%%%%%%%%%%%%%%%%%%%%%%%%%%
\section{形式逻辑：频域共振的随机变量簇机制}
\label{sec:tex55-resonance-random-variable-cluster}
%%%%%%%%%%%%%%%%%%%%%%%%%%%%%%%%%%%%%%%%%%%%%%%%%%%%%%%%%%%%

\begin{remark}[严格对象口径]
本文统一采用“分布律、状态随机变量、命中指示变量”三层对象口径。
频域共振采样由分布律参数索引，真正承载每次命中判定的是随机变量簇。
定义为：
\[
  \boxed{
  \text{频域共振}
  =
  \text{每次命中判定上的、由分布类型与参数索引的随机变量簇机制}.
  }
\]
也就是说，频域共振的严格对象是
\[
  \boxed{
  \mathfrak R
  \DefEq
  \left\{
  X_i^{\lambda}:i\in\NN,\ \lambda\in\Lambda
  \right\},
  \qquad
  \lambda=(\tau,\vartheta).
  }
\]
其中 \(X_i^\lambda\) 是第 \(i\) 次采样、尝试或命中判定对应的随机变量；
\(\pi_\lambda\) 只是它的分布律。
因此，全文中“优化 \(\pi_{\tau,\vartheta}(S_\varepsilon)\)”的分布律简写，
都应理解为“优化由 \(\pi_{\tau,\vartheta}\) 索引的命中指示随机变量 \(Y_i^{(\tau,\vartheta)}\) 的命中概率”。
\end{remark}

\subsection{对象域：分布律索引、状态随机变量与命中指示}
\label{subsec:tex55-resonance-rv-objects}

\begin{definition}[分布类型与参数空间]
\label{def:tex55-law-index-space}
令
\[
  \Lambda
  \DefEq
  \{(\tau,\vartheta)\}
\]
为分布律索引空间。
其中 \(\tau\) 表示分布类型，例如
\[
  \text{Gibbs},
  \quad
  \text{退火 Gibbs},
  \quad
  \text{频域共振 Gibbs 分布律},
  \quad
  \text{混合分布},
  \quad
  \text{重尾分布}.
\]
\(\vartheta\) 表示参数，例如
\[
  \vartheta=(\beta,w,K,\theta_1,\ldots,\theta_K),
\]
其中 \(\beta\) 是逆温度，\(w\) 是均匀或非均匀先验权重，\(K\) 是保留的频域模态数，
\(\theta_j\) 是第 \(j\) 个频域模态的共振参数。
\end{definition}

\begin{definition}[频域共振分布律]
\label{def:tex55-resonance-law-corrected}
给定
\[
  \lambda=(\tau,\vartheta)\in\Lambda,
\]
定义分布律
\[
  \pi_\lambda(x)\DefEq\pi_{\tau,\vartheta}(x).
\]
典型频域共振 Gibbs 分布律为
\[
  \pi_{\beta,\theta,K}(x)
  =
  \frac{
  w(x)
  \exp\left(
  -\beta E(x)+\sum_{j=1}^{K}\theta_j\phi_j(x)
  \right)
  }{
  Z_{\beta,\theta,K}
  }.
\]
其中 \(\phi_1,\ldots,\phi_K\) 是频域模态。
严格地说，\(\pi_\lambda\) 是随机变量的\textbf{分布律}，不是随机变量本身。
\end{definition}

\begin{definition}[每次命中判定的状态随机变量簇]
\label{def:tex55-state-rv-cluster}
令
\[
  (\Xi,\mathcal F,\mathbb P)
\]
为基础概率空间。
对每次尝试
\[
  i=1,2,\ldots
\]
以及每个分布律索引
\[
  \lambda=(\tau,\vartheta)\in\Lambda,
\]
定义状态随机变量
\[
  X_i^\lambda:\Xi\to\Omega.
\]
若
\[
  \mathcal L(X_i^\lambda)=\pi_\lambda,
\]
也就是
\[
  \mathbb P(X_i^\lambda=x)=\pi_\lambda(x),
  \qquad
  x\in\Omega,
\]
则频域共振状态层对象为
\[
  \boxed{
  \mathfrak X
  \DefEq
  \{X_i^\lambda:i\in\NN,\ \lambda\in\Lambda\}.
  }
\]
\end{definition}

\begin{definition}[近似全局最优集合、命中事件与命中指示变量]
\label{def:tex55-hit-event-indicator}
令
\[
  E_*\DefEq\min_{x\in\Omega}E(x),
  \qquad
  \varepsilon\ge0.
\]
定义近似全局最优集合为
\[
  S_\varepsilon
  \DefEq
  \{x\in\Omega:E(x)\le E_*+\varepsilon\}.
\]
当 \(\varepsilon=0\) 时，\(S_\varepsilon=S_*\)。

对每个随机变量 \(X_i^\lambda\)，定义第 \(i\) 次命中事件为
\[
  \mathcal H_i^\lambda
  \DefEq
  \{X_i^\lambda\in S_\varepsilon\}.
\]
命中事件簇为
\[
  \boxed{
  \mathfrak H
  \DefEq
  \{\mathcal H_i^\lambda:i\in\NN,\ \lambda\in\Lambda\}.
  }
\]

定义命中指示随机变量
\[
  Y_i^\lambda
  \DefEq
  \mathbf 1_{S_\varepsilon}(X_i^\lambda),
\]
即
\[
  Y_i^\lambda=
  \begin{cases}
    1,&X_i^\lambda\in S_\varepsilon,\\
    0,&X_i^\lambda\notin S_\varepsilon.
  \end{cases}
\]
命中指示随机变量簇为
\[
  \boxed{
  \mathfrak Y
  \DefEq
  \{Y_i^\lambda:i\in\NN,\ \lambda\in\Lambda\}.
  }
\]
\end{definition}

\subsection{引理与命题：分布律优化是随机变量簇命中优化的降维表达}
\label{subsec:tex55-law-as-rv-projection}

\begin{lemma}[分布律推出事件概率]
\label{lem:tex55-law-event-probability}
若
\[
  \mathcal L(X_i^\lambda)=\pi_\lambda,
\]
则对任意 \(A\subseteq\Omega\)，有
\[
  \mathbb P(X_i^\lambda\in A)=\pi_\lambda(A).
\]
\end{lemma}

\begin{Certificate}[证书：分布律定义证书]
\label{cert:tex55-law-event-probability}
证书为
\[
  \mathcal L(X_i^\lambda)=\pi_\lambda
  \Longleftrightarrow
  \forall A\subseteq\Omega,\quad
  \mathbb P((X_i^\lambda)^{-1}(A))=\pi_\lambda(A).
\]
\end{Certificate}

\begin{proof}[证明（基于数学符号推演逻辑的谓词因果链）]
由随机变量分布律的定义，
\[
  \mathcal L(X_i^\lambda)(A)
  =
  \mathbb P((X_i^\lambda)^{-1}(A)).
\]
又因为
\[
  \mathcal L(X_i^\lambda)=\pi_\lambda,
\]
所以
\[
  \mathbb P(X_i^\lambda\in A)
  =
  \pi_\lambda(A).
\]
引理得证。
\end{proof}

\begin{proposition}[分布律优化是随机变量簇命中优化的降维表达]
\label{prop:tex55-law-optimization-projection}
若
\[
  \mathcal L(X_i^\lambda)=\pi_\lambda,
\]
且目标事件为
\[
  \mathcal H_i^\lambda=\{X_i^\lambda\in S_\varepsilon\},
\]
则
\[
  \pi_\lambda(S_\varepsilon)
  =
  \mathbb P(\mathcal H_i^\lambda)
  =
  \mathbb P(Y_i^\lambda=1).
\]
因此
\[
  \max_{\lambda\in\Lambda}\pi_\lambda(S_\varepsilon)
\]
只是
\[
  \max_{\lambda\in\Lambda}\mathbb P(Y_i^\lambda=1)
\]
的分布律简写。
\end{proposition}

\begin{Certificate}[证书：事件与指示变量等价]
\label{cert:tex55-law-optimization-projection}
证书由两项组成：
\[
\begin{aligned}
  P_1&:\quad
  \mathbb P(X_i^\lambda\in S_\varepsilon)=\pi_\lambda(S_\varepsilon);\\
  P_2&:\quad
  Y_i^\lambda=1
  \Longleftrightarrow
  X_i^\lambda\in S_\varepsilon.
\end{aligned}
\]
\end{Certificate}

\begin{proof}[证明（基于数学符号推演逻辑的谓词因果链）]
由引理~\ref{lem:tex55-law-event-probability}，取
\[
  A=S_\varepsilon,
\]
得到
\[
  \mathbb P(X_i^\lambda\in S_\varepsilon)
  =
  \pi_\lambda(S_\varepsilon).
\]
又由命中指示变量定义，
\[
  Y_i^\lambda=\mathbf 1_{S_\varepsilon}(X_i^\lambda),
\]
所以
\[
  Y_i^\lambda=1
  \Longleftrightarrow
  X_i^\lambda\in S_\varepsilon.
\]
因此
\[
  \mathbb P(Y_i^\lambda=1)
  =
  \mathbb P(X_i^\lambda\in S_\varepsilon)
  =
  \pi_\lambda(S_\varepsilon).
\]
命题得证。
\end{proof}

\subsection{公理降级：多次命中判定的归纳公式}
\label{subsec:tex55-rv-hit-induction}

\begin{definition}[单次与多次命中泛函]
\label{def:tex55-hit-functional}
定义单次命中泛函为
\[
  \mathcal J_1(\lambda)
  \DefEq
  \mathbb E[Y_1^\lambda]
  =
  \mathbb P(Y_1^\lambda=1).
\]
由于 \(Y_1^\lambda\in\{0,1\}\)，由命题~\ref{prop:tex55-law-optimization-projection}，
\[
  \mathcal J_1(\lambda)=\pi_\lambda(S_\varepsilon).
\]

定义多次命中泛函为
\[
  \mathcal J_m(\lambda)
  \DefEq
  \mathbb P\left(
  \sum_{i=1}^{m}Y_i^\lambda\ge1
  \right).
\]
\end{definition}

\begin{axiom}[公理（降级为定理）：独立同分布命中泛函归纳公式]
\label{ax:tex55-rv-iid-hit-functional}
若
\[
  X_1^\lambda,\ldots,X_m^\lambda
  \quad\text{独立同分布，且}\quad
  \mathcal L(X_i^\lambda)=\pi_\lambda,
\]
并令
\[
  p_\lambda\DefEq\pi_\lambda(S_\varepsilon),
\]
则
\[
  \mathcal J_m(\lambda)
  =
  1-(1-p_\lambda)^m.
\]
\end{axiom}

\begin{Certificate}[数学归纳法证书：独立命中指示变量乘法链]
\label{cert:tex55-rv-iid-hit-functional}
定义谓词
\[
  P(m):\quad
  \mathcal J_m(\lambda)=1-(1-p_\lambda)^m.
\]
证书为
\[
  \pi_{\mathrm{rv-hit}}
  =
  \bigl(P(1),\ \forall m\ge1,\ P(m)\Rightarrow P(m+1)\bigr).
\]
\end{Certificate}

\begin{proof}[证明（基于数学符号推演逻辑的谓词因果链）]
基步 \(m=1\)：
\[
  \mathcal J_1(\lambda)
  =
  \mathbb P(Y_1^\lambda=1)
  =
  p_\lambda
  =
  1-(1-p_\lambda)^1.
\]
故 \(P(1)\) 成立。

归纳步：假设 \(P(m)\) 成立。
由于
\[
  \sum_{i=1}^{m+1}Y_i^\lambda\ge1
\]
的补事件是
\[
  Y_1^\lambda=0,\ldots,Y_{m+1}^\lambda=0,
\]
且 \(Y_i^\lambda\) 独立同分布，得到
\[
\begin{aligned}
  \mathcal J_{m+1}(\lambda)
  &=
  1-\mathbb P(Y_1^\lambda=0,\ldots,Y_{m+1}^\lambda=0)\\
  &=
  1-\prod_{i=1}^{m+1}\mathbb P(Y_i^\lambda=0)\\
  &=
  1-(1-p_\lambda)^{m+1}.
\end{aligned}
\]
故 \(P(m+1)\) 成立。
由数学归纳法，
\[
  \forall m\ge1,\quad
  \mathcal J_m(\lambda)=1-(1-p_\lambda)^m.
\]
证毕。
\end{proof}

\subsection{定理：实际采样、混合误差与自适应随机变量簇}
\label{subsec:tex55-adaptive-rv-theorems}

\begin{definition}[带混合误差的实际命中随机变量]
\label{def:tex55-mixing-rv}
若第 \(i\) 次样本并非严格服从目标律 \(\pi_\lambda\)，而是服从实际律
\[
  \mu_{i,\lambda},
\]
则
\[
  \mathcal L(X_i^\lambda)=\mu_{i,\lambda}.
\]
此时单次命中概率为
\[
  p_{i,\lambda}
  \DefEq
  \mu_{i,\lambda}(S_\varepsilon).
\]
若
\[
  \|\mu_{i,\lambda}-\pi_\lambda\|_{\TV}
  \le
  \delta_{i,\lambda},
\]
则
\[
  |p_{i,\lambda}-\pi_\lambda(S_\varepsilon)|
  \le
  \delta_{i,\lambda}.
\]
因此实际命中随机变量是
\[
  Y_i^\lambda
  =
  \mathbf 1_{S_\varepsilon}(X_i^\lambda),
  \qquad
  X_i^\lambda\sim\mu_{i,\lambda},
\]
而理想分布律 \(\pi_\lambda\) 只是其目标律。
\end{definition}

\begin{theorem}[自适应参数命中公式]
\label{thm:tex55-adaptive-hit}
若每次命中尝试允许使用不同参数
\[
  \lambda_i=(\tau_i,\vartheta_i),
\]
并且
\[
  X_i^{\lambda_i}\sim\pi_{\lambda_i}
\]
相互独立，则
\[
  \mathbb P(\exists i\le m:X_i^{\lambda_i}\in S_\varepsilon)
  =
  1-
  \prod_{i=1}^{m}
  (1-p_{\lambda_i}),
\]
其中
\[
  p_{\lambda_i}\DefEq\pi_{\lambda_i}(S_\varepsilon).
\]
于是自适应但独立的优化问题为
\[
  \boxed{
  \max_{\lambda_1,\ldots,\lambda_m}
  \left[
  1-
  \prod_{i=1}^{m}
  (1-\pi_{\lambda_i}(S_\varepsilon))
  \right].
  }
\]
\end{theorem}

\begin{Certificate}[证书：非同分布独立未命中乘积]
\label{cert:tex55-adaptive-hit}
证书为
\[
  \mathbb P(X_i^{\lambda_i}\notin S_\varepsilon)
  =
  1-p_{\lambda_i},
  \qquad
  i=1,\ldots,m,
\]
以及各次尝试的独立性。
\end{Certificate}

\begin{proof}[证明（基于数学符号推演逻辑的谓词因果链）]
至少一次命中的补事件为
\[
  \bigcap_{i=1}^{m}\{X_i^{\lambda_i}\notin S_\varepsilon\}.
\]
由独立性，
\[
\begin{aligned}
  \mathbb P(\forall i\le m,\ X_i^{\lambda_i}\notin S_\varepsilon)
  &=
  \prod_{i=1}^{m}
  \mathbb P(X_i^{\lambda_i}\notin S_\varepsilon)\\
  &=
  \prod_{i=1}^{m}(1-p_{\lambda_i}).
\end{aligned}
\]
取补事件得到
\[
  \mathbb P(\exists i\le m:X_i^{\lambda_i}\in S_\varepsilon)
  =
  1-\prod_{i=1}^{m}(1-p_{\lambda_i}).
\]
定理得证。
\end{proof}

\begin{theorem}[非线性遍历的条件命中公式]
\label{thm:tex55-nonlinear-conditional-hit}
在非线性遍历中，令参数由历史决定：
\[
  \lambda_i=\Gamma_i(\mathcal F_{i-1}),
  \qquad
  \mathcal F_{i-1}\DefEq\sigma(X_1,\ldots,X_{i-1}).
\]
若
\[
  X_i\mid\mathcal F_{i-1}\sim\pi_{\lambda_i},
  \qquad
  Y_i\DefEq\mathbf 1_{S_\varepsilon}(X_i),
\]
并定义
\[
  h_i
  \DefEq
  \mathbb P(Y_i=1\mid Y_1=0,\ldots,Y_{i-1}=0),
\]
则
\[
  \boxed{
  \mathbb P(\exists i\le m:Y_i=1)
  =
  1-
  \prod_{i=1}^{m}(1-h_i).
  }
\]
因此非线性遍历的严格目标是
\[
  \boxed{
  \max_{\Gamma_1,\ldots,\Gamma_m}
  \left[
  1-
  \prod_{i=1}^{m}(1-h_i)
  \right].
  }
\]
\end{theorem}

\begin{Certificate}[证书：条件未命中链式分解]
\label{cert:tex55-nonlinear-conditional-hit}
证书为链式分解
\[
  \mathbb P(Y_1=0,\ldots,Y_m=0)
  =
  \prod_{i=1}^{m}
  \mathbb P(Y_i=0\mid Y_1=0,\ldots,Y_{i-1}=0),
\]
以及
\[
  \mathbb P(Y_i=0\mid Y_1=0,\ldots,Y_{i-1}=0)=1-h_i.
\]
\end{Certificate}

\begin{proof}[证明（基于数学符号推演逻辑的谓词因果链）]
至少一次命中事件的补事件是
\[
  \{Y_1=0,\ldots,Y_m=0\}.
\]
由条件概率链式法则，
\[
\begin{aligned}
  \mathbb P(Y_1=0,\ldots,Y_m=0)
  &=
  \prod_{i=1}^{m}
  \mathbb P(Y_i=0\mid Y_1=0,\ldots,Y_{i-1}=0)\\
  &=
  \prod_{i=1}^{m}(1-h_i).
\end{aligned}
\]
取补事件，得到
\[
  \mathbb P(\exists i\le m:Y_i=1)
  =
  1-\prod_{i=1}^{m}(1-h_i).
\]
定理得证。
\end{proof}

\subsection{推论：频域共振随机变量簇定义与全篇统一口径}
\label{subsec:tex55-resonance-correction-corollary}

\begin{corollary}[频域共振的随机变量簇定义]
\label{cor:tex55-resonance-corrected-definition}
频域共振不应被定义为单个分布函数
\[
  \pi_{\beta,\theta,K}(x),
\]
而应定义为三层对象：
\[
  \boxed{
  \mathsf{Res}
  \DefEq
  \left(
  \{X_i^{\beta,\theta,K}\}_{i\ge1},
  \{Y_i^{\beta,\theta,K}\}_{i\ge1},
  \{\mathcal L(X_i^{\beta,\theta,K})=\pi_{\beta,\theta,K}\}_{i\ge1}
  \right).
  }
\]
其中
\[
  X_i^{\beta,\theta,K}\sim\pi_{\beta,\theta,K},
  \qquad
  Y_i^{\beta,\theta,K}
  =
  \mathbf 1_{S_\varepsilon}(X_i^{\beta,\theta,K}).
\]
因此
\[
  \boxed{
  \begin{gathered}
  \text{频域共振}\\
  =
  \text{由频域参数控制分布律、由分布律生成每次命中随机变量、}\\
  \text{由命中指示变量评价命中事件的随机变量簇机制。}
  \end{gathered}
  }
\]
\end{corollary}

\begin{Certificate}[证书：分布律、状态随机变量、命中指示三层闭合]
\label{cert:tex55-resonance-corrected-definition}
证书为
\[
\begin{aligned}
  C_1&:\quad \pi_{\beta,\theta,K}\text{ 是分布律};\\
  C_2&:\quad X_i^{\beta,\theta,K}\sim\pi_{\beta,\theta,K}\text{ 是状态随机变量};\\
  C_3&:\quad Y_i^{\beta,\theta,K}=\mathbf 1_{S_\varepsilon}(X_i^{\beta,\theta,K})\text{ 是命中指示变量};\\
  C_4&:\quad
  \mathbb P(Y_i^{\beta,\theta,K}=1)=\pi_{\beta,\theta,K}(S_\varepsilon).
\end{aligned}
\]
\end{Certificate}

\begin{proof}[证明（基于数学符号推演逻辑的谓词因果链）]
由 \(C_1\)，\(\pi_{\beta,\theta,K}\) 只确定概率法律。
由 \(C_2\)，每次命中判定需要一个由该法律生成的状态随机变量
\[
  X_i^{\beta,\theta,K}.
\]
由 \(C_3\)，命中事件通过指示变量
\[
  Y_i^{\beta,\theta,K}
  =
  \mathbf 1_{S_\varepsilon}(X_i^{\beta,\theta,K})
\]
表达。
由命题~\ref{prop:tex55-law-optimization-projection}，
\[
  \mathbb P(Y_i^{\beta,\theta,K}=1)
  =
  \pi_{\beta,\theta,K}(S_\varepsilon).
\]
因此，\(\pi_{\beta,\theta,K}\) 是频域共振机制的分布律层，
\(X_i^{\beta,\theta,K}\) 是状态随机变量层，
\(Y_i^{\beta,\theta,K}\) 是命中判定层。
三层合在一起才是严格的频域共振对象。
推论得证。
\end{proof}

\begin{remark}[统一后的优化表达]
分布律质量写法
\[
  \max_{\tau,\vartheta}\pi_{\tau,\vartheta}(S_\varepsilon)
\]
应严格读作
\[
  \boxed{
  \max_{\tau,\vartheta}
  \mathbb P\left(
  Y_i^{(\tau,\vartheta)}=1
  \right)
  }.
\]
多次命中版本应写为
\[
  \boxed{
  \max_{\lambda_1,\ldots,\lambda_m}
  \mathbb P\left(
  \sum_{i=1}^{m}Y_i^{\lambda_i}\ge1
  \right).
  }
\]
其中
\[
  Y_i^{\lambda_i}
  =
  \mathbf 1_{S_\varepsilon}(X_i^{\lambda_i}),
  \qquad
  X_i^{\lambda_i}\sim\pi_{\lambda_i}.
\]
独立同分布情形只是特殊退化：
\[
  \mathbb P\left(
  \sum_{i=1}^{m}Y_i^\lambda\ge1
  \right)
  =
  1-(1-\pi_\lambda(S_\varepsilon))^m.
\]
因此，只有在独立同分布、混合误差可控、估计误差可控时，
才可以把随机变量簇命中泛函简写为分布律质量优化。
\end{remark}

\begin{remark}[本章统一口径]
本章采用的统一结论是：
\[
  \boxed{
  \text{时域局部算法优选问题}
  \Longrightarrow
  \text{频域参数化随机变量簇的命中泛函优化问题}.
  }
\]
该命题强于单纯的
\[
  \text{单层分布函数优化问题}.
\]
分布律 \(\pi_\lambda\) 只是随机变量簇的法律；真正承载“每次命中”的对象是
\[
  \boxed{
  X_i^\lambda
  \quad\text{与}\quad
  Y_i^\lambda.
  }
\]
这把对象层级统一为“分布律索引的随机变量簇机制”，也是全篇关于频域共振表述的统一口径。
\end{remark}

\clearpage

%%%%%%%%%%%%%%%%%%%%%%%%%%%%%%%%%%%%%%%%%%%%%%%%%%%%%%%%%%%%
\section{形式逻辑：同伦扭曲轨迹函数与首次命中泛函优化}
\label{sec:tex55-homotopy-twist-hitting-functional}
%%%%%%%%%%%%%%%%%%%%%%%%%%%%%%%%%%%%%%%%%%%%%%%%%%%%%%%%%%%%

\begin{remark}[严格命中口径]
本文把同伦扭曲的最快命中目标统一表述为：
\[
  \boxed{
  \text{使同伦扭曲诱导的某一随机变量取值落入命中集合。}
  }
\]
严格的对象链条是
\[
  \boxed{
  \begin{gathered}
  \text{同伦扭曲}
  \Longrightarrow
  \text{分布律轨迹}
  \Longrightarrow
  \text{随机变量簇}\\
  \Longrightarrow
  \text{命中指示变量簇}
  \Longrightarrow
  \text{首次命中时间}
  \Longrightarrow
  \text{命中泛函}.
  \end{gathered}
  }
\]
因此，同伦扭曲算法的优化方向应写成：
\[
  \boxed{
  \begin{gathered}
  \text{同伦扭曲算法的优化方向}\\
  =
  \text{优选命中泛函过程的轨迹函数，使首次命中时间尽可能小。}
  \end{gathered}
  }
\]
\end{remark}

\subsection{对象域：命中集合、分布律与同伦扭曲轨迹}
\label{subsec:tex55-homtwist-objects}

\begin{definition}[状态空间与命中集合]
\label{def:tex55-homtwist-state-hitset}
令
\[
  \Omega=\{x_1,\ldots,x_N\}
\]
为有限状态空间，令
\[
  E:\Omega\to\RR
\]
为能量函数或目标函数。
定义
\[
  E_*\DefEq\min_{x\in\Omega}E(x),
  \qquad
  S_\varepsilon
  \DefEq
  \{x\in\Omega:E(x)\le E_*+\varepsilon\},
  \qquad
  \varepsilon\ge0.
\]
命中事件定义为随机变量取值满足
\[
  X\in S_\varepsilon.
\]
\end{definition}

\begin{definition}[分布律与参数空间]
\label{def:tex55-homtwist-law-space}
令 \(\Lambda\) 为分布类型与参数的索引空间。
每个
\[
  \lambda\in\Lambda
\]
代表一个分布律
\[
  \pi_\lambda.
\]
例如频域统计力学分布律可写为
\[
  \pi_\lambda(x)
  =
  \frac{
  w_\lambda(x)
  \exp\left(
  -\beta_\lambda E(x)
  +
  \sum_{j=1}^{K_\lambda}\theta_{\lambda,j}\phi_j(x)
  \right)
  }{
  Z_\lambda
  },
\]
其中
\[
  \lambda
  =
  (\beta_\lambda,w_\lambda,K_\lambda,
  \theta_{\lambda,1},\ldots,\theta_{\lambda,K_\lambda}).
\]
\end{definition}

\begin{definition}[同伦扭曲轨迹函数]
\label{def:tex55-homtwist-trajectory}
所谓同伦扭曲在本章中定义为一条参数轨迹。
非自适应版本为
\[
  \gamma:\{1,\ldots,T\}\to\Lambda,
  \qquad
  t\mapsto\lambda_t=\gamma(t).
\]
自适应版本为
\[
  \Gamma_t:\mathcal F_{t-1}\to\Lambda,
  \qquad
  \mathcal F_{t-1}\DefEq\sigma(X_1,\ldots,X_{t-1}),
\]
并令
\[
  \lambda_t=\Gamma_t(\mathcal F_{t-1}).
\]
因此同伦扭曲算法严格应理解为轨迹函数族
\[
  \boxed{
  \Gamma=(\Gamma_1,\ldots,\Gamma_T),
  }
\]
这使算法对象落在分布律参数空间的轨迹层。
\end{definition}

\begin{definition}[同伦扭曲诱导的状态随机变量簇]
\label{def:tex55-homtwist-state-rv-cluster}
给定同伦扭曲轨迹函数 \(\Gamma\)，第 \(t\) 次尝试产生状态随机变量
\[
  X_t^\Gamma:\Xi\to\Omega.
\]
其条件分布为
\[
  \mathcal L(X_t^\Gamma\mid\mathcal F_{t-1})
  =
  \pi_{\lambda_t},
  \qquad
  \lambda_t=\Gamma_t(\mathcal F_{t-1}).
\]
因此同伦扭曲形成的状态随机变量簇为
\[
  \boxed{
  \mathfrak X_\Gamma
  \DefEq
  \{X_t^\Gamma:t=1,\ldots,T\}.
  }
\]
若把所有可能参数纳入，则得到更大的母簇
\[
  \boxed{
  \mathfrak X
  \DefEq
  \{X_t^\lambda:t\in\NN,\ \lambda\in\Lambda\},
  \qquad
  \mathcal L(X_t^\lambda)=\pi_\lambda.
  }
\]
\end{definition}

\begin{definition}[命中指示随机变量簇]
\label{def:tex55-homtwist-hit-indicator}
对每个 \(X_t^\Gamma\)，定义命中指示变量
\[
  Y_t^\Gamma
  \DefEq
  \mathbf 1_{S_\varepsilon}(X_t^\Gamma).
\]
即
\[
  Y_t^\Gamma=
  \begin{cases}
    1,&X_t^\Gamma\in S_\varepsilon,\\
    0,&X_t^\Gamma\notin S_\varepsilon.
  \end{cases}
\]
命中指示随机变量簇为
\[
  \boxed{
  \mathfrak Y_\Gamma
  \DefEq
  \{Y_t^\Gamma:t=1,\ldots,T\}.
  }
\]
这里的命中严格对应
\[
  \exists t,\quad Y_t^\Gamma=1,
\]
也就是
\[
  \exists t,\quad X_t^\Gamma\in S_\varepsilon.
\]
\end{definition}

\subsection{首次命中时间与命中泛函}
\label{subsec:tex55-homtwist-hitting-time-functional}

\begin{definition}[首次命中时间]
\label{def:tex55-homtwist-first-hit-time}
定义首次命中时间为
\[
  \tau_{S_\varepsilon}^\Gamma
  \DefEq
  \inf\{t\ge1:X_t^\Gamma\in S_\varepsilon\}.
\]
等价地，
\[
  \tau_{S_\varepsilon}^\Gamma
  =
  \inf\{t\ge1:Y_t^\Gamma=1\}.
\]
若直到 \(T\) 步都没有命中，则定义截断版本
\[
  \tau_{S_\varepsilon,T}^\Gamma
  \DefEq
  \min(\tau_{S_\varepsilon}^\Gamma,T+1).
\]
因此“最快命中”的严格说法是：
\[
  \boxed{
  \text{使首次出现 }Y_t^\Gamma=1\text{ 的时间 }\tau_{S_\varepsilon}^\Gamma\text{ 尽可能小。}
  }
\]
\end{definition}

\begin{definition}[有限时间命中泛函与最快命中泛函]
\label{def:tex55-homtwist-hit-functional}
定义有限时间命中泛函为
\[
  \mathcal J_T(\Gamma)
  \DefEq
  \mathbb P_\Gamma(\tau_{S_\varepsilon}^\Gamma\le T).
\]
等价地，
\[
  \mathcal J_T(\Gamma)
  =
  \mathbb P_\Gamma
  \left(
  \sum_{t=1}^{T}Y_t^\Gamma\ge1
  \right),
\]
也可写为
\[
  \mathcal J_T(\Gamma)
  =
  \mathbb E_\Gamma
  \left[
  \mathbf 1
  \left\{
  \exists t\le T:X_t^\Gamma\in S_\varepsilon
  \right\}
  \right].
\]

若优化目标是“最快”，定义
\[
  \mathcal T(\Gamma)
  \DefEq
  \mathbb E_\Gamma[\tau_{S_\varepsilon,T}^\Gamma].
\]
于是算法优化可写为
\[
  \Gamma^\star
  \in
  \arg\min_{\Gamma}\mathcal T(\Gamma),
\]
或在固定时间预算 \(T\) 下写为
\[
  \Gamma^\star
  \in
  \arg\max_{\Gamma}\mathcal J_T(\Gamma).
\]
\end{definition}

\subsection{引理与命题：轨迹函数诱导命中泛函}
\label{subsec:tex55-homtwist-propositions}

\begin{lemma}[命中事件等价于命中指示变量取值为 \(1\)]
\label{lem:tex55-homtwist-hit-equivalence}
对任意 \(t\le T\)，有
\[
  X_t^\Gamma\in S_\varepsilon
  \Longleftrightarrow
  Y_t^\Gamma=1.
\]
因此
\[
  \{\tau_{S_\varepsilon}^\Gamma\le T\}
  \Longleftrightarrow
  \{\exists t\le T,\ Y_t^\Gamma=1\}.
\]
\end{lemma}

\begin{Certificate}[证书：指示函数定义]
\label{cert:tex55-homtwist-hit-equivalence}
证书为
\[
  Y_t^\Gamma=\mathbf 1_{S_\varepsilon}(X_t^\Gamma)
\]
与首次命中时间定义
\[
  \tau_{S_\varepsilon}^\Gamma
  =
  \inf\{t\ge1:X_t^\Gamma\in S_\varepsilon\}.
\]
\end{Certificate}

\begin{proof}[证明（基于数学符号推演逻辑的谓词因果链）]
由命中指示变量定义，
\[
  Y_t^\Gamma=1
  \Longleftrightarrow
  \mathbf 1_{S_\varepsilon}(X_t^\Gamma)=1
  \Longleftrightarrow
  X_t^\Gamma\in S_\varepsilon.
\]
由首次命中时间定义，
\[
  \tau_{S_\varepsilon}^\Gamma\le T
\]
当且仅当存在 \(t\le T\) 使
\[
  X_t^\Gamma\in S_\varepsilon.
\]
再由前一等价式，得到
\[
  \tau_{S_\varepsilon}^\Gamma\le T
  \Longleftrightarrow
  \exists t\le T,\quad Y_t^\Gamma=1.
\]
引理得证。
\end{proof}

\begin{proposition}[同伦扭曲诱导命中泛函对应的随机变量簇]
\label{prop:tex55-homtwist-induces-functional}
给定同伦扭曲轨迹函数
\[
  \Gamma=(\Gamma_1,\ldots,\Gamma_T),
\]
其诱导分布律轨迹为
\[
  \pi_{\lambda_1},\ldots,\pi_{\lambda_T}.
\]
该分布律轨迹诱导状态随机变量簇
\[
  \mathfrak X_\Gamma=\{X_t^\Gamma\}_{t=1}^{T}.
\]
命中泛函
\[
  \mathcal J_T(\Gamma)
\]
是该随机变量簇，或等价地命中指示簇 \(\mathfrak Y_\Gamma\)，诱导出的泛函。
\end{proposition}

\begin{Certificate}[证书：轨迹到泛函的组合链]
\label{cert:tex55-homtwist-induces-functional}
证书为
\[
  \Gamma
  \Longrightarrow
  (\lambda_1,\ldots,\lambda_T)
  \Longrightarrow
  (X_1^\Gamma,\ldots,X_T^\Gamma)
  \Longrightarrow
  (Y_1^\Gamma,\ldots,Y_T^\Gamma)
  \Longrightarrow
  \mathcal J_T(\Gamma).
\]
\end{Certificate}

\begin{proof}[证明（基于数学符号推演逻辑的谓词因果链）]
由定义~\ref{def:tex55-homtwist-trajectory}，
\[
  \lambda_t=\Gamma_t(\mathcal F_{t-1}).
\]
由定义~\ref{def:tex55-homtwist-state-rv-cluster}，
\[
  \mathcal L(X_t^\Gamma\mid\mathcal F_{t-1})=\pi_{\lambda_t}.
\]
所以每个同伦扭曲参数 \(\lambda_t\) 确定第 \(t\) 次随机变量的条件分布律。
因此
\[
  \Gamma
  \Longrightarrow
  (\lambda_1,\ldots,\lambda_T)
  \Longrightarrow
  (X_1^\Gamma,\ldots,X_T^\Gamma).
\]
由命中指示变量定义，
\[
  Y_t^\Gamma=\mathbf 1_{S_\varepsilon}(X_t^\Gamma),
\]
从而
\[
  (X_1^\Gamma,\ldots,X_T^\Gamma)
  \Longrightarrow
  (Y_1^\Gamma,\ldots,Y_T^\Gamma).
\]
由命中泛函定义，
\[
  \mathcal J_T(\Gamma)
  =
  \mathbb P_\Gamma
  \left(
  \sum_{t=1}^{T}Y_t^\Gamma\ge1
  \right).
\]
所以 \(\mathcal J_T(\Gamma)\) 确实由随机变量簇或命中指示簇诱导。
命题得证。
\end{proof}

\subsection{公理降级：条件命中强度的生存概率递推}
\label{subsec:tex55-homtwist-induction}

\begin{definition}[条件命中强度与生存概率]
\label{def:tex55-homtwist-conditional-intensity}
定义第 \(t\) 步在尚未命中条件下的命中强度为
\[
  h_t^\Gamma
  \DefEq
  \mathbb P_\Gamma
  (
  Y_t^\Gamma=1
  \mid
  \tau_{S_\varepsilon}^\Gamma>t-1
  ).
\]
等价地，
\[
  h_t^\Gamma
  =
  \mathbb P_\Gamma
  (
  X_t^\Gamma\in S_\varepsilon
  \mid
  \tau_{S_\varepsilon}^\Gamma>t-1
  ).
\]
定义生存概率
\[
  Q_t^\Gamma
  \DefEq
  \mathbb P_\Gamma(\tau_{S_\varepsilon}^\Gamma>t).
\]
\end{definition}

\begin{axiom}[公理（降级为定理）：生存概率条件递推闭合]
\label{ax:tex55-homtwist-survival-recursion}
若条件命中强度 \(h_t^\Gamma\) 按定义~\ref{def:tex55-homtwist-conditional-intensity} 给出，
则对任意 \(1\le t\le T\)，
\[
  Q_t^\Gamma
  =
  \prod_{s=1}^{t}(1-h_s^\Gamma).
\]
因此
\[
  \mathcal J_T(\Gamma)
  =
  1-
  \prod_{t=1}^{T}(1-h_t^\Gamma).
\]
\end{axiom}

\begin{Certificate}[数学归纳法证书：条件未命中乘法链]
\label{cert:tex55-homtwist-survival-recursion}
定义谓词
\[
  P(t):\quad
  Q_t^\Gamma=\prod_{s=1}^{t}(1-h_s^\Gamma).
\]
证书为
\[
  \pi_{\mathrm{surv}}
  =
  \bigl(P(1),\ \forall t<T,\ P(t)\Rightarrow P(t+1)\bigr).
\]
\end{Certificate}

\begin{proof}[证明（基于数学符号推演逻辑的谓词因果链）]
基步 \(t=1\)：
\[
  Q_1^\Gamma
  =
  \mathbb P_\Gamma(Y_1^\Gamma=0)
  =
  1-h_1^\Gamma.
\]
故 \(P(1)\) 成立。

归纳步：假设
\[
  Q_t^\Gamma
  =
  \prod_{s=1}^{t}(1-h_s^\Gamma).
\]
由条件概率链式法则，
\[
\begin{aligned}
  Q_{t+1}^\Gamma
  &=
  \mathbb P_\Gamma(\tau_{S_\varepsilon}^\Gamma>t+1)\\
  &=
  \mathbb P_\Gamma(\tau_{S_\varepsilon}^\Gamma>t)
  \mathbb P_\Gamma(
  Y_{t+1}^\Gamma=0
  \mid
  \tau_{S_\varepsilon}^\Gamma>t
  )\\
  &=
  Q_t^\Gamma(1-h_{t+1}^\Gamma)\\
  &=
  \prod_{s=1}^{t+1}(1-h_s^\Gamma).
\end{aligned}
\]
故 \(P(t+1)\) 成立。
由数学归纳法，
\[
  Q_T^\Gamma
  =
  \prod_{t=1}^{T}(1-h_t^\Gamma).
\]
又
\[
  \mathcal J_T(\Gamma)
  =
  \mathbb P_\Gamma(\tau_{S_\varepsilon}^\Gamma\le T)
  =
  1-Q_T^\Gamma,
\]
所以
\[
  \mathcal J_T(\Gamma)
  =
  1-
  \prod_{t=1}^{T}(1-h_t^\Gamma).
\]
证毕。
\end{proof}

\subsection{定理与命题：最快命中的严格含义}
\label{subsec:tex55-homtwist-fastest-hit}

\begin{theorem}[命中泛函由条件命中强度决定]
\label{thm:tex55-homtwist-hit-functional-intensity}
若
\[
  h_t^\Gamma
  =
  \mathbb P_\Gamma(
  Y_t^\Gamma=1
  \mid
  \tau_{S_\varepsilon}^\Gamma>t-1),
\]
则
\[
  \boxed{
  \mathcal J_T(\Gamma)
  =
  1-
  \prod_{t=1}^{T}(1-h_t^\Gamma).
  }
\]
\end{theorem}

\begin{Certificate}[证书：命中泛函与生存概率互补]
\label{cert:tex55-homtwist-hit-functional-intensity}
证书为
\[
  \mathcal J_T(\Gamma)=1-Q_T^\Gamma,
  \qquad
  Q_T^\Gamma=\prod_{t=1}^{T}(1-h_t^\Gamma).
\]
\end{Certificate}

\begin{proof}[证明（基于数学符号推演逻辑的谓词因果链）]
由公理~\ref{ax:tex55-homtwist-survival-recursion}，
\[
  Q_T^\Gamma
  =
  \prod_{t=1}^{T}(1-h_t^\Gamma).
\]
由命中泛函定义，
\[
  \mathcal J_T(\Gamma)
  =
  \mathbb P_\Gamma(\tau_{S_\varepsilon}^\Gamma\le T)
  =
  1-\mathbb P_\Gamma(\tau_{S_\varepsilon}^\Gamma>T)
  =
  1-Q_T^\Gamma.
\]
代入上式即得
\[
  \mathcal J_T(\Gamma)
  =
  1-
  \prod_{t=1}^{T}(1-h_t^\Gamma).
\]
定理得证。
\end{proof}

\begin{proposition}[最快命中等价于压低生存概率]
\label{prop:tex55-homtwist-fast-hit-survival}
对截断首次命中时间
\[
  \tau_{S_\varepsilon,T}^\Gamma
  =
  \min(\tau_{S_\varepsilon}^\Gamma,T+1),
\]
有
\[
  \mathcal T(\Gamma)
  =
  \mathbb E_\Gamma[\tau_{S_\varepsilon,T}^\Gamma]
  =
  \sum_{t=0}^{T}Q_t^\Gamma.
\]
因此，最快命中等价于尽可能压低前缀生存概率序列
\[
  Q_1^\Gamma,Q_2^\Gamma,\ldots,Q_T^\Gamma.
\]
\end{proposition}

\begin{Certificate}[证书：离散尾和公式]
\label{cert:tex55-homtwist-fast-hit-survival}
证书为离散非负整数值随机变量恒等式
\[
  \mathbb E[Z]=\sum_{t\ge0}\mathbb P(Z>t),
\]
其中
\[
  Z=\tau_{S_\varepsilon,T}^\Gamma\in\{1,\ldots,T+1\}.
\]
\end{Certificate}

\begin{proof}[证明（基于数学符号推演逻辑的谓词因果链）]
对
\[
  Z=\tau_{S_\varepsilon,T}^\Gamma
\]
使用离散尾和公式：
\[
  \mathbb E_\Gamma[Z]
  =
  \sum_{t=0}^{T}\mathbb P_\Gamma(Z>t).
\]
当 \(0\le t\le T\) 时，
\[
  Z>t
  \Longleftrightarrow
  \tau_{S_\varepsilon}^\Gamma>t.
\]
因此
\[
  \mathbb P_\Gamma(Z>t)
  =
  \mathbb P_\Gamma(\tau_{S_\varepsilon}^\Gamma>t)
  =
  Q_t^\Gamma.
\]
于是
\[
  \mathcal T(\Gamma)
  =
  \mathbb E_\Gamma[\tau_{S_\varepsilon,T}^\Gamma]
  =
  \sum_{t=0}^{T}Q_t^\Gamma.
\]
所以若要最小化 \(\mathcal T(\Gamma)\)，就要压低各个前缀生存概率 \(Q_t^\Gamma\)。
命题得证。
\end{proof}

\subsection{频域统计力学版本与轨迹成本}
\label{subsec:tex55-homtwist-frequency-version}

\begin{theorem}[频域统计力学轨迹优化的命中下界]
\label{thm:tex55-homtwist-frequency-trajectory}
在频域统计力学版本中，令参数轨迹为
\[
  \lambda_t=(\beta_t,w_t,K_t,\theta_t),
\]
对应分布律为
\[
  \pi_{\lambda_t}(x)
  =
  \frac{
  w_t(x)
  \exp\left(
  -\beta_tE(x)
  +
  \sum_{j=1}^{K_t}\theta_{t,j}\phi_j(x)
  \right)
  }{
  Z_t
  }.
\]
若第 \(t\) 步实际采样分布为 \(\mu_t^\Gamma\)，且
\[
  \|\mu_t^\Gamma-\pi_{\lambda_t}\|_{\TV}\le\delta_t,
\]
则条件命中强度可由
\[
  h_t^\Gamma
  \ge
  \left[
  \pi_{\lambda_t}(S_\varepsilon)-\delta_t
  \right]_+
\]
保守控制。
因此
\[
  \mathcal J_T(\Gamma)
  \ge
  1-
  \prod_{t=1}^{T}
  \left(
  1-
  \left[
  \pi_{\lambda_t}(S_\varepsilon)-\delta_t
  \right]_+
  \right).
\]
\end{theorem}

\begin{Certificate}[证书：全变差误差扣除]
\label{cert:tex55-homtwist-frequency-trajectory}
证书为
\[
  |\mu_t^\Gamma(S_\varepsilon)-\pi_{\lambda_t}(S_\varepsilon)|
  \le
  \|\mu_t^\Gamma-\pi_{\lambda_t}\|_{\TV}
  \le
  \delta_t
\]
以及定理~\ref{thm:tex55-homtwist-hit-functional-intensity}。
\end{Certificate}

\begin{proof}[证明（基于数学符号推演逻辑的谓词因果链）]
由全变差距离定义，对事件 \(S_\varepsilon\) 有
\[
  |\mu_t^\Gamma(S_\varepsilon)-\pi_{\lambda_t}(S_\varepsilon)|
  \le
  \delta_t.
\]
因此实际第 \(t\) 步命中概率满足保守下界
\[
  \mu_t^\Gamma(S_\varepsilon)
  \ge
  \pi_{\lambda_t}(S_\varepsilon)-\delta_t.
\]
由于概率非负，
\[
  h_t^\Gamma
  \ge
  \left[
  \pi_{\lambda_t}(S_\varepsilon)-\delta_t
  \right]_+.
\]
由定理~\ref{thm:tex55-homtwist-hit-functional-intensity}，
\[
  \mathcal J_T(\Gamma)
  =
  1-\prod_{t=1}^{T}(1-h_t^\Gamma).
\]
将 \(h_t^\Gamma\) 的保守下界代入，得到
\[
  \mathcal J_T(\Gamma)
  \ge
  1-
  \prod_{t=1}^{T}
  \left(
  1-
  \left[
  \pi_{\lambda_t}(S_\varepsilon)-\delta_t
  \right]_+
  \right).
\]
定理得证。
\end{proof}

\begin{corollary}[带同伦扭曲成本的轨迹优化]
\label{cor:tex55-homtwist-cost-optimization}
若考虑同伦变形成本、参数跃迁成本或计算成本
\[
  c(\lambda_{t-1},\lambda_t)\ge0,
\]
则轨迹函数优化可写为
\[
  \boxed{
  \Gamma^\star
  \in
  \arg\max_{\Gamma}
  \left\{
  \mathcal J_T(\Gamma)
  -
  \sum_{t=1}^{T}
  c(\lambda_{t-1},\lambda_t)
  \right\}.
  }
\]
\end{corollary}

\begin{Certificate}[证书：命中收益减轨迹成本]
\label{cert:tex55-homtwist-cost-optimization}
证书为收益项
\[
  \mathcal J_T(\Gamma)
\]
与成本项
\[
  C_T(\Gamma)\DefEq\sum_{t=1}^{T}c(\lambda_{t-1},\lambda_t).
\]
\end{Certificate}

\begin{proof}[证明（基于数学符号推演逻辑的谓词因果链）]
若只优化命中概率，则目标为
\[
  \max_\Gamma \mathcal J_T(\Gamma).
\]
当参数扭曲存在成本时，每一步从 \(\lambda_{t-1}\) 到 \(\lambda_t\) 产生代价
\[
  c(\lambda_{t-1},\lambda_t).
\]
总成本为
\[
  C_T(\Gamma)=\sum_{t=1}^{T}c(\lambda_{t-1},\lambda_t).
\]
因此净收益目标为
\[
  \mathcal J_T(\Gamma)-C_T(\Gamma),
\]
对应优化问题即为
\[
  \Gamma^\star
  \in
  \arg\max_{\Gamma}
  \left\{
  \mathcal J_T(\Gamma)
  -
  \sum_{t=1}^{T}
  c(\lambda_{t-1},\lambda_t)
  \right\}.
\]
推论得证。
\end{proof}

\subsection{谓词因果链与最终口径}
\label{subsec:tex55-homtwist-predicate-chain}

\begin{remark}[同伦扭曲的谓词因果链]
本章的谓词因果链为
\[
\begin{aligned}
  &\text{同伦扭曲轨迹函数 }\Gamma\\
  &\Longrightarrow
  \text{分布参数轨迹 }\lambda_t=\Gamma_t(\mathcal F_{t-1})\\
  &\Longrightarrow
  \text{分布律轨迹 }\pi_{\lambda_t}\\
  &\Longrightarrow
  \text{状态随机变量簇 }X_t^\Gamma\\
  &\Longrightarrow
  \text{命中指示变量簇 }
  Y_t^\Gamma=\mathbf 1_{S_\varepsilon}(X_t^\Gamma)\\
  &\Longrightarrow
  \text{首次命中时间 }
  \tau_{S_\varepsilon}^\Gamma=\inf\{t:Y_t^\Gamma=1\}\\
  &\Longrightarrow
  \text{命中泛函 }
  \mathcal J_T(\Gamma)=
  \mathbb P(\tau_{S_\varepsilon}^\Gamma\le T)\\
  &\Longrightarrow
  \text{优化问题 }
  \Gamma^\star\in\arg\max_\Gamma\mathcal J_T(\Gamma),
  \quad
  \text{或 }
  \Gamma^\star\in
  \arg\min_\Gamma\mathbb E[\tau_{S_\varepsilon,T}^\Gamma].
\end{aligned}
\]
\end{remark}

\begin{remark}[最终严格表述]
原句可严格改写为：
\[
  \boxed{
  \text{同伦扭曲形成由分布律及其参数轨迹诱导的随机变量簇。}
  }
\]
\[
  \boxed{
  \begin{gathered}
  \text{命中对应于该随机变量簇中某个 }X_t^\Gamma\text{ 的取值落入 }S_\varepsilon,\\
  \text{等价于 }Y_t^\Gamma=1.
  \end{gathered}
  }
\]
\[
  \boxed{
  \begin{gathered}
  \text{同伦扭曲算法的优化方向，是优选参数轨迹函数 }\Gamma,\\
  \text{使首次命中时间 }\tau_{S_\varepsilon}^\Gamma\text{ 最小，或使命中泛函 }\mathcal J_T(\Gamma)\text{ 最大。}
  \end{gathered}
  }
\]
最压缩的数学表达是
\[
  \boxed{
  \Gamma^\star
  \in
  \arg\min_{\Gamma}
  \mathbb E_\Gamma
  \left[
  \inf\{t\ge1:X_t^\Gamma\in S_\varepsilon\}
  \right].
  }
\]
有限时间版本是
\[
  \boxed{
  \Gamma^\star
  \in
  \arg\max_{\Gamma}
  \mathbb P_\Gamma
  \left(
  \exists t\le T:
  X_t^\Gamma\in S_\varepsilon
  \right).
  }
\]
其中
\[
  X_t^\Gamma\sim\pi_{\Gamma_t(\mathcal F_{t-1})}.
\]
这就是“优选命中泛函过程的轨迹函数”的严格形式。
\end{remark}

%%%%%%%%%%%%%%%%%%%%%%%%%%%%%%%%%%%%%%%%%%%%%%%%%%%%%%%%%%%%
\section{形式逻辑：目标正交同伦扭曲与幂化加速命中}
\label{sec:tex55-orthogonal-twist-power-hitting}
%%%%%%%%%%%%%%%%%%%%%%%%%%%%%%%%%%%%%%%%%%%%%%%%%%%%%%%%%%%%

\begin{remark}[严格加速口径]
本文把目标正交同伦扭曲的加速口径统一为一个可证明命题。
关键限定是：
\[
  \boxed{
  \begin{gathered}
  \text{与目标扭曲正交本身不必然加速命中；}\\
  \text{真正加速的是与目标扭曲正交、但与命中梯度不正交的扭曲方向。}
  \end{gathered}
  }
\]
因此严格口径是：把命中泛函参数化后，在目标扭曲的正交补空间中，
沿命中泛函梯度的正交投影方向做同伦扭曲；
若该投影非零，则命中对数胜算一阶提高，未命中率下降。
若继续把每一步剩余命中梯度投影到前序扭曲方向的正交补中，
并额外满足每层生存概率幂压缩条件，则得到高阶幂化加速命中。
\end{remark}

\subsection{对象域：命中变量、参数化分布律与同伦扭曲}
\label{subsec:tex55-orthogonal-objects}

\begin{definition}[状态空间、命中集合与命中变量]
\label{def:tex55-orthogonal-hit-variable}
令 \(\Omega\) 为有限状态空间，令 \(S\subseteq\Omega\) 为命中集合。
在优化语境中可取
\[
  S=S_\varepsilon
  =
  \{x\in\Omega:E(x)\le E_*+\varepsilon\},
  \qquad
  E_*=\min_{x\in\Omega}E(x).
\]
对第 \(t\) 次尝试，令
\[
  X_t:\Xi\to\Omega
\]
为状态随机变量，并定义命中指示变量
\[
  Y_t=\mathbf 1_S(X_t).
\]
也就是说
\[
  Y_t=
  \begin{cases}
  1,&X_t\in S,\\
  0,&X_t\notin S.
  \end{cases}
\]
因此命中事件严格是
\[
  \{Y_t=1\}=\{X_t\in S\},
\]
这一定义把命中绑定到随机变量取值与集合 \(S\) 的关系。
\end{definition}

\begin{definition}[参数化分布律与命中泛函]
\label{def:tex55-orthogonal-param-law}
令 \(\lambda\in\Lambda\) 表示分布律的类型与参数。
例如在频域统计力学口径下可写成
\[
  \lambda=(\beta_\lambda,w_\lambda,K_\lambda,
  \theta_{\lambda,1},\ldots,\theta_{\lambda,K_\lambda}),
\]
并定义
\[
  \pi_\lambda(x)
  =
  \frac{
  w_\lambda(x)
  \exp\left(
  -\beta_\lambda E(x)
  +
  \sum_{j=1}^{K_\lambda}
  \theta_{\lambda,j}\phi_j(x)
  \right)
  }{
  Z_\lambda
  }.
\]
由此生成随机变量
\[
  X^\lambda\sim\pi_\lambda,
  \qquad
  Y^\lambda=\mathbf 1_S(X^\lambda).
\]
命中泛函定义为
\[
  \mathcal J(\lambda)
  \DefEq
  \mathbb P(X^\lambda\in S)
  =
  \pi_\lambda(S)
  =
  \mathbb E[Y^\lambda].
\]
所以“命中泛函参数化”的严格对象链为
\[
  \boxed{
  \lambda
  \Longrightarrow
  X^\lambda
  \Longrightarrow
  Y^\lambda
  \Longrightarrow
  \mathcal J(\lambda)=\mathbb E[Y^\lambda].
  }
\]
\end{definition}

\begin{definition}[指数同伦扭曲路径]
\label{def:tex55-orthogonal-exponential-twist}
给定当前分布 \(\mu\) 与扭曲方向 \(v:\Omega\to\mathbb R\)，定义指数同伦扭曲路径
\[
  \mu_{\eta,v}(x)
  =
  \frac{\mu(x)e^{\eta v(x)}}{Z_v(\eta)},
  \qquad
  \eta\in[0,1],
\]
其中
\[
  Z_v(\eta)
  =
  \sum_{z\in\Omega}\mu(z)e^{\eta v(z)}.
\]
当 \(\eta=0\) 时，\(\mu_{0,v}=\mu\)；
当 \(\eta=1\) 时，\(\mu_{1,v}\) 是沿 \(v\) 扭曲后的分布。
因此同伦扭曲在此处表示分布律路径
\[
  \eta\mapsto\mu_{\eta,v}.
\]
\end{definition}

\begin{definition}[目标扭曲与正交扭曲]
\label{def:tex55-orthogonal-target}
令 \(u:\Omega\to\mathbb R\) 为目标扭曲方向。
例如 \(u(x)=-E(x)\) 表示能量下降方向，也可取 \(u=\phi_{\mathrm{target}}\) 表示已有谱目标方向。
在当前分布 \(\mu\) 下定义中心化内积
\[
  \langle f,g\rangle_\mu
  \DefEq
  \mathbb E_\mu
  \left[
  (f-\mathbb E_\mu f)(g-\mathbb E_\mu g)
  \right]
  =
  \operatorname{Cov}_\mu(f,g).
\]
若
\[
  \langle v,u\rangle_\mu=0,
\]
则称 \(v\) 与目标扭曲 \(u\) 在分布 \(\mu\) 下正交。
这只表示沿 \(v\) 扭曲时，目标扭曲 \(u\) 的期望一阶不变；
它本身不保证命中泛函上升。
\end{definition}

\begin{definition}[命中对数胜算]
\label{def:tex55-orthogonal-log-odds}
设
\[
  p_\mu=\mu(S),
  \qquad
  q_\mu=\mu(S^c)=1-p_\mu.
\]
在 \(0<p_\mu<1\) 时，定义命中胜算与命中对数胜算为
\[
  O_\mu
  =
  \frac{p_\mu}{q_\mu},
  \qquad
  L_\mu
  =
  \log O_\mu
  =
  \log\frac{\mu(S)}{\mu(S^c)}.
\]
指数扭曲天然改变胜算比，因此用 \(L_\mu\) 表达一阶加速比直接使用 \(p_\mu\) 更清晰。
\end{definition}

\subsection{公理降级：指数扭曲的一阶导数公式}
\label{subsec:tex55-orthogonal-axiom-downgrade}

\begin{theorem}[公理降级：指数扭曲导数公式]
\label{thm:tex55-exponential-twist-derivative}
令 \(\Omega\) 有限，\(\mu(x)>0\)，并令
\[
  \mu_{\eta,v}(x)
  =
  \frac{\mu(x)e^{\eta v(x)}}{Z_v(\eta)}.
\]
则对任意函数 \(f:\Omega\to\mathbb R\)，有
\[
  \left.
  \frac{d}{d\eta}
  \mathbb E_{\mu_{\eta,v}}[f]
  \right|_{\eta=0}
  =
  \operatorname{Cov}_\mu(f,v)
  =
  \langle f,v\rangle_\mu.
\]
\end{theorem}

\begin{Certificate}[符号推演证书]
\label{cert:tex55-exponential-twist-derivative}
本证书把常用的指数族得分公式降级为有限和上的直接计算。
目标谓词为
\[
  P(f):
  \left.
  \frac{d}{d\eta}
  \mathbb E_{\mu_{\eta,v}}[f]
  \right|_{\eta=0}
  =
  \operatorname{Cov}_\mu(f,v).
\]
证明只依赖有限状态空间上的求和、求导与归一化常数
\[
  Z_v(\eta)=\sum_x\mu(x)e^{\eta v(x)}.
\]
\end{Certificate}

\begin{proof}
由定义，
\[
  \mathbb E_{\mu_{\eta,v}}[f]
  =
  \frac{
  \sum_x f(x)\mu(x)e^{\eta v(x)}
  }{
  Z_v(\eta)
  }.
\]
记
\[
  A_f(\eta)=\sum_x f(x)\mu(x)e^{\eta v(x)}.
\]
则
\[
  \mathbb E_{\mu_{\eta,v}}[f]
  =
  \frac{A_f(\eta)}{Z_v(\eta)}.
\]
对 \(\eta\) 求导得到
\[
  \frac{d}{d\eta}
  \mathbb E_{\mu_{\eta,v}}[f]
  =
  \frac{A_f'(\eta)Z_v(\eta)-A_f(\eta)Z_v'(\eta)}{Z_v(\eta)^2}.
\]
在 \(\eta=0\) 处，
\[
  Z_v(0)=1,
  \qquad
  A_f(0)=\mathbb E_\mu[f],
\]
\[
  A_f'(0)=\mathbb E_\mu[fv],
  \qquad
  Z_v'(0)=\mathbb E_\mu[v].
\]
因此
\[
  \left.
  \frac{d}{d\eta}
  \mathbb E_{\mu_{\eta,v}}[f]
  \right|_{\eta=0}
  =
  \mathbb E_\mu[fv]
  -
  \mathbb E_\mu[f]\mathbb E_\mu[v]
  =
  \operatorname{Cov}_\mu(f,v).
\]
又由定义
\[
  \operatorname{Cov}_\mu(f,v)=\langle f,v\rangle_\mu.
\]
证毕。
\end{proof}

\subsection{定理：正交扭曲的一阶命中加速条件}
\label{subsec:tex55-orthogonal-first-order}

\begin{theorem}[目标正交扭曲的一阶命中加速条件]
\label{thm:tex55-orthogonal-first-order-hit}
令
\[
  p(\eta)=\mu_{\eta,v}(S),
  \qquad
  L(\eta)=\log\frac{p(\eta)}{1-p(\eta)}.
\]
若
\[
  \langle v,u\rangle_\mu=0,
\]
则目标扭曲 \(u\) 的期望在 \(\eta=0\) 处一阶不变：
\[
  \left.
  \frac{d}{d\eta}
  \mathbb E_{\mu_{\eta,v}}[u]
  \right|_{\eta=0}
  =
  0.
\]
同时命中概率的一阶导数为
\[
  \left.
  \frac{d}{d\eta}
  p(\eta)
  \right|_{\eta=0}
  =
  p_\mu q_\mu
  \left(
  \mathbb E_\mu[v\mid S]
  -
  \mathbb E_\mu[v\mid S^c]
  \right),
\]
命中对数胜算的一阶导数为
\[
  L'(0)
  =
  \mathbb E_\mu[v\mid S]
  -
  \mathbb E_\mu[v\mid S^c].
\]
因此，若
\[
  \mathbb E_\mu[v\mid S]
  >
  \mathbb E_\mu[v\mid S^c],
\]
则沿 \(v\) 的同伦扭曲一阶提高命中概率与命中胜算。
\end{theorem}

\begin{Certificate}[谓词因果链证书]
\label{cert:tex55-orthogonal-first-order-hit}
本定理的谓词因果链为
\[
\begin{aligned}
  &\mu_{\eta,v}(x)=\mu(x)e^{\eta v(x)}/Z_v(\eta)\\
  &\Longrightarrow
  \frac{d}{d\eta}\mathbb E_{\mu_{\eta,v}}[f]\big|_{\eta=0}
  =
  \operatorname{Cov}_\mu(f,v)\\
  &\Longrightarrow
  f=u
  \Longrightarrow
  \text{目标期望一阶变化 }=\langle u,v\rangle_\mu\\
  &\Longrightarrow
  \langle u,v\rangle_\mu=0
  \Longrightarrow
  \text{目标扭曲一阶不变}\\
  &\Longrightarrow
  f=\mathbf 1_S
  \Longrightarrow
  p'(0)=\operatorname{Cov}_\mu(\mathbf 1_S,v)\\
  &\Longrightarrow
  p'(0)=p_\mu q_\mu
  \left(\mathbb E_\mu[v\mid S]-\mathbb E_\mu[v\mid S^c]\right)\\
  &\Longrightarrow
  L'(0)=
  \mathbb E_\mu[v\mid S]-\mathbb E_\mu[v\mid S^c].
\end{aligned}
\]
\end{Certificate}

\begin{proof}
由定理 \ref{thm:tex55-exponential-twist-derivative}，取 \(f=u\)，得到
\[
  \left.
  \frac{d}{d\eta}
  \mathbb E_{\mu_{\eta,v}}[u]
  \right|_{\eta=0}
  =
  \operatorname{Cov}_\mu(u,v)
  =
  \langle u,v\rangle_\mu.
\]
若 \(\langle u,v\rangle_\mu=0\)，则目标扭曲 \(u\) 的期望一阶不变。

再取 \(f=\mathbf 1_S\)，得到
\[
  p'(0)
  =
  \operatorname{Cov}_\mu(\mathbf 1_S,v).
\]
由 \(p_\mu=\mu(S)\)、\(q_\mu=\mu(S^c)\)，可写
\[
  \operatorname{Cov}_\mu(\mathbf 1_S,v)
  =
  p_\mu
  \left(
  \mathbb E_\mu[v\mid S]
  -
  \mathbb E_\mu[v]
  \right).
\]
而
\[
  \mathbb E_\mu[v]
  =
  p_\mu\mathbb E_\mu[v\mid S]
  +
  q_\mu\mathbb E_\mu[v\mid S^c].
\]
因此
\[
  \mathbb E_\mu[v\mid S]-\mathbb E_\mu[v]
  =
  q_\mu
  \left(
  \mathbb E_\mu[v\mid S]
  -
  \mathbb E_\mu[v\mid S^c]
  \right),
\]
从而
\[
  p'(0)
  =
  p_\mu q_\mu
  \left(
  \mathbb E_\mu[v\mid S]
  -
  \mathbb E_\mu[v\mid S^c]
  \right).
\]
又
\[
  L(\eta)=\log\frac{p(\eta)}{1-p(\eta)}
\]
给出
\[
  L'(\eta)
  =
  \frac{p'(\eta)}{p(\eta)(1-p(\eta))}.
\]
在 \(\eta=0\) 处代入上式，得到
\[
  L'(0)
  =
  \mathbb E_\mu[v\mid S]
  -
  \mathbb E_\mu[v\mid S^c].
\]
若该差值为正，则 \(p'(0)>0\) 且 \(L'(0)>0\)。
证毕。
\end{proof}

\subsection{引理与推论：最优正交命中扭曲方向}
\label{subsec:tex55-orthogonal-optimal-direction}

\begin{definition}[命中得分函数]
\label{def:tex55-hit-score-function}
在 \(0<p_\mu<1\) 时，定义命中得分函数
\[
  a_\mu(x)
  \DefEq
  \frac{\mathbf 1_S(x)}{p_\mu}
  -
  \frac{\mathbf 1_{S^c}(x)}{q_\mu}.
\]
则对任意方向 \(v\)，有
\[
  \langle v,a_\mu\rangle_\mu
  =
  \mathbb E_\mu[v\mid S]
  -
  \mathbb E_\mu[v\mid S^c].
\]
因此 \(a_\mu\) 是命中对数胜算 \(L_\mu\) 的得分方向。
\end{definition}

\begin{lemma}[目标正交条件下的最优命中扭曲方向]
\label{lem:tex55-optimal-orthogonal-hit-direction}
令目标扭曲子空间为
\[
  \mathcal U=\operatorname{span}\{u\},
\]
令
\[
  \mathcal U^\perp
  =
  \{v:\langle v,u\rangle_\mu=0\}.
\]
在约束
\[
  v\in\mathcal U^\perp,
  \qquad
  \lVert v\rVert_\mu=1
\]
下，使命中对数胜算一阶增长最大的方向为
\[
  v^\star
  =
  \frac{
  \Pi_{\mathcal U^\perp}a_\mu
  }{
  \lVert\Pi_{\mathcal U^\perp}a_\mu\rVert_\mu
  },
\]
前提是
\[
  \Pi_{\mathcal U^\perp}a_\mu\neq0.
\]
最大一阶增长率为
\[
  \lVert\Pi_{\mathcal U^\perp}a_\mu\rVert_\mu.
\]
\end{lemma}

\begin{Certificate}[投影最优化证书]
\label{cert:tex55-optimal-orthogonal-hit-direction}
目标谓词为
\[
  P(v):
  L'(0;v)=\langle v,a_\mu\rangle_\mu.
\]
在合法空间 \(v\in\mathcal U^\perp\) 上，
\[
  \langle v,a_\mu\rangle_\mu
  =
  \langle v,\Pi_{\mathcal U^\perp}a_\mu\rangle_\mu.
\]
于是 Cauchy--Schwarz 不等式给出最大值，等号方向即为归一化投影方向。
\end{Certificate}

\begin{proof}
由定义 \ref{def:tex55-hit-score-function} 与定理
\ref{thm:tex55-orthogonal-first-order-hit}，
\[
  \left.
  \frac{d}{d\eta}
  L_{\mu_{\eta,v}}
  \right|_{\eta=0}
  =
  \langle v,a_\mu\rangle_\mu.
\]
若 \(v\in\mathcal U^\perp\)，则
\[
  \langle v,a_\mu\rangle_\mu
  =
  \langle v,\Pi_{\mathcal U^\perp}a_\mu\rangle_\mu.
\]
由 Cauchy--Schwarz 不等式，
\[
  \langle v,\Pi_{\mathcal U^\perp}a_\mu\rangle_\mu
  \le
  \lVert v\rVert_\mu
  \lVert\Pi_{\mathcal U^\perp}a_\mu\rVert_\mu.
\]
在 \(\lVert v\rVert_\mu=1\) 下，右端为
\[
  \lVert\Pi_{\mathcal U^\perp}a_\mu\rVert_\mu.
\]
等号成立当且仅当
\[
  v
  =
  \frac{
  \Pi_{\mathcal U^\perp}a_\mu
  }{
  \lVert\Pi_{\mathcal U^\perp}a_\mu\rVert_\mu
  }.
\]
证毕。
\end{proof}

\begin{corollary}[正交扭曲能否加速的判定条件]
\label{cor:tex55-orthogonal-acceleration-criterion}
若
\[
  \Pi_{\mathcal U^\perp}a_\mu=0,
\]
则所有与目标扭曲 \(u\) 正交的方向都无法一阶提高命中对数胜算。
若
\[
  \Pi_{\mathcal U^\perp}a_\mu\neq0,
\]
则存在与目标扭曲正交的扭曲方向 \(v^\star\)，使
\[
  \left.
  \frac{d}{d\eta}
  L_{\mu_{\eta,v^\star}}
  \right|_{\eta=0}
  >
  0.
\]
因此“目标扭曲正交加速”应严格写成：
\[
  \boxed{
  \text{在目标扭曲正交补中，沿命中得分函数的投影方向扭曲。}
  }
\]
\end{corollary}

\begin{proof}
若 \(\Pi_{\mathcal U^\perp}a_\mu=0\)，则对所有 \(v\in\mathcal U^\perp\)，
\[
  \langle v,a_\mu\rangle_\mu
  =
  \langle v,\Pi_{\mathcal U^\perp}a_\mu\rangle_\mu
  =
  0.
\]
所以所有合法方向的一阶命中对数胜算变化均为零。

若 \(\Pi_{\mathcal U^\perp}a_\mu\neq0\)，取引理
\ref{lem:tex55-optimal-orthogonal-hit-direction} 中的 \(v^\star\)，则
\[
  \left.
  \frac{d}{d\eta}
  L_{\mu_{\eta,v^\star}}
  \right|_{\eta=0}
  =
  \lVert\Pi_{\mathcal U^\perp}a_\mu\rVert_\mu
  >
  0.
\]
证毕。
\end{proof}

\subsection{定理：命中泛函的幂化表示}
\label{subsec:tex55-power-form}

\begin{definition}[条件命中强度与条件未命中率]
\label{def:tex55-conditional-hit-power}
对同伦扭曲轨迹 \(\Gamma\)，定义首次命中时间
\[
  \tau_S^\Gamma
  =
  \inf\{t\ge1:Y_t^\Gamma=1\}.
\]
第 \(t\) 步条件命中强度定义为
\[
  h_t^\Gamma
  \DefEq
  \mathbb P_\Gamma
  \left(
  Y_t^\Gamma=1
  \mid
  \tau_S^\Gamma>t-1
  \right).
\]
对应条件未命中率为
\[
  r_t^\Gamma
  \DefEq
  1-h_t^\Gamma.
\]
\end{definition}

\begin{theorem}[命中泛函的幂化表示]
\label{thm:tex55-power-hit-functional}
有限时间命中泛函为
\[
  \mathcal J_T(\Gamma)
  =
  \mathbb P_\Gamma(\tau_S^\Gamma\le T).
\]
则
\[
  \boxed{
  \mathcal J_T(\Gamma)
  =
  1-
  \prod_{t=1}^{T}r_t^\Gamma.
  }
\]
取基准未命中率 \(q_0\in(0,1)\)，并定义幂化指数
\[
  A_T(\Gamma)
  \DefEq
  \sum_{t=1}^{T}
  \frac{\log r_t^\Gamma}{\log q_0},
\]
其中 \(0<r_t^\Gamma<1\)。
则
\[
  \prod_{t=1}^{T}r_t^\Gamma
  =
  q_0^{A_T(\Gamma)},
\]
因此
\[
  \boxed{
  \mathcal J_T(\Gamma)
  =
  1-q_0^{A_T(\Gamma)}.
  }
\]
\end{theorem}

\begin{Certificate}[条件概率链式证书]
\label{cert:tex55-power-hit-functional}
核心谓词链为
\[
\begin{aligned}
  &h_t^\Gamma
  =
  \mathbb P_\Gamma(Y_t^\Gamma=1\mid\tau_S^\Gamma>t-1)\\
  &\Longrightarrow
  r_t^\Gamma
  =
  \mathbb P_\Gamma(Y_t^\Gamma=0\mid\tau_S^\Gamma>t-1)\\
  &\Longrightarrow
  \mathbb P_\Gamma(\tau_S^\Gamma>T)
  =
  \prod_{t=1}^{T}r_t^\Gamma\\
  &\Longrightarrow
  \mathcal J_T(\Gamma)
  =
  1-\prod_{t=1}^{T}r_t^\Gamma\\
  &\Longrightarrow
  A_T(\Gamma)
  =
  \sum_{t=1}^{T}
  \frac{\log r_t^\Gamma}{\log q_0}
  \Longrightarrow
  \prod_{t=1}^{T}r_t^\Gamma=q_0^{A_T(\Gamma)}.
\end{aligned}
\]
\end{Certificate}

\begin{proof}
未命中到第 \(T\) 步的事件为 \(\{\tau_S^\Gamma>T\}\)。
由条件概率链式法则，
\[
  \mathbb P_\Gamma(\tau_S^\Gamma>T)
  =
  \prod_{t=1}^{T}
  \mathbb P_\Gamma
  \left(
  Y_t^\Gamma=0
  \mid
  \tau_S^\Gamma>t-1
  \right).
\]
由定义
\[
  \mathbb P_\Gamma
  \left(
  Y_t^\Gamma=0
  \mid
  \tau_S^\Gamma>t-1
  \right)
  =
  1-h_t^\Gamma
  =
  r_t^\Gamma.
\]
因此
\[
  \mathbb P_\Gamma(\tau_S^\Gamma>T)
  =
  \prod_{t=1}^{T}r_t^\Gamma.
\]
所以
\[
  \mathcal J_T(\Gamma)
  =
  \mathbb P_\Gamma(\tau_S^\Gamma\le T)
  =
  1-\mathbb P_\Gamma(\tau_S^\Gamma>T)
  =
  1-\prod_{t=1}^{T}r_t^\Gamma.
\]
又由 \(A_T(\Gamma)\) 的定义，
\[
  q_0^{A_T(\Gamma)}
  =
  q_0^{
  \sum_{t=1}^{T}
  \frac{\log r_t^\Gamma}{\log q_0}
  }
  =
  \prod_{t=1}^{T}
  q_0^{\frac{\log r_t^\Gamma}{\log q_0}}
  =
  \prod_{t=1}^{T}r_t^\Gamma.
\]
代回即可得到
\[
  \mathcal J_T(\Gamma)=1-q_0^{A_T(\Gamma)}.
\]
证毕。
\end{proof}

\begin{corollary}[幂化加速判据]
\label{cor:tex55-power-acceleration-criterion}
若基准算法每步未命中率为 \(q_0\)，则基准 \(T\) 步命中概率为
\[
  \mathcal J_T^{(0)}
  =
  1-q_0^T.
\]
若同伦扭曲轨迹 \(\Gamma\) 满足
\[
  A_T(\Gamma)>T,
\]
则
\[
  \mathcal J_T(\Gamma)
  >
  \mathcal J_T^{(0)}.
\]
因此 \(A_T(\Gamma)>T\) 是相对于基准算法的幂化加速命中条件。
\end{corollary}

\begin{proof}
由于 \(0<q_0<1\)，函数 \(a\mapsto q_0^a\) 严格递减。
若 \(A_T(\Gamma)>T\)，则
\[
  q_0^{A_T(\Gamma)}<q_0^T.
\]
由定理 \ref{thm:tex55-power-hit-functional}，
\[
  \mathcal J_T(\Gamma)
  =
  1-q_0^{A_T(\Gamma)}
  >
  1-q_0^T
  =
  \mathcal J_T^{(0)}.
\]
证毕。
\end{proof}

\subsection{定理：目标正交命中扭曲导致幂化指数提升}
\label{subsec:tex55-orthogonal-power-improvement}

\begin{theorem}[目标正交命中扭曲导致未命中率下降]
\label{thm:tex55-orthogonal-power-improvement}
设第 \(t\) 步当前分布为 \(\mu_t\)，当前命中得分函数为
\[
  a_{\mu_t}
  =
  \frac{\mathbf 1_S}{\mu_t(S)}
  -
  \frac{\mathbf 1_{S^c}}{\mu_t(S^c)}.
\]
令目标扭曲方向为 \(u_t\)，并设
\[
  \mathcal U_t=\operatorname{span}\{u_t\}.
\]
若
\[
  \Pi_{\mathcal U_t^\perp}a_{\mu_t}\neq0,
\]
取正交命中方向
\[
  v_t
  =
  \frac{
  \Pi_{\mathcal U_t^\perp}a_{\mu_t}
  }{
  \lVert\Pi_{\mathcal U_t^\perp}a_{\mu_t}\rVert_{\mu_t}
  }.
\]
则存在足够小的正步长 \(\eta_t>0\)，使同伦扭曲更新
\[
  \mu_{t+1}(x)
  =
  \frac{
  \mu_t(x)e^{\eta_t v_t(x)}
  }{
  Z_t
  }
\]
满足
\[
  L_{\mu_{t+1}}>L_{\mu_t},
  \qquad
  \mu_{t+1}(S)>\mu_t(S),
  \qquad
  1-\mu_{t+1}(S)<1-\mu_t(S).
\]
于是命中泛函中的生存乘积下降；以基准未命中率 \(q_0\) 衡量时，对应幂化指数相对于未扭曲过程上升。
\end{theorem}

\begin{Certificate}[谓词因果链证书]
\label{cert:tex55-orthogonal-power-improvement}
证明链为
\[
\begin{aligned}
  &\Pi_{\mathcal U_t^\perp}a_{\mu_t}\neq0\\
  &\Longrightarrow
  v_t=
  \Pi_{\mathcal U_t^\perp}a_{\mu_t}/
  \lVert\Pi_{\mathcal U_t^\perp}a_{\mu_t}\rVert_{\mu_t}\\
  &\Longrightarrow
  \left.
  \frac{d}{d\eta}L_{\mu_{\eta,v_t}}
  \right|_{\eta=0}
  =
  \lVert\Pi_{\mathcal U_t^\perp}a_{\mu_t}\rVert_{\mu_t}>0\\
  &\Longrightarrow
  \exists\eta_t>0,\quad L_{\mu_{t+1}}>L_{\mu_t}\\
  &\Longrightarrow
  \mu_{t+1}(S)>\mu_t(S)\\
  &\Longrightarrow
  r_{t+1}<r_t\\
  &\Longrightarrow
  \prod r_t\text{ 下降，}A_T\text{ 上升。}
\end{aligned}
\]
\end{Certificate}

\begin{proof}
由引理 \ref{lem:tex55-optimal-orthogonal-hit-direction}，方向
\[
  v_t
  =
  \frac{
  \Pi_{\mathcal U_t^\perp}a_{\mu_t}
  }{
  \lVert\Pi_{\mathcal U_t^\perp}a_{\mu_t}\rVert_{\mu_t}
  }
\]
是在目标正交约束下使命中对数胜算一阶增长最大的方向，且
\[
  \left.
  \frac{d}{d\eta}
  L_{\mu_{\eta,v_t}}
  \right|_{\eta=0}
  =
  \lVert\Pi_{\mathcal U_t^\perp}a_{\mu_t}\rVert_{\mu_t}
  >
  0.
\]
由于 \(L_{\mu_{\eta,v_t}}\) 关于 \(\eta\) 连续可微，所以存在足够小的
\(\eta_t>0\)，使
\[
  L_{\mu_{t+1}}>L_{\mu_t}.
\]
又因为
\[
  L_\mu=\log\frac{\mu(S)}{1-\mu(S)}
\]
是 \(\mu(S)\in(0,1)\) 的严格单调函数，故
\[
  L_{\mu_{t+1}}>L_{\mu_t}
  \Longrightarrow
  \mu_{t+1}(S)>\mu_t(S).
\]
于是
\[
  1-\mu_{t+1}(S)<1-\mu_t(S).
\]
这说明未命中率下降，因而有限时间生存乘积下降。
由定理 \ref{thm:tex55-power-hit-functional} 的表示
\[
  \prod_{t=1}^{T}r_t^\Gamma=q_0^{A_T(\Gamma)},
\]
在 \(0<q_0<1\) 下，生存乘积下降对应幂化指数上升。
证毕。
\end{proof}

\subsection{依次正交泛迭代与高阶幂化加速}
\label{subsec:tex55-sequential-orthogonal-iteration}

\begin{definition}[依次正交扭曲子空间]
\label{def:tex55-sequential-orthogonal-space}
第 \(k\) 步之前已经选择 \(v_1,\ldots,v_{k-1}\)。
定义已用方向子空间
\[
  \mathcal W_{k-1}
  =
  \operatorname{span}\{u,v_1,\ldots,v_{k-1}\}.
\]
第 \(k\) 步的合法扭曲方向满足
\[
  v_k\in\mathcal W_{k-1}^\perp.
\]
即在当前分布 \(\mu_k\) 下，
\[
  \langle v_k,u\rangle_{\mu_k}=0,
  \qquad
  \langle v_k,v_i\rangle_{\mu_k}=0,
  \quad i=1,\ldots,k-1.
\]
若
\[
  \Pi_{\mathcal W_{k-1}^\perp}a_{\mu_k}\neq0,
\]
则取
\[
  v_k
  =
  \frac{
  \Pi_{\mathcal W_{k-1}^\perp}a_{\mu_k}
  }{
  \lVert\Pi_{\mathcal W_{k-1}^\perp}a_{\mu_k}\rVert_{\mu_k}
  }.
\]
这就是命中得分函数在依次正交约束下的残差方向。
\end{definition}

\begin{theorem}[依次正交泛迭代的单步增益]
\label{thm:tex55-sequential-orthogonal-gain}
若每一步满足
\[
  \Pi_{\mathcal W_{k-1}^\perp}a_{\mu_k}\neq0,
\]
并沿
\[
  v_k
  =
  \frac{
  \Pi_{\mathcal W_{k-1}^\perp}a_{\mu_k}
  }{
  \lVert\Pi_{\mathcal W_{k-1}^\perp}a_{\mu_k}\rVert_{\mu_k}
  }
\]
做足够小的正向同伦扭曲，则
\[
  L_{\mu_{k+1}}>L_{\mu_k},
  \qquad
  \mu_{k+1}(S)>\mu_k(S).
\]
因此依次正交泛迭代在每一步都给出一阶命中增益。
\end{theorem}

\begin{proof}
将引理 \ref{lem:tex55-optimal-orthogonal-hit-direction} 中的
\(\mathcal U^\perp\) 替换为 \(\mathcal W_{k-1}^\perp\)。
若投影
\[
  \Pi_{\mathcal W_{k-1}^\perp}a_{\mu_k}
\]
非零，则归一化投影方向使 \(L_\mu\) 的一阶导数为
\[
  \lVert\Pi_{\mathcal W_{k-1}^\perp}a_{\mu_k}\rVert_{\mu_k}>0.
\]
由连续可微性，足够小的正步长给出
\[
  L_{\mu_{k+1}}>L_{\mu_k}.
\]
再由 \(L_\mu\) 对 \(\mu(S)\) 的严格单调性得到
\[
  \mu_{k+1}(S)>\mu_k(S).
\]
证毕。
\end{proof}

\begin{theorem}[每层幂压缩推出高阶幂化加速]
\label{thm:tex55-higher-order-power-acceleration}
设 \(Q_k\) 表示第 \(k\) 层迭代后的生存概率。
若存在 \(A_k>1\)，使每层满足
\[
  Q_k\le Q_{k-1}^{A_k},
  \qquad
  k=1,\ldots,m,
\]
则
\[
  \boxed{
  Q_m
  \le
  Q_0^{A_1A_2\cdots A_m}.
  }
\]
因此命中泛函满足
\[
  \boxed{
  \mathcal J_m
  =
  1-Q_m
  \ge
  1-Q_0^{A_1A_2\cdots A_m}.
  }
\]
这就是高阶幂化加速命中。
\end{theorem}

\begin{Certificate}[数学归纳法证书]
\label{cert:tex55-higher-order-power-acceleration}
目标谓词为
\[
  P(m):
  Q_m\le Q_0^{\prod_{k=1}^{m}A_k}.
\]
基例 \(m=1\) 由假设
\[
  Q_1\le Q_0^{A_1}
\]
直接成立。

归纳步：假设 \(P(m)\) 成立，即
\[
  Q_m\le Q_0^{\prod_{k=1}^{m}A_k}.
\]
又由第 \(m+1\) 层幂压缩假设
\[
  Q_{m+1}\le Q_m^{A_{m+1}},
\]
代入归纳假设得
\[
  Q_{m+1}
  \le
  \left(
  Q_0^{\prod_{k=1}^{m}A_k}
  \right)^{A_{m+1}}
  =
  Q_0^{\prod_{k=1}^{m+1}A_k}.
\]
故 \(P(m+1)\) 成立。
由数学归纳法，
\[
  \forall m\ge1,\quad
  Q_m\le Q_0^{\prod_{k=1}^{m}A_k}.
\]
\end{Certificate}

\begin{proof}
证书 \ref{cert:tex55-higher-order-power-acceleration} 已给出完整归纳链。
由
\[
  Q_m\le Q_0^{A_1A_2\cdots A_m}
\]
立刻得到
\[
  \mathcal J_m=1-Q_m
  \ge
  1-Q_0^{A_1A_2\cdots A_m}.
\]
证毕。
\end{proof}

\begin{corollary}[高阶幂化的条件口径]
\label{cor:tex55-higher-order-power-boundary}
依次正交泛迭代只保证在非零投影与足够小正步长下得到一阶命中增益。
要推出
\[
  Q_m
  \le
  Q_0^{A_1A_2\cdots A_m},
\]
还必须额外满足每层生存概率幂压缩条件
\[
  Q_k\le Q_{k-1}^{A_k},
  \qquad
  A_k>1.
\]
因此严格命题是
\[
  \boxed{
  \begin{gathered}
  \text{依次正交泛迭代}
  +
  \text{每层生存概率幂压缩}\\
  \Longrightarrow
  \text{高阶幂化加速命中。}
  \end{gathered}
  }
\]
\end{corollary}

\begin{proof}
由定理 \ref{thm:tex55-sequential-orthogonal-gain}，依次正交泛迭代的非零投影只给出
\[
  \mu_{k+1}(S)>\mu_k(S),
\]
从而给出未命中率单调下降。
但单调下降一般只能推出 \(Q_k<Q_{k-1}\)，不能推出存在 \(A_k>1\) 使
\[
  Q_k\le Q_{k-1}^{A_k}.
\]
所以高阶幂化必须把每层幂压缩作为额外假设。
在该假设成立时，定理 \ref{thm:tex55-higher-order-power-acceleration} 给出高阶幂化结论。
证毕。
\end{proof}

\subsection{谓词因果链、边界与最终口径}
\label{subsec:tex55-orthogonal-chain-boundary}

\begin{remark}[核心谓词因果链]
本章的完整因果链为
\[
\begin{aligned}
  &\text{命中集合 }S\\
  &\Longrightarrow
  \text{命中指示变量 }Y^\lambda=\mathbf 1_S(X^\lambda)\\
  &\Longrightarrow
  \text{命中泛函 }\mathcal J(\lambda)=\mathbb E[Y^\lambda]\\
  &\Longrightarrow
  \text{命中得分函数 }
  a_\mu=
  \frac{\mathbf 1_S}{\mu(S)}
  -
  \frac{\mathbf 1_{S^c}}{\mu(S^c)}\\
  &\Longrightarrow
  \text{目标正交合法空间 }\mathcal U^\perp\\
  &\Longrightarrow
  \text{最优正交扭曲方向 }
  v^\star=
  \frac{\Pi_{\mathcal U^\perp}a_\mu}
  {\lVert\Pi_{\mathcal U^\perp}a_\mu\rVert_\mu}\\
  &\Longrightarrow
  \text{命中对数胜算提升 }L_{\mu_+}>L_\mu\\
  &\Longrightarrow
  \text{未命中率下降 }Q_+<Q\\
  &\Longrightarrow
  \text{生存概率幂化表示 }Q_T=q_0^{A_T}\\
  &\Longrightarrow
  \text{幂化加速命中 }\mathcal J_T=1-q_0^{A_T}.
\end{aligned}
\]
若继续做依次正交残差迭代，则
\[
  \mathcal W_{k-1}
  =
  \operatorname{span}\{u,v_1,\ldots,v_{k-1}\},
\]
\[
  v_k
  =
  \frac{\Pi_{\mathcal W_{k-1}^\perp}a_{\mu_k}}
  {\lVert\Pi_{\mathcal W_{k-1}^\perp}a_{\mu_k}\rVert_{\mu_k}}.
\]
在每层幂压缩条件下得到
\[
  Q_m
  \le
  Q_0^{\prod_{k=1}^{m}A_k}.
\]
\end{remark}

\begin{remark}[必要边界]
第一，正交于目标不等于必然加速。
若
\[
  \langle v,u\rangle_\mu=0
  \quad\text{但}\quad
  \langle v,a_\mu\rangle_\mu=0,
\]
则
\[
  \left.
  \frac{d}{d\eta}
  L_{\mu_{\eta,v}}
  \right|_{\eta=0}
  =
  0.
\]
若 \(\langle v,a_\mu\rangle_\mu<0\)，则该正交扭曲反而降低命中概率。

第二，正交性应随分布更新而重算。
每次扭曲后分布由 \(\mu_k\) 变为 \(\mu_{k+1}\)，内积也从
\[
  \langle\cdot,\cdot\rangle_{\mu_k}
\]
变为
\[
  \langle\cdot,\cdot\rangle_{\mu_{k+1}}.
\]
因此每一步应重新计算投影方向。

第三，高阶幂化需要层级压缩条件。
从 \(Q_k<Q_{k-1}\) 只能推出单调改善；
要推出 \(Q_k\le Q_{k-1}^{A_k}\) 且 \(A_k>1\)，必须补充更强假设。
\end{remark}

\begin{remark}[最终严格表述]
用户命题可严格写成：
\[
  \boxed{
  \begin{gathered}
  \text{命中泛函参数化后，其梯度给出命中得分函数；}\\
  \text{将该得分函数投影到目标扭曲的正交补空间，}\\
  \text{并沿该投影方向做同伦扭曲，可在不一阶扰动目标扭曲的条件下提高命中胜算。}
  \end{gathered}
  }
\]
进一步，
\[
  \boxed{
  \begin{gathered}
  \text{由该扭曲诱导的条件命中强度序列 }h_t
  \text{ 使生存概率}\\
  \prod_{t=1}^{T}(1-h_t)
  \text{ 化为 }q_0^{A_T},
  \text{ 从而形成幂化加速命中。}
  \end{gathered}
  }
\]
再进一步，
\[
  \boxed{
  \begin{gathered}
  \text{若每一步在目标扭曲与前序扭曲方向的正交补中}\\
  \text{继续选择命中得分残差方向，则形成依次正交泛迭代；}\\
  \text{若每层满足生存概率幂压缩，则得到高阶幂化加速命中。}
  \end{gathered}
  }
\]
最压缩的数学表达是
\[
  \boxed{
  v_k
  =
  \frac{
  \Pi_{\operatorname{span}\{u,v_1,\ldots,v_{k-1}\}^{\perp}}
  a_{\mu_k}
  }{
  \left\lVert
  \Pi_{\operatorname{span}\{u,v_1,\ldots,v_{k-1}\}^{\perp}}
  a_{\mu_k}
  \right\rVert_{\mu_k}
  },
  }
\]
其中
\[
  a_{\mu_k}
  =
  \frac{\mathbf 1_S}{\mu_k(S)}
  -
  \frac{\mathbf 1_{S^c}}{\mu_k(S^c)}.
\]
扭曲更新为
\[
  \mu_{k+1}(x)
  =
  \frac{\mu_k(x)e^{\eta_k v_k(x)}}{Z_k}.
\]
命中加速为
\[
  \boxed{
  \mathcal J_T
  =
  1-q_0^{A_T}.
  }
\]
高阶幂化形式为
\[
  \boxed{
  \mathcal J_m
  \ge
  1-Q_0^{A_1A_2\cdots A_m}.
  }
\]
这就是“将命中泛函参数化，并与目标扭曲正交进行扭曲，实现幂化加速命中；
依次正交泛迭代形成高阶幂化加速命中”的严格形式。
\end{remark}

\clearpage

%%%%%%%%%%%%%%%%%%%%%%%%%%%%%%%%%%%%%%%%%%%%%%%%%%%%%%%%%%%%
\section{形式逻辑：高阶正交雅可比间隙下降与幂化命中加速}
\label{sec:tex55-jacobian-gap-power-hitting}
%%%%%%%%%%%%%%%%%%%%%%%%%%%%%%%%%%%%%%%%%%%%%%%%%%%%%%%%%%%%

\begin{remark}[核心陈述]
本章把“高阶正交命中加速”表述为带收敛指引的版本：
\[
  \boxed{
  \begin{gathered}
  \text{高阶正交命中加速}
  +
  \text{每阶雅可比间隙梯度下降}\\
  \Longrightarrow
  \text{带收敛指引的高阶幂化命中加速。}
  \end{gathered}
  }
\]
这里的雅可比间隙不是普通目标函数梯度。
它衡量的是：
\[
  \boxed{
  \begin{gathered}
  \text{当前同伦扭曲实际能产生的轨迹变化}\\
  \text{与理想高阶正交命中变化之间的差。}
  \end{gathered}
  }
\]
因此每一阶先把命中梯度投影到高阶正交补空间，
再让同伦扭曲族的雅可比方向尽量逼近这个投影方向。
梯度下降只作用在“雅可比间隙”上，用作每阶同伦扭曲的收敛指引。
标准光滑下降引理可参考 \cite{NocedalWright2006NumericalOptimization}；
本章仍在有限维符号框架中给出完整推演。
\end{remark}

\subsection{对象域：命中随机过程、命中指数与轨迹梯度}
\label{subsec:tex55-jacgap-objects}

\begin{definition}[命中随机过程与命中泛函]
\label{def:tex55-jacgap-hit-process}
令 \(\Gamma\) 为轨迹函数，并令
\[
  X_t^\Gamma:\Xi\to\Omega
\]
为由 \(\Gamma\) 诱导的第 \(t\) 个状态随机变量。
令命中集合为
\[
  S\subseteq\Omega.
\]
定义命中指示变量
\[
  Y_t^\Gamma=\mathbf 1_S(X_t^\Gamma).
\]
首次命中时间为
\[
  \tau_S^\Gamma
  =
  \inf\{t\ge1:Y_t^\Gamma=1\}.
\]
有限时间命中泛函定义为
\[
  \mathcal J_T(\Gamma)
  =
  \mathbb P_\Gamma(\tau_S^\Gamma\le T).
\]
等价地，令生存概率为
\[
  Q_T(\Gamma)
  =
  \mathbb P_\Gamma(\tau_S^\Gamma>T),
\]
则
\[
  \mathcal J_T(\Gamma)=1-Q_T(\Gamma).
\]
\end{definition}

\begin{definition}[命中指数泛函]
\label{def:tex55-jacgap-hit-index}
在 \(0<Q_T(\Gamma)<1\) 时，定义命中指数泛函为
\[
  \mathcal C_T(\Gamma)
  =
  -\log Q_T(\Gamma).
\]
于是
\[
  Q_T(\Gamma)=e^{-\mathcal C_T(\Gamma)},
  \qquad
  \mathcal J_T(\Gamma)=1-e^{-\mathcal C_T(\Gamma)}.
\]
由于函数 \(c\mapsto 1-e^{-c}\) 严格递增，所以
\[
  \boxed{
  \max \mathcal J_T(\Gamma)
  \Longleftrightarrow
  \max \mathcal C_T(\Gamma).
  }
\]
幂化加速
\[
  Q_T(\Gamma_1)\le Q_T(\Gamma_0)^A,
  \qquad A>1,
\]
等价于
\[
  \mathcal C_T(\Gamma_1)\ge A\mathcal C_T(\Gamma_0).
\]
因此 \(\mathcal C_T\) 是研究生存概率幂化下降的自然坐标。
\end{definition}

\begin{definition}[轨迹函数空间与命中梯度]
\label{def:tex55-jacgap-trajectory-gradient}
令所有可行轨迹函数组成有限维光滑空间
\[
  \mathfrak T.
\]
在切空间 \(T_\Gamma\mathfrak T\) 上设内积
\[
  \langle\cdot,\cdot\rangle_\Gamma.
\]
若 \(\mathcal C_T\) 在 \(\Gamma\) 处可微，则存在命中梯度
\[
  g_\Gamma
  =
  \nabla_\Gamma\mathcal C_T(\Gamma)
  \in T_\Gamma\mathfrak T
\]
满足对任意轨迹扰动方向
\[
  \delta\in T_\Gamma\mathfrak T
\]
有
\[
  D\mathcal C_T(\Gamma)[\delta]
  =
  \langle g_\Gamma,\delta\rangle_\Gamma.
\]
这里 \(\delta\) 表示轨迹函数的一阶变化方向，而不是状态点本身。
\end{definition}

\subsection{高阶正交命中结构}
\label{subsec:tex55-jacgap-orthogonal-structure}

\begin{definition}[零阶目标扭曲方向]
\label{def:tex55-jacgap-zero-target}
令
\[
  u_0\in T_\Gamma\mathfrak T
\]
表示原始目标扭曲方向。
它可以来自能量方向 \(-E\)、频域目标方向、能量梯度诱导方向，或既有目标分布的扭曲方向。
定义第一层已用方向空间为
\[
  \mathcal W_0
  =
  \operatorname{span}\{u_0\}.
\]
\end{definition}

\begin{definition}[一阶正交命中加速方向]
\label{def:tex55-jacgap-first-order-direction}
给定初始轨迹
\[
  \Gamma^{[0]}=\Gamma,
\]
定义一阶正交命中残差为
\[
  b_1
  =
  \Pi_{\mathcal W_0^\perp}
  g_{\Gamma^{[0]}}.
\]
若 \(b_1\neq0\)，则一阶归一化理想命中方向为
\[
  \widehat b_1
  =
  \frac{b_1}{\lVert b_1\rVert}.
\]
它满足
\[
  \langle \widehat b_1,u_0\rangle=0,
\]
并且
\[
  D\mathcal C_T(\Gamma^{[0]})[\widehat b_1]
  =
  \langle g_{\Gamma^{[0]}},\widehat b_1\rangle
  =
  \lVert b_1\rVert
  >
  0.
\]
因此一阶正交命中方向是在不一阶扰动目标扭曲方向 \(u_0\) 的条件下，
提高命中指数的方向。
\end{definition}

\begin{definition}[高阶正交命中加速方向]
\label{def:tex55-jacgap-higher-order-direction}
假设已经得到前 \(k-1\) 阶实际扭曲方向
\[
  \delta_1,\ldots,\delta_{k-1}.
\]
定义已用方向空间
\[
  \mathcal W_{k-1}
  =
  \operatorname{span}\{u_0,\delta_1,\ldots,\delta_{k-1}\}.
\]
第 \(k\) 阶命中梯度为
\[
  g_k
  =
  g_{\Gamma^{[k-1]}}
  =
  \nabla\mathcal C_T(\Gamma^{[k-1]}).
\]
第 \(k\) 阶理想正交命中方向定义为
\[
  b_k
  =
  \Pi_{\mathcal W_{k-1}^{\perp}}g_k.
\]
若 \(b_k\neq0\)，则归一化理想方向为
\[
  \widehat b_k
  =
  \frac{b_k}{\lVert b_k\rVert}.
\]
这表示第 \(k\) 阶不重复前 \(k-1\) 阶已经使用过的命中加速方向，
而是在剩余正交补空间中继续寻找命中梯度残差。
\end{definition}

\begin{remark}[高阶正交的递归含义]
由定义 \ref{def:tex55-jacgap-first-order-direction} 与
\ref{def:tex55-jacgap-higher-order-direction}，可概括为
\[
  \boxed{
  1\text{阶}
  =
  \text{直接目标正交命中加速};
  }
\]
\[
  \boxed{
  n\text{阶}
  =
  \text{对前 }n-1\text{ 阶命中加速后的剩余命中梯度继续正交命中加速。}
  }
\]
这就是“针对命中泛函随机过程轨迹函数进行正交命中”的递归结构。
\end{remark}

\subsection{每阶雅可比间隙}
\label{subsec:tex55-jacgap-definition}

\begin{definition}[每阶同伦扭曲族与雅可比]
\label{def:tex55-jacgap-homotopy-family}
第 \(k\) 阶同伦扭曲由参数
\[
  \theta_k\in\Theta_k
\]
控制，记为
\[
  \mathcal H_k(\theta_k):\mathfrak T\to\mathfrak T.
\]
从当前轨迹 \(\Gamma^{[k-1]}\) 出发，得到扭曲后的轨迹
\[
  \Gamma^{[k]}
  =
  \mathcal H_k(\theta_k)(\Gamma^{[k-1]}).
\]
若考虑小参数速度
\[
  d_k\in T_{\theta_k}\Theta_k,
\]
则第 \(k\) 阶雅可比方向定义为
\[
  J_k(\theta_k)d_k
  =
  D_{\theta_k}
  \left[
  \mathcal H_k(\theta_k)(\Gamma^{[k-1]})
  \right][d_k].
\]
它表示第 \(k\) 阶参数扰动 \(d_k\) 实际在轨迹函数空间中造成的变化方向。
\end{definition}

\begin{definition}[第 \(k\) 阶雅可比间隙]
\label{def:tex55-jacgap-gap}
第 \(k\) 阶理想方向为
\[
  b_k
  =
  \Pi_{\mathcal W_{k-1}^{\perp}}g_k.
\]
实际可实现方向为
\[
  \delta_k
  =
  J_k(\theta_k)d_k.
\]
给定惩罚系数 \(\rho_k>0\)，定义第 \(k\) 阶雅可比间隙为
\[
  \boxed{
  \mathfrak G_k(\theta_k,d_k)
  =
  \frac12
  \left\lVert
  \Pi_{\mathcal W_{k-1}^{\perp}}
  J_k(\theta_k)d_k
  -
  b_k
  \right\rVert^2
  +
  \frac{\rho_k}{2}
  \left\lVert
  \Pi_{\mathcal W_{k-1}}
  J_k(\theta_k)d_k
  \right\rVert^2.
  }
\]
第一项表示实际雅可比方向与理想高阶正交命中方向之间的逼近误差。
第二项表示实际方向泄漏到已用方向空间中的非正交误差。
因此
\[
  \boxed{
  \mathfrak G_k
  =
  \text{命中方向逼近误差}
  +
  \text{高阶正交泄漏误差}.
  }
\]
\end{definition}

\begin{definition}[每阶同伦扭曲的收敛指引]
\label{def:tex55-jacgap-gradient-step}
第 \(k\) 阶不直接求解全局问题
\[
  \max_\Gamma \mathcal C_T(\Gamma),
\]
而是先求
\[
  (\theta_k^\star,d_k^\star)
  \approx
  \arg\min_{\theta_k,d_k}
  \mathfrak G_k(\theta_k,d_k).
\]
令
\[
  z_k=(\theta_k,d_k).
\]
一种梯度下降形式为
\[
  z_{k,s+1}
  =
  \Pi_{\mathcal K_k}
  \left[
  z_{k,s}
  -
  \alpha_{k,s}
  \nabla \mathfrak G_k(z_{k,s})
  \right],
\]
其中 \(\mathcal K_k\) 是第 \(k\) 阶约束集，例如单位速度、温度范围、谱模态范围或稳定性约束。
这说明每一阶同伦扭曲的收敛方向由雅可比间隙下降指引。
\end{definition}

\subsection{定理：雅可比间隙小推出命中指数增长}
\label{subsec:tex55-jacgap-growth-theorem}

\begin{theorem}[第 \(k\) 阶雅可比逼近控制命中指数增长]
\label{thm:tex55-jacgap-growth}
设
\[
  g_k
  =
  \nabla\mathcal C_T(\Gamma^{[k-1]}),
  \qquad
  b_k
  =
  \Pi_{\mathcal W_{k-1}^{\perp}}g_k
  \neq0.
\]
令实际同伦扭曲方向为
\[
  \delta_k=J_k(\theta_k)d_k.
\]
若
\[
  \lVert\delta_k-b_k\rVert
  <
  \frac{\lVert b_k\rVert^2}{\lVert g_k\rVert},
\]
则
\[
  D\mathcal C_T(\Gamma^{[k-1]})[\delta_k]>0.
\]
因此取足够小的正步长 \(\eta_k>0\)，定义
\[
  \Gamma^{[k]}
  =
  \operatorname{Hom}_{\eta_k,\delta_k}
  (\Gamma^{[k-1]}),
\]
则
\[
  \mathcal C_T(\Gamma^{[k]})
  >
  \mathcal C_T(\Gamma^{[k-1]}).
\]
于是
\[
  Q_T(\Gamma^{[k]})
  <
  Q_T(\Gamma^{[k-1]}),
\]
即第 \(k\) 阶命中加速成立。
\end{theorem}

\begin{Certificate}[谓词因果链证书]
\label{cert:tex55-jacgap-growth}
本定理的谓词链为
\[
\begin{aligned}
  &b_k=\Pi_{\mathcal W_{k-1}^{\perp}}g_k\neq0\\
  &\Longrightarrow
  \langle g_k,b_k\rangle=\lVert b_k\rVert^2\\
  &\Longrightarrow
  \delta_k=b_k+e_k,\quad e_k=\delta_k-b_k\\
  &\Longrightarrow
  D\mathcal C_T[\delta_k]
  =
  \langle g_k,b_k\rangle+\langle g_k,e_k\rangle\\
  &\Longrightarrow
  D\mathcal C_T[\delta_k]
  \ge
  \lVert b_k\rVert^2-\lVert g_k\rVert\lVert e_k\rVert\\
  &\Longrightarrow
  \lVert e_k\rVert<\lVert b_k\rVert^2/\lVert g_k\rVert\\
  \Longrightarrow
  D\mathcal C_T[\delta_k]>0\\
  &\Longrightarrow
  \mathcal C_T(\Gamma^{[k]})
  >
  \mathcal C_T(\Gamma^{[k-1]})\\
  &\Longrightarrow
  Q_T(\Gamma^{[k]})
  <
  Q_T(\Gamma^{[k-1]}).
\end{aligned}
\]
\end{Certificate}

\begin{proof}
因为
\[
  b_k
  =
  \Pi_{\mathcal W_{k-1}^{\perp}}g_k,
\]
所以
\[
  g_k
  =
  \Pi_{\mathcal W_{k-1}}g_k+b_k,
  \qquad
  b_k\in\mathcal W_{k-1}^{\perp}.
\]
于是
\[
  \langle g_k,b_k\rangle
  =
  \langle b_k,b_k\rangle
  =
  \lVert b_k\rVert^2.
\]
令
\[
  e_k=\delta_k-b_k.
\]
则
\[
  \delta_k=b_k+e_k.
\]
命中指数的一阶变化为
\[
  D\mathcal C_T(\Gamma^{[k-1]})[\delta_k]
  =
  \langle g_k,\delta_k\rangle.
\]
代入 \(\delta_k=b_k+e_k\)，得到
\[
  \langle g_k,\delta_k\rangle
  =
  \langle g_k,b_k\rangle+\langle g_k,e_k\rangle.
\]
由上式与 Cauchy--Schwarz 不等式，
\[
  \langle g_k,\delta_k\rangle
  \ge
  \lVert b_k\rVert^2-\lVert g_k\rVert\lVert e_k\rVert.
\]
若
\[
  \lVert e_k\rVert
  <
  \frac{\lVert b_k\rVert^2}{\lVert g_k\rVert},
\]
则
\[
  \langle g_k,\delta_k\rangle>0.
\]
所以
\[
  D\mathcal C_T(\Gamma^{[k-1]})[\delta_k]>0.
\]
由于 \(\mathcal C_T\) 可微，沿正方向取足够小的步长 \(\eta_k>0\)，可得
\[
  \mathcal C_T(\Gamma^{[k]})
  >
  \mathcal C_T(\Gamma^{[k-1]}).
\]
再由
\[
  Q_T(\Gamma)=e^{-\mathcal C_T(\Gamma)}
\]
得到
\[
  Q_T(\Gamma^{[k]})
  <
  Q_T(\Gamma^{[k-1]}).
\]
证毕。
\end{proof}

\begin{corollary}[雅可比间隙小推出雅可比逼近条件]
\label{cor:tex55-jacgap-small-gap-implies-growth}
令
\[
  M_k=\max\{1,\rho_k^{-1}\}.
\]
若
\[
  \mathfrak G_k(\theta_k,d_k)
  <
  \frac{\lVert b_k\rVert^4}{2M_k\lVert g_k\rVert^2},
\]
则
\[
  \lVert\delta_k-b_k\rVert
  <
  \frac{\lVert b_k\rVert^2}{\lVert g_k\rVert},
\]
从而由定理 \ref{thm:tex55-jacgap-growth} 得到第 \(k\) 阶命中指数增长。
\end{corollary}

\begin{proof}
记
\[
  a=
  \Pi_{\mathcal W_{k-1}^{\perp}}\delta_k-b_k,
  \qquad
  c=
  \Pi_{\mathcal W_{k-1}}\delta_k.
\]
因为 \(b_k\in\mathcal W_{k-1}^\perp\)，且
\[
  \delta_k-b_k=a+c
\]
是正交分解，所以
\[
  \lVert\delta_k-b_k\rVert^2
  =
  \lVert a\rVert^2+\lVert c\rVert^2.
\]
由雅可比间隙定义，
\[
  \mathfrak G_k
  =
  \frac12\lVert a\rVert^2
  +
  \frac{\rho_k}{2}\lVert c\rVert^2.
\]
因此
\[
  \lVert a\rVert^2+\lVert c\rVert^2
  \le
  2\max\{1,\rho_k^{-1}\}\mathfrak G_k
  =
  2M_k\mathfrak G_k.
\]
若
\[
  \mathfrak G_k
  <
  \frac{\lVert b_k\rVert^4}{2M_k\lVert g_k\rVert^2},
\]
则
\[
  \lVert\delta_k-b_k\rVert^2
  <
  \frac{\lVert b_k\rVert^4}{\lVert g_k\rVert^2}.
\]
取平方根即得结论。
证毕。
\end{proof}

\subsection{公理降级：每阶雅可比间隙梯度下降}
\label{subsec:tex55-jacgap-gradient-descent}

\begin{theorem}[公理降级：光滑雅可比间隙的下降引理]
\label{thm:tex55-jacgap-smooth-descent}
设第 \(k\) 阶雅可比间隙 \(\mathfrak G_k\) 是 \(L_k\)-光滑函数，即
\[
  \lVert\nabla\mathfrak G_k(z)-\nabla\mathfrak G_k(z')\rVert
  \le
  L_k\lVert z-z'\rVert.
\]
若采用无约束梯度下降
\[
  z_{s+1}=z_s-\alpha_k\nabla\mathfrak G_k(z_s),
\]
并且
\[
  0<\alpha_k\le\frac1{L_k},
\]
则
\[
  \mathfrak G_k(z_{s+1})
  \le
  \mathfrak G_k(z_s)
  -
  \frac{\alpha_k}{2}
  \lVert\nabla\mathfrak G_k(z_s)\rVert^2.
\]
因此 \(\mathfrak G_k(z_s)\) 单调不增。
\end{theorem}

\begin{Certificate}[符号推演证书]
\label{cert:tex55-jacgap-smooth-descent}
目标谓词为
\[
  P(s):
  \mathfrak G_k(z_{s+1})
  \le
  \mathfrak G_k(z_s)
  -
  \frac{\alpha_k}{2}
  \lVert\nabla\mathfrak G_k(z_s)\rVert^2.
\]
证书使用 \(L_k\)-光滑性的下降引理：
\[
  \mathfrak G_k(z')
  \le
  \mathfrak G_k(z)
  +
  \langle\nabla\mathfrak G_k(z),z'-z\rangle
  +
  \frac{L_k}{2}\lVert z'-z\rVert^2.
\]
\end{Certificate}

\begin{proof}
由 \(L_k\)-光滑性，对
\[
  z'=z_{s+1},
  \qquad
  z=z_s
\]
有
\[
  \mathfrak G_k(z_{s+1})
  \le
  \mathfrak G_k(z_s)
  +
  \langle\nabla\mathfrak G_k(z_s),z_{s+1}-z_s\rangle
  +
  \frac{L_k}{2}\lVert z_{s+1}-z_s\rVert^2.
\]
代入
\[
  z_{s+1}-z_s
  =
  -\alpha_k\nabla\mathfrak G_k(z_s),
\]
得到
\[
  \mathfrak G_k(z_{s+1})
  \le
  \mathfrak G_k(z_s)
  -
  \alpha_k\lVert\nabla\mathfrak G_k(z_s)\rVert^2
  +
  \frac{L_k\alpha_k^2}{2}
  \lVert\nabla\mathfrak G_k(z_s)\rVert^2.
\]
整理得
\[
  \mathfrak G_k(z_{s+1})
  \le
  \mathfrak G_k(z_s)
  -
  \alpha_k
  \left(
  1-\frac{L_k\alpha_k}{2}
  \right)
  \lVert\nabla\mathfrak G_k(z_s)\rVert^2.
\]
若
\[
  0<\alpha_k\le\frac1{L_k},
\]
则
\[
  1-\frac{L_k\alpha_k}{2}\ge\frac12.
\]
因此
\[
  \mathfrak G_k(z_{s+1})
  \le
  \mathfrak G_k(z_s)
  -
  \frac{\alpha_k}{2}
  \lVert\nabla\mathfrak G_k(z_s)\rVert^2.
\]
证毕。
\end{proof}

\begin{proposition}[带投影约束的下降口径]
\label{prop:tex55-jacgap-projected-descent}
若采用投影更新
\[
  z_{s+1}
  =
  \Pi_{\mathcal K_k}
  \left[
  z_s-\alpha_k\nabla\mathfrak G_k(z_s)
  \right],
\]
则严格下降结论需要约束集 \(\mathcal K_k\) 的凸性、投影非扩张性以及相应的一阶最优性条件。
在这些条件满足时，投影梯度下降以
\[
  \mathfrak G_k
\]
作为每阶同伦扭曲的收敛指引；
在不满足这些条件时，投影更新只能作为启发式参数修正，不能直接推出单调下降。
\end{proposition}

\begin{proof}
投影算子若为闭凸集上的欧氏投影，则具有非扩张性，并满足投影变分不等式。
这些条件可把无约束下降引理中的方向
\[
  -\nabla\mathfrak G_k(z_s)
\]
替换为可行投影步方向，从而得到投影梯度的标准充分下降式。
若约束集非凸或投影不唯一，则该变分不等式不再自动成立；
此时不能从 \(L_k\)-光滑性单独推出单调下降。
故命题成立。
\end{proof}

\subsection{定理：高阶幂化加速}
\label{subsec:tex55-jacgap-high-order-power}

\begin{definition}[每阶幂化因子]
\label{def:tex55-jacgap-power-factor}
第 \(k\) 阶完成后，若存在
\[
  A_k>1
\]
使得
\[
  Q_T(\Gamma^{[k]})
  \le
  Q_T(\Gamma^{[k-1]})^{A_k},
\]
则称第 \(k\) 阶实现幂化加速，幂化因子为 \(A_k\)。
等价地，
\[
  \mathcal C_T(\Gamma^{[k]})
  \ge
  A_k\mathcal C_T(\Gamma^{[k-1]}).
\]
\end{definition}

\begin{theorem}[依次正交的高阶幂化加速]
\label{thm:tex55-jacgap-high-order-power}
若对每一阶
\[
  k=1,\ldots,n
\]
都有
\[
  Q_T(\Gamma^{[k]})
  \le
  Q_T(\Gamma^{[k-1]})^{A_k},
  \qquad
  A_k>1,
\]
则
\[
  \boxed{
  Q_T(\Gamma^{[n]})
  \le
  Q_T(\Gamma^{[0]})^{A_1A_2\cdots A_n}.
  }
\]
因此
\[
  \boxed{
  \mathcal J_T(\Gamma^{[n]})
  \ge
  1-
  Q_T(\Gamma^{[0]})^{A_1A_2\cdots A_n}.
  }
\]
\end{theorem}

\begin{Certificate}[数学归纳法证书]
\label{cert:tex55-jacgap-high-order-power}
目标谓词为
\[
  P(n):
  Q_T(\Gamma^{[n]})
  \le
  Q_T(\Gamma^{[0]})^{\prod_{k=1}^{n}A_k}.
\]

基例：当 \(n=1\) 时，由假设
\[
  Q_T(\Gamma^{[1]})
  \le
  Q_T(\Gamma^{[0]})^{A_1}
\]
可知 \(P(1)\) 成立。

归纳步：假设 \(P(n)\) 成立，即
\[
  Q_T(\Gamma^{[n]})
  \le
  Q_T(\Gamma^{[0]})^{\prod_{k=1}^{n}A_k}.
\]
第 \(n+1\) 阶满足
\[
  Q_T(\Gamma^{[n+1]})
  \le
  Q_T(\Gamma^{[n]})^{A_{n+1}}.
\]
代入归纳假设得到
\[
  Q_T(\Gamma^{[n+1]})
  \le
  \left(
  Q_T(\Gamma^{[0]})^{\prod_{k=1}^{n}A_k}
  \right)^{A_{n+1}}
  =
  Q_T(\Gamma^{[0]})^{\prod_{k=1}^{n+1}A_k}.
\]
所以 \(P(n+1)\) 成立。
由数学归纳法，命题对所有 \(n\ge1\) 成立。
\end{Certificate}

\begin{proof}
证书 \ref{cert:tex55-jacgap-high-order-power} 已给出完整归纳证明。
由
\[
  \mathcal J_T(\Gamma)=1-Q_T(\Gamma)
\]
得到
\[
  \mathcal J_T(\Gamma^{[n]})
  =
  1-Q_T(\Gamma^{[n]})
  \ge
  1-Q_T(\Gamma^{[0]})^{A_1A_2\cdots A_n}.
\]
证毕。
\end{proof}

\subsection{总定理、谓词因果链与边界}
\label{subsec:tex55-jacgap-total-chain-boundary}

\begin{theorem}[高阶正交雅可比间隙下降推出高阶幂化命中加速]
\label{thm:tex55-jacgap-total-theorem}
若满足以下条件：
\[
  \text{一、}\quad
  \mathcal C_T\text{ 可微};
\]
\[
  \text{二、}\quad
  b_k=
  \Pi_{\mathcal W_{k-1}^{\perp}}
  \nabla\mathcal C_T(\Gamma^{[k-1]})
  \neq0;
\]
\[
  \text{三、}\quad
  \text{每阶雅可比间隙下降得到 }
  \delta_k=J_k(\theta_k)d_k
  \text{ 且 }
  \lVert\delta_k-b_k\rVert
  <
  \frac{\lVert b_k\rVert^2}
  {\lVert\nabla\mathcal C_T(\Gamma^{[k-1]})\rVert};
\]
\[
  \text{四、}\quad
  \text{每阶步长 }\eta_k>0\text{ 足够小},
\]
则每阶都有
\[
  \mathcal C_T(\Gamma^{[k]})
  >
  \mathcal C_T(\Gamma^{[k-1]}),
\]
从而
\[
  Q_T(\Gamma^{[k]})
  <
  Q_T(\Gamma^{[k-1]}).
\]
若进一步每阶满足幂化压缩条件
\[
  Q_T(\Gamma^{[k]})
  \le
  Q_T(\Gamma^{[k-1]})^{A_k},
  \qquad
  A_k>1,
\]
则
\[
  \boxed{
  Q_T(\Gamma^{[n]})
  \le
  Q_T(\Gamma^{[0]})^{A_1A_2\cdots A_n}
  }
\]
以及
\[
  \boxed{
  \mathcal J_T(\Gamma^{[n]})
  \ge
  1-
  Q_T(\Gamma^{[0]})^{A_1A_2\cdots A_n}.
  }
\]
\end{theorem}

\begin{Certificate}[总定理证书]
\label{cert:tex55-jacgap-total-theorem}
总链条分两段。
第一段由雅可比逼近推出单阶命中加速：
\[
  \mathfrak G_k\downarrow
  \Longrightarrow
  J_k(\theta_k)d_k\approx b_k
  \Longrightarrow
  D\mathcal C_T[J_k(\theta_k)d_k]>0
  \Longrightarrow
  Q_T(\Gamma^{[k]})<Q_T(\Gamma^{[k-1]}).
\]
第二段由每阶幂压缩推出高阶幂化：
\[
  Q_T(\Gamma^{[k]})
  \le
  Q_T(\Gamma^{[k-1]})^{A_k}
  \Longrightarrow
  Q_T(\Gamma^{[n]})
  \le
  Q_T(\Gamma^{[0]})^{\prod_{k=1}^{n}A_k}.
\]
\end{Certificate}

\begin{proof}
由条件二、三与定理 \ref{thm:tex55-jacgap-growth}，每一阶实际雅可比方向
\[
  \delta_k=J_k(\theta_k)d_k
\]
满足
\[
  D\mathcal C_T(\Gamma^{[k-1]})[\delta_k]>0.
\]
由条件四，取足够小正步长 \(\eta_k\) 可得
\[
  \mathcal C_T(\Gamma^{[k]})
  >
  \mathcal C_T(\Gamma^{[k-1]}).
\]
因此
\[
  Q_T(\Gamma^{[k]})
  =
  e^{-\mathcal C_T(\Gamma^{[k]})}
  <
  e^{-\mathcal C_T(\Gamma^{[k-1]})}
  =
  Q_T(\Gamma^{[k-1]}).
\]
若进一步每阶满足幂化压缩条件，则由定理
\ref{thm:tex55-jacgap-high-order-power} 得到
\[
  Q_T(\Gamma^{[n]})
  \le
  Q_T(\Gamma^{[0]})^{A_1A_2\cdots A_n},
\]
并推出命中泛函下界。
证毕。
\end{proof}

\begin{remark}[完整谓词因果链]
本章的谓词因果链为
\[
\begin{aligned}
  &\text{命中泛函随机过程轨迹函数 }\Gamma\\
  &\Longrightarrow
  \text{命中指数 }\mathcal C_T(\Gamma)=-\log Q_T(\Gamma)\\
  &\Longrightarrow
  \text{命中梯度 }g_\Gamma=\nabla\mathcal C_T(\Gamma)\\
  &\Longrightarrow
  \text{第 }k\text{ 阶已用方向空间 }\mathcal W_{k-1}\\
  &\Longrightarrow
  \text{高阶正交命中残差 }
  b_k=\Pi_{\mathcal W_{k-1}^{\perp}}g_{\Gamma^{[k-1]}}\\
  &\Longrightarrow
  \text{同伦扭曲雅可比 }J_k(\theta_k)d_k\\
  &\Longrightarrow
  \text{雅可比间隙 }\mathfrak G_k(\theta_k,d_k)\\
  &\Longrightarrow
  \text{每阶雅可比间隙梯度下降}\\
  &\Longrightarrow
  J_k(\theta_k)d_k\approx b_k\\
  &\Longrightarrow
  D\mathcal C_T[J_k(\theta_k)d_k]>0\\
  &\Longrightarrow
  Q_T(\Gamma^{[k]})<Q_T(\Gamma^{[k-1]})\\
  &\Longrightarrow
  \text{若满足幂压缩，则 }
  Q_T(\Gamma^{[k]})
  \le
  Q_T(\Gamma^{[k-1]})^{A_k}\\
  &\Longrightarrow
  Q_T(\Gamma^{[n]})
  \le
  Q_T(\Gamma^{[0]})^{\prod_{k=1}^{n}A_k}\\
  &\Longrightarrow
  \text{高阶幂化加速命中。}
\end{aligned}
\]
\end{remark}

\begin{remark}[必要边界]
第一，高阶正交不自动等于加速。
若
\[
  b_k=
  \Pi_{\mathcal W_{k-1}^{\perp}}g_{\Gamma^{[k-1]}}
  =
  0,
\]
则第 \(k\) 阶正交补空间中没有剩余命中梯度；
继续做高阶正交扭曲不能推出命中加速。

第二，雅可比间隙下降不自动等于命中加速。
只有当实际雅可比方向足够接近理想方向，即
\[
  \lVert\delta_k-b_k\rVert
  <
  \frac{\lVert b_k\rVert^2}{\lVert g_k\rVert},
\]
才能保证
\[
  D\mathcal C_T[\delta_k]>0.
\]
若间隙仍然很大，实际方向可能偏离理想命中方向，甚至降低命中概率。

第三，单调加速不等于幂化加速。
若只证明
\[
  Q_T(\Gamma^{[k]})<Q_T(\Gamma^{[k-1]}),
\]
只能得到单调改善。
要证明幂化加速，必须额外满足
\[
  Q_T(\Gamma^{[k]})
  \le
  Q_T(\Gamma^{[k-1]})^{A_k},
  \qquad A_k>1.
\]
因此严谨说法是
\[
  \boxed{
  \text{高阶正交}
  +
  \text{雅可比间隙下降}
  \Longrightarrow
  \text{命中加速};
  }
\]
\[
  \boxed{
  \text{若再满足每阶幂压缩条件}
  \Longrightarrow
  \text{高阶幂化命中加速。}
  }
\]
\end{remark}

\begin{remark}[最终严格表述]
用户命题可最终写成：
\[
  \boxed{
  \begin{gathered}
  \text{在第 }k\text{ 阶，将命中指数梯度投影到目标扭曲}\\
  \text{与前序 }k-1\text{ 阶扭曲方向的正交补空间，}\\
  \text{得到理想高阶正交命中方向 }b_k.
  \end{gathered}
  }
\]
\[
  \boxed{
  \begin{gathered}
  \text{再通过最小化雅可比间隙 }\mathfrak G_k,\\
  \text{使第 }k\text{ 阶同伦扭曲的实际雅可比方向 }J_kd_k
  \text{ 收敛到 }b_k.
  \end{gathered}
  }
\]
\[
  \boxed{
  \begin{gathered}
  \text{若雅可比间隙足够小，则该阶同伦扭曲提高命中指数、降低未命中概率；}\\
  \text{若每阶还满足幂压缩，则形成高阶幂化加速命中。}
  \end{gathered}
  }
\]
最压缩的数学表达是
\[
  \boxed{
  b_k
  =
  \Pi_{\mathcal W_{k-1}^{\perp}}
  \nabla\mathcal C_T(\Gamma^{[k-1]}).
  }
\]
\[
  \boxed{
  \mathfrak G_k(\theta_k,d_k)
  =
  \frac12
  \left\lVert
  \Pi_{\mathcal W_{k-1}^{\perp}}
  J_k(\theta_k)d_k
  -
  b_k
  \right\rVert^2
  +
  \frac{\rho_k}{2}
  \left\lVert
  \Pi_{\mathcal W_{k-1}}
  J_k(\theta_k)d_k
  \right\rVert^2.
  }
\]
\[
  \boxed{
  (\theta_k,d_k)
  \leftarrow
  (\theta_k,d_k)
  -
  \alpha_k
  \nabla\mathfrak G_k(\theta_k,d_k).
  }
\]
\[
  \boxed{
  Q_T(\Gamma^{[n]})
  \le
  Q_T(\Gamma^{[0]})^{A_1A_2\cdots A_n}.
  }
\]
这就是“高阶正交命中加速沿每阶雅可比间隙梯度下降，并以雅可比间隙作为每阶同伦扭曲的收敛指引”的严格形式。
\end{remark}

\clearpage

%%%%%%%%%%%%%%%%%%%%%%%%%%%%%%%%%%%%%%%%%%%%%%%%%%%%%%%%%%%%
\section{形式逻辑：从局部最优算法到高阶正交同伦命中轨迹的总体增益}
\label{sec:tex55-overall-gain-evaluation}
%%%%%%%%%%%%%%%%%%%%%%%%%%%%%%%%%%%%%%%%%%%%%%%%%%%%%%%%%%%%

\begin{remark}[总评价]
前述讨论的核心增益是：把“局部最优算法如何走向全局命中”这个容易停留在直觉层面的说法，
逐步提升为一个可定义、可分解、可度量、可优化且具有可证明边界的结构。
压缩成一句话：
\[
  \boxed{
  \begin{gathered}
  \text{增益}\\
  =
  \text{把“寻找最优算法”转化为}\\
  \text{“优化命中泛函随机过程的高阶正交同伦轨迹”。}
  \end{gathered}
  }
\]
更完整地说：
\[
  \boxed{
  \begin{gathered}
  \text{时域局部搜索}
  \Longrightarrow
  \text{分布参数化}
  \Longrightarrow
  \text{随机变量簇}\\
  \Longrightarrow
  \text{命中泛函}
  \Longrightarrow
  \text{轨迹函数}
  \Longrightarrow
  \text{正交命中加速}\\
  \Longrightarrow
  \text{雅可比间隙收敛指引}
  \Longrightarrow
  \text{高阶幂化命中加速。}
  \end{gathered}
  }
\]
它比“用频域全局优化替代时域局部优化”更严格，因为它同时回答：
优化对象是什么、命中如何定义、加速如何度量、每阶扭曲如何收敛。
\end{remark}

\subsection{对象层级的增益}
\label{subsec:tex55-gain-object-level}

\begin{proposition}[增益 G1：区分分布律与命中随机变量]
\label{prop:tex55-gain-distribution-rv}
原始表达容易写成
\[
  \max_{\lambda}\pi_\lambda(S).
\]
严格对象层级应写成
\[
  X_t^\lambda\sim\pi_\lambda,
  \qquad
  Y_t^\lambda=\mathbf 1_S(X_t^\lambda),
\]
以及
\[
  \mathcal J_T(\Gamma)
  =
  \mathbb P_\Gamma(\exists t\le T:Y_t^\Gamma=1).
\]
因此增益是：
\[
  \boxed{
  \text{从“优化分布”提升为“优化由分布生成的命中随机变量簇”。}
  }
\]
\end{proposition}

\begin{Certificate}[谓词因果链证书]
\label{cert:tex55-gain-distribution-rv}
本命题的谓词链为
\[
\begin{aligned}
  &\pi_\lambda
  \text{ 是分布律}\\
  &\Longrightarrow
  X_t^\lambda\sim\pi_\lambda
  \text{ 是承载取值的随机变量}\\
  &\Longrightarrow
  X_t^\lambda\in S
  \text{ 是命中事件}\\
  &\Longrightarrow
  Y_t^\lambda=\mathbf 1_S(X_t^\lambda)
  \text{ 是命中指示变量}\\
  &\Longrightarrow
  \pi_\lambda(S)=\mathbb P(X_t^\lambda\in S)
  =
  \mathbb E[Y_t^\lambda].
\end{aligned}
\]
\end{Certificate}

\begin{proof}
分布律 \(\pi_\lambda\) 只是随机变量 \(X_t^\lambda\) 的法律。
真正发生命中的对象不是 \(\pi_\lambda\) 本身，而是随机变量取值满足
\[
  X_t^\lambda\in S.
\]
由命中指示变量定义
\[
  Y_t^\lambda=\mathbf 1_S(X_t^\lambda),
\]
可知
\[
  Y_t^\lambda=1
  \Longleftrightarrow
  X_t^\lambda\in S.
\]
因此
\[
  \pi_\lambda(S)
  =
  \mathbb P(X_t^\lambda\in S)
  =
  \mathbb P(Y_t^\lambda=1)
  =
  \mathbb E[Y_t^\lambda].
\]
所以“分布优化”只是“命中随机变量簇优化”的降维表达。
证毕。
\end{proof}

\subsection{泛函层级的增益}
\label{subsec:tex55-gain-functional-level}

\begin{proposition}[增益 G2：从单次命中提升为命中泛函过程]
\label{prop:tex55-gain-functional-process}
单次命中概率
\[
  p_\lambda=\mathbb P(X^\lambda\in S)
\]
被提升为有限时间命中泛函
\[
  \mathcal J_T(\Gamma)
  =
  \mathbb P_\Gamma(\tau_S^\Gamma\le T),
\]
其中
\[
  \tau_S^\Gamma
  =
  \inf\{t\ge1:X_t^\Gamma\in S\}.
\]
进一步引入生存概率
\[
  Q_T(\Gamma)
  =
  \mathbb P_\Gamma(\tau_S^\Gamma>T),
\]
以及命中指数
\[
  \mathcal C_T(\Gamma)
  =
  -\log Q_T(\Gamma).
\]
该层级增益是：
\[
  \boxed{
  \text{把“命中一次”变成“优化首次命中时间与未命中生存概率”。}
  }
\]
\end{proposition}

\begin{Certificate}[符号推演证书]
\label{cert:tex55-gain-functional-process}
目标谓词为
\[
  \mathcal J_T(\Gamma)
  =
  1-e^{-\mathcal C_T(\Gamma)}.
\]
证明链为
\[
  \mathcal J_T
  =
  \mathbb P(\tau_S\le T)
  =
  1-\mathbb P(\tau_S>T)
  =
  1-Q_T
  =
  1-e^{-\mathcal C_T}.
\]
\end{Certificate}

\begin{proof}
由定义
\[
  \mathcal J_T(\Gamma)
  =
  \mathbb P_\Gamma(\tau_S^\Gamma\le T).
\]
其补事件为 \(\{\tau_S^\Gamma>T\}\)，所以
\[
  \mathcal J_T(\Gamma)
  =
  1-
  \mathbb P_\Gamma(\tau_S^\Gamma>T)
  =
  1-Q_T(\Gamma).
\]
由命中指数定义
\[
  \mathcal C_T(\Gamma)=-\log Q_T(\Gamma),
\]
得到
\[
  Q_T(\Gamma)=e^{-\mathcal C_T(\Gamma)}.
\]
因此
\[
  \mathcal J_T(\Gamma)
  =
  1-e^{-\mathcal C_T(\Gamma)}.
\]
函数 \(c\mapsto1-e^{-c}\) 严格递增，故
\[
  \max\mathcal J_T(\Gamma)
  \Longleftrightarrow
  \max\mathcal C_T(\Gamma).
\]
证毕。
\end{proof}

\subsection{轨迹函数层级的增益}
\label{subsec:tex55-gain-trajectory-level}

\begin{proposition}[增益 G3：从选择算法转化为选择同伦扭曲轨迹]
\label{prop:tex55-gain-trajectory-selection}
原始问题可写成
\[
  \max_{A\in\mathcal A}\mathcal J_T(A).
\]
经严格化后，问题变成
\[
  \max_{\Gamma}\mathcal J_T(\Gamma),
\]
其中
\[
  \Gamma=(\Gamma_1,\ldots,\Gamma_T),
  \qquad
  \lambda_t=\Gamma_t(\mathcal F_{t-1}),
\]
并且
\[
  X_t^\Gamma\mid\mathcal F_{t-1}\sim\pi_{\lambda_t}.
\]
该层级增益是：
\[
  \boxed{
  \text{算法不再只是固定规则，而是参数化分布律在时间上的同伦轨迹函数。}
  }
\]
\end{proposition}

\begin{Certificate}[轨迹诱导证书]
\label{cert:tex55-gain-trajectory-selection}
轨迹函数的诱导链为
\[
\begin{aligned}
  &\Gamma\\
  &\Longrightarrow
  \lambda_t=\Gamma_t(\mathcal F_{t-1})\\
  &\Longrightarrow
  \pi_{\lambda_t}\\
  &\Longrightarrow
  X_t^\Gamma\mid\mathcal F_{t-1}\sim\pi_{\lambda_t}\\
  &\Longrightarrow
  Y_t^\Gamma=\mathbf 1_S(X_t^\Gamma)\\
  &\Longrightarrow
  \mathcal J_T(\Gamma).
\end{aligned}
\]
\end{Certificate}

\begin{proof}
每一步同伦扭曲由轨迹函数选择参数 \(\lambda_t\)。
参数 \(\lambda_t\) 决定分布律 \(\pi_{\lambda_t}\)。
分布律生成随机变量 \(X_t^\Gamma\)。
随机变量决定命中指示变量
\[
  Y_t^\Gamma=\mathbf 1_S(X_t^\Gamma).
\]
最终由
\[
  (Y_1^\Gamma,\ldots,Y_T^\Gamma)
\]
决定
\[
  \mathcal J_T(\Gamma)
  =
  \mathbb P_\Gamma(\exists t\le T:Y_t^\Gamma=1).
\]
所以优化算法可转化为优化轨迹函数
\[
  \Gamma^\star\in\arg\max_\Gamma\mathcal J_T(\Gamma).
\]
证毕。
\end{proof}

\subsection{正交结构与高阶残差的增益}
\label{subsec:tex55-gain-orthogonal-level}

\begin{proposition}[增益 G4：加速分解为目标正交与命中梯度投影]
\label{prop:tex55-gain-orthogonal-projection}
设
\[
  \mathcal W_0=\operatorname{span}\{u_0\},
\]
其中 \(u_0\) 为目标扭曲方向。
有效的一阶正交命中方向不是任意正交方向，而是
\[
  b_1
  =
  \Pi_{\mathcal W_0^\perp}
  \nabla\mathcal C_T(\Gamma).
\]
若 \(b_1\neq0\)，则
\[
  \widehat b_1=\frac{b_1}{\lVert b_1\rVert}
\]
满足
\[
  D\mathcal C_T(\Gamma)[\widehat b_1]
  =
  \lVert b_1\rVert>0.
\]
该层级增益是：
\[
  \boxed{
  \text{把“正交扭曲”从几何描述变成可计算的投影梯度方向。}
  }
\]
\end{proposition}

\begin{Certificate}[投影梯度证书]
\label{cert:tex55-gain-orthogonal-projection}
在约束 \(\delta\in\mathcal W_0^\perp\) 下，
\[
  D\mathcal C_T(\Gamma)[\delta]
  =
  \langle\nabla\mathcal C_T(\Gamma),\delta\rangle
  =
  \langle
  \Pi_{\mathcal W_0^\perp}\nabla\mathcal C_T(\Gamma),
  \delta
  \rangle.
\]
所以 Cauchy--Schwarz 不等式给出最优方向为归一化投影方向。
\end{Certificate}

\begin{proof}
若方向 \(\delta\in\mathcal W_0^\perp\)，则它与目标扭曲方向正交。
命中指数的一阶变化为
\[
  D\mathcal C_T(\Gamma)[\delta]
  =
  \langle\nabla\mathcal C_T(\Gamma),\delta\rangle.
\]
由于 \(\delta\in\mathcal W_0^\perp\)，有
\[
  \langle\nabla\mathcal C_T(\Gamma),\delta\rangle
  =
  \langle
  \Pi_{\mathcal W_0^\perp}\nabla\mathcal C_T(\Gamma),
  \delta
  \rangle.
\]
记
\[
  b_1=
  \Pi_{\mathcal W_0^\perp}\nabla\mathcal C_T(\Gamma).
\]
在单位方向 \(\lVert\delta\rVert=1\) 下，Cauchy--Schwarz 给出最大值
\[
  \langle b_1,\delta\rangle\le\lVert b_1\rVert,
\]
等号方向为
\[
  \delta=\widehat b_1=\frac{b_1}{\lVert b_1\rVert}.
\]
若 \(b_1\neq0\)，则该方向给出正的一阶命中指数增长。
证毕。
\end{proof}

\begin{proposition}[增益 G5：从一次正交加速推广为高阶正交残差加速]
\label{prop:tex55-gain-higher-order-residual}
第 \(k\) 阶已用方向空间为
\[
  \mathcal W_{k-1}
  =
  \operatorname{span}\{u_0,\delta_1,\ldots,\delta_{k-1}\}.
\]
第 \(k\) 阶命中残差方向为
\[
  b_k
  =
  \Pi_{\mathcal W_{k-1}^\perp}
  \nabla\mathcal C_T(\Gamma^{[k-1]}).
\]
该层级增益是：
\[
  \boxed{
  \begin{gathered}
  \text{每一阶都避开已经使用过的命中加速方向，}\\
  \text{并在剩余正交补空间中寻找新的命中增益。}
  \end{gathered}
  }
\]
\end{proposition}

\begin{proof}
第 \(k-1\) 阶之前已经使用
\[
  \delta_1,\ldots,\delta_{k-1}.
\]
若第 \(k\) 阶继续沿这些方向，则只是重复已有加速方向，不能形成独立高阶增益。
因此要求
\[
  \delta_k\in\mathcal W_{k-1}^\perp.
\]
在该约束下，命中指数最大一阶增长方向由投影给出：
\[
  b_k
  =
  \Pi_{\mathcal W_{k-1}^\perp}
  \nabla\mathcal C_T(\Gamma^{[k-1]}).
\]
若 \(b_k\neq0\)，则存在新的正交命中增益方向。
这相当于把命中梯度分解为
\[
  \nabla\mathcal C_T
  =
  \text{已用分量}
  +
  \text{剩余可加速分量}.
\]
证毕。
\end{proof}

\subsection{雅可比间隙与幂化加速的增益}
\label{subsec:tex55-gain-jacobian-power}

\begin{proposition}[增益 G6：雅可比间隙提供每阶同伦扭曲的收敛指引]
\label{prop:tex55-gain-jacobian-gap}
第 \(k\) 阶理想方向为 \(b_k\)，实际同伦扭曲由参数控制，实际可实现方向为
\[
  J_k(\theta_k)d_k.
\]
因此引入雅可比间隙
\[
  \mathfrak G_k(\theta_k,d_k)
  =
  \frac12
  \left\lVert
  \Pi_{\mathcal W_{k-1}^{\perp}}
  J_k(\theta_k)d_k
  -
  b_k
  \right\rVert^2
  +
  \frac{\rho_k}{2}
  \left\lVert
  \Pi_{\mathcal W_{k-1}}
  J_k(\theta_k)d_k
  \right\rVert^2.
\]
该层级增益是：
\[
  \boxed{
  \begin{gathered}
  \text{把“我想沿某个高阶命中方向扭曲”变成}\\
  \text{“实际参数雅可比如何逼近该方向”的可优化问题。}
  \end{gathered}
  }
\]
\end{proposition}

\begin{Certificate}[雅可比间隙证书]
\label{cert:tex55-gain-jacobian-gap}
雅可比间隙分解为
\[
  \mathfrak G_k
  =
  \text{命中方向逼近误差}
  +
  \text{高阶正交泄漏误差}.
\]
其中第一项控制
\[
  \Pi_{\mathcal W_{k-1}^\perp}J_kd_k\approx b_k,
\]
第二项控制
\[
  \Pi_{\mathcal W_{k-1}}J_kd_k\approx0.
\]
\end{Certificate}

\begin{proof}
第 \(k\) 阶同伦扭曲参数为 \(\theta_k\)，参数速度为 \(d_k\)。
它在轨迹函数空间中产生的实际一阶变化为
\[
  J_k(\theta_k)d_k.
\]
理想高阶正交命中方向为 \(b_k\)。
若实际方向偏离 \(b_k\)，则实际扭曲可能无法提高命中指数。
因此需要最小化
\[
  \left\lVert
  \Pi_{\mathcal W_{k-1}^{\perp}}
  J_k(\theta_k)d_k
  -
  b_k
  \right\rVert^2.
\]
同时，为防止实际方向泄漏到已用方向空间，需要惩罚
\[
  \left\lVert
  \Pi_{\mathcal W_{k-1}}
  J_k(\theta_k)d_k
  \right\rVert^2.
\]
二者合成 \(\mathfrak G_k\)，于是雅可比间隙提供每阶同伦扭曲的可下降目标。
证毕。
\end{proof}

\begin{theorem}[增益 G7：幂化加速给出更快命中的量化形式]
\label{thm:tex55-gain-power-acceleration}
若每阶满足
\[
  Q_T(\Gamma^{[k]})
  \le
  Q_T(\Gamma^{[k-1]})^{A_k},
  \qquad
  A_k>1,
\]
则经过 \(n\) 阶后
\[
  \boxed{
  Q_T(\Gamma^{[n]})
  \le
  Q_T(\Gamma^{[0]})^{A_1A_2\cdots A_n}.
  }
\]
于是
\[
  \boxed{
  \mathcal J_T(\Gamma^{[n]})
  \ge
  1-
  Q_T(\Gamma^{[0]})^{A_1A_2\cdots A_n}.
  }
\]
该层级增益是：
\[
  \boxed{
  \text{把“更快命中”量化为“未命中生存概率的幂化压缩”。}
  }
\]
\end{theorem}

\begin{Certificate}[数学归纳法证书]
\label{cert:tex55-gain-power-acceleration}
目标谓词为
\[
  P(n):
  Q_T(\Gamma^{[n]})
  \le
  Q_T(\Gamma^{[0]})^{\prod_{k=1}^{n}A_k}.
\]

基例：当 \(n=1\) 时，由假设
\[
  Q_T(\Gamma^{[1]})
  \le
  Q_T(\Gamma^{[0]})^{A_1}
\]
得 \(P(1)\) 成立。

归纳步：若 \(P(n)\) 成立，且
\[
  Q_T(\Gamma^{[n+1]})
  \le
  Q_T(\Gamma^{[n]})^{A_{n+1}},
\]
则
\[
  Q_T(\Gamma^{[n+1]})
  \le
  \left(
  Q_T(\Gamma^{[0]})^{\prod_{k=1}^{n}A_k}
  \right)^{A_{n+1}}
  =
  Q_T(\Gamma^{[0]})^{\prod_{k=1}^{n+1}A_k}.
\]
故 \(P(n+1)\) 成立。
\end{Certificate}

\begin{proof}
证书 \ref{cert:tex55-gain-power-acceleration} 给出归纳链，故
\[
  Q_T(\Gamma^{[n]})
  \le
  Q_T(\Gamma^{[0]})^{A_1A_2\cdots A_n}.
\]
又因为
\[
  \mathcal J_T(\Gamma)=1-Q_T(\Gamma),
\]
所以
\[
  \mathcal J_T(\Gamma^{[n]})
  =
  1-Q_T(\Gamma^{[n]})
  \ge
  1-
  Q_T(\Gamma^{[0]})^{A_1A_2\cdots A_n}.
\]
证毕。
\end{proof}

\subsection{总体方法论增益}
\label{subsec:tex55-gain-methodology}

\begin{remark}[从局部最优到命中过程]
原问题是
\[
  \text{如何避免局部最优？}
\]
严格化后变成
\[
  \text{如何使首次命中时间 }\tau_S\text{ 尽可能小？}
\]
这是对象上的增益：局部最优是状态性质，而首次命中时间是过程性质。
优化对象从单个状态 \(x_t\) 提升为随机过程
\[
  (X_t^\Gamma)_{t\le T}.
\]
\end{remark}

\begin{remark}[从算法优选到轨迹函数优选]
原问题可写成
\[
  \max_A\mathcal J_T(A).
\]
转化后写成
\[
  \max_\Gamma\mathcal J_T(\Gamma).
\]
其中 \(\Gamma\) 可控制
\[
  \beta_t,\quad
  w_t,\quad
  K_t,\quad
  \theta_t,\quad
  \text{频域模态},\quad
  \text{同伦扭曲路径}.
\]
这使算法设计从选择固定规则，提升为选择参数轨迹。
\end{remark}

\begin{remark}[从直觉加速到梯度--投影--雅可比间隙]
加速方向不再凭直觉给出，而是由三层对象确定：
\[
  g_k=\nabla\mathcal C_T(\Gamma^{[k-1]}),
\]
\[
  b_k=\Pi_{\mathcal W_{k-1}^\perp}g_k,
\]
\[
  J_kd_k\approx b_k.
\]
于是每阶设计都有明确的命中梯度、正交残差和雅可比实现方向。
这使“同伦扭曲”从概念变成算法步骤。
\end{remark}

\begin{remark}[从是否有效到间隙是否足够小]
雅可比间隙判据为
\[
  \lVert\delta_k-b_k\rVert
  <
  \frac{\lVert b_k\rVert^2}{\lVert g_k\rVert},
  \qquad
  \delta_k=J_kd_k.
\]
若该条件成立，则
\[
  D\mathcal C_T[\delta_k]>0.
\]
这说明每阶雅可比间隙承担命中加速是否发生的判定功能。
\end{remark}

\subsection{主要理论增益与综合结论}
\label{subsec:tex55-gain-main-conclusions}

\begin{theorem}[总增益定理]
\label{thm:tex55-overall-gain-theorem}
前述形式系统把局部最优算法选择问题提升为以下对象链：
\[
  \boxed{
  \begin{gathered}
  \text{局部最优算法}
  \Longrightarrow
  \text{分布参数}
  \Longrightarrow
  \text{随机变量簇}\\
  \Longrightarrow
  \text{命中泛函}
  \Longrightarrow
  \text{轨迹函数}
  \Longrightarrow
  \text{命中指数梯度}\\
  \Longrightarrow
  \text{高阶正交残差}
  \Longrightarrow
  \text{雅可比间隙下降}
  \Longrightarrow
  \text{高阶幂化命中加速。}
  \end{gathered}
  }
\]
其理论增益是
\[
  \boxed{
  \begin{gathered}
  \text{把“如何更快找到全局可命中区域”转化为}\\
  \text{“如何在轨迹函数空间中降低未命中生存概率”的问题。}
  \end{gathered}
  }
\]
其算法增益是
\[
  \boxed{
  \begin{gathered}
  \text{每阶都有明确的方向 }b_k,\quad
  \text{明确的实现误差 }\mathfrak G_k,\\
  \text{以及明确的收敛指引 }\nabla\mathfrak G_k.
  \end{gathered}
  }
\]
其证明增益是
\[
  \boxed{
  \begin{gathered}
  \text{把“加速直觉”组织为条件定理：}\\
  \text{残差非零、间隙足够小、幂压缩成立}
  \Longrightarrow
  \text{高阶幂化加速。}
  \end{gathered}
  }
\]
\end{theorem}

\begin{Certificate}[总增益谓词链证书]
\label{cert:tex55-overall-gain-theorem}
总谓词链为
\[
\begin{aligned}
  &\text{分布律 }\pi_\lambda\\
  &\Longrightarrow
  \text{随机变量 }X_t^\lambda\\
  &\Longrightarrow
  \text{命中指示 }Y_t^\lambda\\
  &\Longrightarrow
  \text{命中泛函 }\mathcal J_T\\
  &\Longrightarrow
  \text{生存概率 }Q_T\\
  &\Longrightarrow
  \text{命中指数 }\mathcal C_T=-\log Q_T\\
  &\Longrightarrow
  \text{轨迹梯度 }g_k=\nabla\mathcal C_T(\Gamma^{[k-1]})\\
  &\Longrightarrow
  \text{正交残差 }b_k=\Pi_{\mathcal W_{k-1}^\perp}g_k\\
  &\Longrightarrow
  \text{同伦雅可比方向 }J_kd_k\\
  &\Longrightarrow
  \text{雅可比间隙 }\mathfrak G_k\\
  &\Longrightarrow
  \text{若 }\mathfrak G_k\text{ 足够小，则 }D\mathcal C_T[J_kd_k]>0\\
  &\Longrightarrow
  Q_T(\Gamma^{[k]})<Q_T(\Gamma^{[k-1]})\\
  &\Longrightarrow
  \text{若每阶幂压缩，则高阶幂化命中加速。}
\end{aligned}
\]
\end{Certificate}

\begin{proof}
命题 \ref{prop:tex55-gain-distribution-rv} 证明分布律必须提升为命中随机变量簇。
命题 \ref{prop:tex55-gain-functional-process} 证明命中概率可提升为生存概率与命中指数。
命题 \ref{prop:tex55-gain-trajectory-selection} 证明算法选择可转化为轨迹函数选择。
命题 \ref{prop:tex55-gain-orthogonal-projection} 与
\ref{prop:tex55-gain-higher-order-residual} 证明正交命中方向可由命中梯度投影与高阶残差给出。
命题 \ref{prop:tex55-gain-jacobian-gap} 证明雅可比间隙提供实际同伦扭曲逼近理想方向的可优化目标。
定理 \ref{thm:tex55-gain-power-acceleration} 证明每阶幂压缩推出高阶幂化命中加速。
这些命题按证书中的谓词链组合，即得到总增益定理。
证毕。
\end{proof}

\begin{corollary}[从经验启发式到结构化理论]
\label{cor:tex55-gain-heuristic-to-structure}
原本的局部最优算法选择问题常写成
\[
  \text{哪个启发式更有效？}
\]
现在可写成
\[
  \text{哪条同伦参数轨迹能最大化命中指数并最小化雅可比间隙？}
\]
即
\[
  \Gamma^\star\in\arg\max_{\Gamma}\mathcal C_T(\Gamma),
  \qquad
  (\theta_k^\star,d_k^\star)
  \approx
  \arg\min_{\theta_k,d_k}\mathfrak G_k(\theta_k,d_k).
\]
因此问题从经验试错转化为结构化优化。
\end{corollary}

\begin{proof}
由定理 \ref{thm:tex55-overall-gain-theorem}，算法优选被提升为轨迹函数优选；
由雅可比间隙定义，每阶轨迹实现问题被转化为 \(\mathfrak G_k\) 的下降问题。
故结论成立。
\end{proof}

\begin{corollary}[从全局最优断言到概率命中框架]
\label{cor:tex55-gain-global-claim-to-hit}
直接说“算法达到全局最优”通常过强。
严格框架应写成
\[
  \mathbb P(\tau_S^\Gamma\le T)
\]
以及
\[
  Q_T(\Gamma)
  \le
  Q_T(\Gamma^{[0]})^A.
\]
这承认有限时间下是概率命中问题，而不是无条件确定命中。
\end{corollary}

\begin{proof}
由定义
\[
  \mathcal J_T(\Gamma)=\mathbb P(\tau_S^\Gamma\le T)
\]
可知，有限预算下的结论是概率事件。
若有幂化压缩
\[
  Q_T(\Gamma)\le Q_T(\Gamma^{[0]})^A,
\]
则得到的是未命中概率下降与命中概率上升，而不是无条件命中。
证毕。
\end{proof}

\subsection{必要边界与最终综合评价}
\label{subsec:tex55-gain-boundaries-final}

\begin{remark}[必要边界]
第一，高阶正交残差必须非零。
若
\[
  b_k=0,
\]
则第 \(k\) 阶正交补空间中没有剩余命中梯度，继续正交扭曲不能推出加速。

第二，雅可比间隙必须足够小。
需要
\[
  \lVert\delta_k-b_k\rVert
  <
  \frac{\lVert b_k\rVert^2}{\lVert g_k\rVert}.
\]
否则实际同伦扭曲方向可能偏离理想命中方向。

第三，命中加速不等于幂化加速。
若只得到
\[
  Q_T(\Gamma^{[k]})<Q_T(\Gamma^{[k-1]}),
\]
只能说明单调改进。
要得到幂化加速，还需要
\[
  Q_T(\Gamma^{[k]})
  \le
  Q_T(\Gamma^{[k-1]})^{A_k},
  \qquad
  A_k>1.
\]

第四，正交结构需要随分布与轨迹更新。
每阶内积、梯度和投影都依赖当前轨迹或当前分布；
因此 \(\mathcal W_{k-1}^\perp\) 不是固定空间，而应随
\[
  \Gamma^{[k-1]}
\]
重新计算。
\end{remark}

\begin{remark}[最终综合评价]
上述讨论的最大增益不是单个公式，而是完成了对象链升级：
\[
  \boxed{
  \begin{gathered}
  \text{局部最优算法}
  \Longrightarrow
  \text{分布参数}
  \Longrightarrow
  \text{随机变量簇}
  \Longrightarrow
  \text{命中泛函}\\
  \Longrightarrow
  \text{轨迹函数}
  \Longrightarrow
  \text{命中指数梯度}
  \Longrightarrow
  \text{高阶正交残差}\\
  \Longrightarrow
  \text{雅可比间隙下降}
  \Longrightarrow
  \text{高阶幂化命中加速。}
  \end{gathered}
  }
\]
最简结论为
\[
  \boxed{
  \begin{gathered}
  \text{上述讨论的增益，在于把全局命中问题从“选择算法”提升为}\\
  \text{“优化命中泛函随机过程的高阶正交同伦轨迹”，}\\
  \text{并引入雅可比间隙作为每阶扭曲是否收敛到有效命中方向的判据。}
  \end{gathered}
  }
\]
\end{remark}

\clearpage

%%%%%%%%%%%%%%%%%%%%%%%%%%%%%%%%%%%%%%%%%%%%%%%%%%%%%%%%%%%%
\section{形式逻辑：统计力学同伦命中随机过程的工程截断确定性收敛}
\label{sec:tex55-engineering-truncation-deterministic-convergence}
%%%%%%%%%%%%%%%%%%%%%%%%%%%%%%%%%%%%%%%%%%%%%%%%%%%%%%%%%%%%

\begin{remark}[严格对象口径]
本文把“对统计力学同伦扭曲命中的随机性做工程截断确定性收敛”统一拆成两层：
\[
  \boxed{
  \text{有限 }(N,L)\text{ 下：工程截断的确定性近似收敛}
  }
\]
以及
\[
  \boxed{
  N\to\infty,\quad L\to\infty
  \text{ 或 }L^\infty\text{ 一致极限下：数学严格的确定性极限收敛。}
  }
\]
确定性极限的对象限定为经验测度、命中泛函、命中指数、轨迹函数与雅可比间隙过程，即
\[
  \boxed{
  \begin{gathered}
  \widehat{\mathcal J}_{T,N,L}(\Gamma)\to\mathcal J_T(\Gamma),\\
  \widehat{\mathfrak G}_{k,N,L}\to\mathfrak G_k,\qquad
  \Gamma_{N,L}\to\Gamma_*.
  \end{gathered}
  }
\]
其中一致大数定律与经验过程的标准背景可参考
\cite{VanDerVaartWellner1996EmpiricalProcesses}；
本章把所需条件显式列出，不把随机性消失当作无条件结论。
\end{remark}

\subsection{对象域：随机命中过程与经验命中泛函}
\label{subsec:tex55-engdet-objects}

\begin{definition}[命中随机过程]
\label{def:tex55-engdet-hit-process}
令
\[
  X_t^\Gamma:\Xi\to\Omega
\]
为由同伦轨迹函数 \(\Gamma\) 诱导的统计力学随机过程。
令命中集合为
\[
  S\subseteq\Omega.
\]
命中指示变量为
\[
  Y_t^\Gamma=\mathbf 1_S(X_t^\Gamma).
\]
首次命中时间为
\[
  \tau_S^\Gamma
  =
  \inf\{t\ge1:Y_t^\Gamma=1\}.
\]
有限时间命中泛函为
\[
  \mathcal J_T(\Gamma)
  =
  \mathbb P_\Gamma(\tau_S^\Gamma\le T).
\]
未命中生存概率与命中指数分别为
\[
  Q_T(\Gamma)
  =
  \mathbb P_\Gamma(\tau_S^\Gamma>T),
  \qquad
  \mathcal C_T(\Gamma)
  =
  -\log Q_T(\Gamma).
\]
于是
\[
  \mathcal J_T(\Gamma)
  =
  1-e^{-\mathcal C_T(\Gamma)}.
\]
\end{definition}

\begin{definition}[有限 \(N\) 的经验命中泛函]
\label{def:tex55-engdet-empirical-functional}
取 \(N\) 条独立或满足一致大数定律的弱相关样本轨迹
\[
  X_t^{\Gamma,1},\ldots,X_t^{\Gamma,N}.
\]
对应首次命中时间记为
\[
  \tau_S^{\Gamma,i}
  =
  \inf\{t\ge1:\mathbf 1_S(X_t^{\Gamma,i})=1\}.
\]
定义经验未命中率
\[
  \widehat Q_{T,N}(\Gamma)
  =
  \frac1N
  \sum_{i=1}^{N}
  \mathbf 1\{\tau_S^{\Gamma,i}>T\}.
\]
定义经验命中泛函
\[
  \widehat{\mathcal J}_{T,N}(\Gamma)
  =
  1-\widehat Q_{T,N}(\Gamma).
\]
定义经验命中指数
\[
  \widehat{\mathcal C}_{T,N}(\Gamma)
  =
  -\log\left(\widehat Q_{T,N}(\Gamma)+\varepsilon_N\right),
\]
其中 \(\varepsilon_N>0\) 用于避免 \(\log 0\)，并要求
\[
  \varepsilon_N\to0.
\]
当 \(N\) 很大且一致收敛条件成立时，经验对象逼近理论对象。
\end{definition}

\begin{definition}[有限 \(L\) 的高阶正交截断]
\label{def:tex55-engdet-L-truncation}
令 \(\mathfrak T\) 为完整同伦轨迹函数空间，令
\[
  \mathfrak T_L\subset\mathfrak T
\]
为有限 \(L\) 阶截断空间。
这里 \(L\) 可以表示高阶正交层数、频域模态数，或二者的统一截断尺度。
第 \(k\) 阶已用方向空间为
\[
  \mathcal W_{k-1,N,L}
  =
  \operatorname{span}\{
  u_{0,N,L},
  \delta_{1,N,L},
  \ldots,
  \delta_{k-1,N,L}
  \}.
\]
经验命中梯度为
\[
  \widehat g_{N,L}
  =
  \nabla
  \widehat{\mathcal C}_{T,N,L}
  (\Gamma^{[k-1]}).
\]
第 \(k\) 阶高阶正交命中残差方向为
\[
  \widehat b_{k,N,L}
  =
  \Pi_{\mathcal W_{k-1,N,L}^{\perp}}
  \widehat g_{N,L}.
\]
\end{definition}

\begin{definition}[有限 \((N,L)\) 的经验雅可比间隙]
\label{def:tex55-engdet-empirical-jacobian-gap}
第 \(k\) 阶同伦扭曲由参数
\[
  z_k=(\theta_k,d_k)
\]
控制，其雅可比实际方向为
\[
  J_{k,N,L}(\theta_k)d_k.
\]
定义经验雅可比间隙为
\[
  \boxed{
  \begin{aligned}
  \widehat{\mathfrak G}_{k,N,L}(\theta_k,d_k)
  &=
  \frac12
  \left\lVert
  \Pi_{\mathcal W_{k-1,N,L}^{\perp}}
  J_{k,N,L}(\theta_k)d_k
  -
  \widehat b_{k,N,L}
  \right\rVert^2\\
  &\quad+
  \frac{\rho_k}{2}
  \left\lVert
  \Pi_{\mathcal W_{k-1,N,L}}
  J_{k,N,L}(\theta_k)d_k
  \right\rVert^2.
  \end{aligned}
  }
\]
第一项表示实际雅可比方向与理想高阶命中方向的差距；
第二项表示实际方向泄漏到已用方向空间的正交破坏误差。
因此
\[
  \boxed{
  \widehat{\mathfrak G}_{k,N,L}
  =
  \text{方向逼近误差}
  +
  \text{正交泄漏误差}.
  }
\]
\end{definition}

\subsection{工程截断确定性收敛}
\label{subsec:tex55-engdet-engineering}

\begin{proposition}[有限 \((N,L)\) 下得到工程截断确定性问题]
\label{prop:tex55-engdet-finite-deterministic}
固定样本集、固定截断空间 \(\mathfrak T_L\)、固定数值精度后，问题
\[
  \min_{\theta_k,d_k}
  \widehat{\mathfrak G}_{k,N,L}(\theta_k,d_k)
\]
是一个有限维确定性优化问题。
因此更新
\[
  z_{k,s+1}
  =
  z_{k,s}
  -
  \alpha_{k,s}
  \nabla\widehat{\mathfrak G}_{k,N,L}(z_{k,s})
\]
在给定样本条件下是确定性的。
\end{proposition}

\begin{Certificate}[条件确定性证书]
\label{cert:tex55-engdet-finite-deterministic}
给定样本后，
\[
  \widehat{\mathcal C}_{T,N,L},
  \quad
  \widehat g_{N,L},
  \quad
  \widehat b_{k,N,L},
  \quad
  J_{k,N,L},
  \quad
  \widehat{\mathfrak G}_{k,N,L}
\]
全部成为固定函数。
因此
\[
  z_{k,s}\mapsto z_{k,s+1}
\]
由确定函数给出，不再含新的随机采样。
\end{Certificate}

\begin{proof}
固定样本集后，经验未命中率 \(\widehat Q_{T,N}\) 是样本上的确定函数。
因此经验命中指数
\[
  \widehat{\mathcal C}_{T,N,L}
\]
也是确定函数。
由确定函数求梯度得到 \(\widehat g_{N,L}\)。
由确定的已用方向空间进行投影，得到 \(\widehat b_{k,N,L}\)。
同伦扭曲族在截断空间 \(\mathfrak T_L\) 中的雅可比 \(J_{k,N,L}\) 也是确定算子。
所以 \(\widehat{\mathfrak G}_{k,N,L}\) 是有限维确定函数，梯度下降映射
\[
  z_{k,s}\mapsto z_{k,s}-\alpha_{k,s}
  \nabla\widehat{\mathfrak G}_{k,N,L}(z_{k,s})
\]
为确定映射。
证毕。
\end{proof}

\begin{theorem}[工程截断收敛到经验雅可比间隙的驻点]
\label{thm:tex55-engdet-empirical-stationary}
若 \(\widehat{\mathfrak G}_{k,N,L}\) 是 \(M_{k,N,L}\)-光滑函数，且步长满足
\[
  0<\alpha_{k,s}\le\frac1{M_{k,N,L}},
\]
则
\[
  \widehat{\mathfrak G}_{k,N,L}(z_{k,s+1})
  \le
  \widehat{\mathfrak G}_{k,N,L}(z_{k,s})
  -
  \frac{\alpha_{k,s}}{2}
  \left\lVert
  \nabla\widehat{\mathfrak G}_{k,N,L}(z_{k,s})
  \right\rVert^2.
\]
因此 \(\widehat{\mathfrak G}_{k,N,L}(z_{k,s})\) 单调不增。
若非凸且没有额外结构，一般只能保证驻点意义上的收敛；
若另有凸性、强凸性、PL 条件或唯一稳定极小点条件，则可得到更强的最小点收敛结论。
\end{theorem}

\begin{Certificate}[数学归纳法证书：累计下降界]
\label{cert:tex55-engdet-empirical-stationary}
目标谓词为
\[
  P(S):
  \widehat{\mathfrak G}_{k,N,L}(z_{k,S})
  \le
  \widehat{\mathfrak G}_{k,N,L}(z_{k,0})
  -
  \sum_{s=0}^{S-1}
  \frac{\alpha_{k,s}}2
  \left\lVert
  \nabla\widehat{\mathfrak G}_{k,N,L}(z_{k,s})
  \right\rVert^2.
\]

基例 \(S=1\) 由单步下降不等式直接成立。

归纳步：若 \(P(S)\) 成立，并且单步下降给出
\[
  \widehat{\mathfrak G}(z_{k,S+1})
  \le
  \widehat{\mathfrak G}(z_{k,S})
  -
  \frac{\alpha_{k,S}}2
  \lVert\nabla\widehat{\mathfrak G}(z_{k,S})\rVert^2,
\]
则代入 \(P(S)\) 得到 \(P(S+1)\)。
因此由数学归纳法，累计下降界对所有 \(S\ge1\) 成立。
\end{Certificate}

\begin{proof}
由 \(M_{k,N,L}\)-光滑性，对
\[
  z_{k,s+1}
  =
  z_{k,s}
  -
  \alpha_{k,s}
  \nabla\widehat{\mathfrak G}_{k,N,L}(z_{k,s})
\]
应用下降引理，得到
\[
  \widehat{\mathfrak G}_{k,N,L}(z_{k,s+1})
  \le
  \widehat{\mathfrak G}_{k,N,L}(z_{k,s})
  -
  \alpha_{k,s}
  \left(
  1-\frac{M_{k,N,L}\alpha_{k,s}}2
  \right)
  \left\lVert
  \nabla\widehat{\mathfrak G}_{k,N,L}(z_{k,s})
  \right\rVert^2.
\]
当
\[
  0<\alpha_{k,s}\le\frac1{M_{k,N,L}}
\]
时，
\[
  1-\frac{M_{k,N,L}\alpha_{k,s}}2\ge\frac12.
\]
故得到单步下降式。
证书 \ref{cert:tex55-engdet-empirical-stationary} 给出累计下降的数学归纳法证明。
若 \(\widehat{\mathfrak G}_{k,N,L}\) 下有界，则累计下降界推出
\[
  \sum_s
  \alpha_{k,s}
  \left\lVert
  \nabla\widehat{\mathfrak G}_{k,N,L}(z_{k,s})
  \right\rVert^2
  <\infty.
\]
在步长不退化的常见条件下，可得到梯度范数趋于零的驻点收敛结论。
若再加凸性、强凸性、PL 条件或唯一稳定极小点条件，则可提升为最小点收敛。
证毕。
\end{proof}

\begin{remark}[有限工程截断的严格口径]
有限 \((N,L)\) 下的确定性收敛应写成
\[
  \boxed{
  \begin{gathered}
  \text{经验雅可比间隙下降到驻点，}\\
  \text{或在附加凸性/稳定性条件下下降到经验最小点。}
  \end{gathered}
  }
\]
它不是无条件地收敛到完整理论问题的全局最优解。
\end{remark}

\subsection{样本数极限下随机性消失的严格含义}
\label{subsec:tex55-engdet-N-limit}

\begin{theorem}[经验命中泛函收敛到理论命中泛函]
\label{thm:tex55-engdet-empirical-to-theory}
若对截断轨迹函数类 \(\mathfrak T_L\) 存在一致大数定律：
\[
  \sup_{\Gamma\in\mathfrak T_L}
  \left|
  \widehat{\mathcal C}_{T,N,L}(\Gamma)
  -
  \mathcal C_{T,L}(\Gamma)
  \right|
  \to0
\]
以概率收敛或几乎必然收敛，并且有梯度一致收敛：
\[
  \sup_{\Gamma\in\mathfrak T_L}
  \left\lVert
  \nabla\widehat{\mathcal C}_{T,N,L}(\Gamma)
  -
  \nabla\mathcal C_{T,L}(\Gamma)
  \right\rVert
  \to0,
\]
则经验命中梯度、经验高阶正交残差方向与经验雅可比间隙在投影稳定、雅可比稳定条件下，
收敛到对应的截断理论对象。
\end{theorem}

\begin{Certificate}[一致收敛证书]
\label{cert:tex55-engdet-empirical-to-theory}
证明链为
\[
\begin{aligned}
  &\widehat{\mathcal C}_{T,N,L}\to\mathcal C_{T,L}
  \text{ 一致}\\
  &\Longrightarrow
  \nabla\widehat{\mathcal C}_{T,N,L}
  \to
  \nabla\mathcal C_{T,L}
  \text{ 一致}\\
  &\Longrightarrow
  \widehat g_{N,L}\to g_L\\
  &\Longrightarrow
  \Pi_{\mathcal W_{k-1,N,L}^{\perp}}\widehat g_{N,L}
  \to
  \Pi_{\mathcal W_{k-1,L}^{\perp}}g_L\\
  &\Longrightarrow
  \widehat b_{k,N,L}\to b_{k,L}\\
  &\Longrightarrow
  J_{k,N,L}\to J_{k,L}\\
  &\Longrightarrow
  \widehat{\mathfrak G}_{k,N,L}\to\mathfrak G_{k,L}.
\end{aligned}
\]
\end{Certificate}

\begin{proof}
由一致大数定律，
\[
  \widehat{\mathcal C}_{T,N,L}
  \to
  \mathcal C_{T,L}
\]
在 \(\mathfrak T_L\) 上一致成立。
若梯度也一致收敛，则
\[
  \widehat g_{N,L}
  =
  \nabla\widehat{\mathcal C}_{T,N,L}
  \to
  \nabla\mathcal C_{T,L}
  =
  g_L.
\]
若高阶正交投影稳定，即
\[
  \Pi_{\mathcal W_{k-1,N,L}^{\perp}}
  \to
  \Pi_{\mathcal W_{k-1,L}^{\perp}},
\]
则
\[
  \widehat b_{k,N,L}
  =
  \Pi_{\mathcal W_{k-1,N,L}^{\perp}}
  \widehat g_{N,L}
  \to
  \Pi_{\mathcal W_{k-1,L}^{\perp}}g_L
  =
  b_{k,L}.
\]
若同时
\[
  J_{k,N,L}\to J_{k,L}
\]
一致成立，则经验雅可比间隙的两个组成项分别收敛到截断理论雅可比间隙的对应项。
因此
\[
  \widehat{\mathfrak G}_{k,N,L}
  \to
  \mathfrak G_{k,L}.
\]
这说明有限样本随机性在 \(N\to\infty\) 下被经验平均消除，留下确定性的截断理论问题。
证毕。
\end{proof}

\subsection{截断阶数与一致极限下的严格确定性收敛}
\label{subsec:tex55-engdet-L-limit}

\begin{definition}[截断极限与一致极限]
\label{def:tex55-engdet-L-infty}
若 \(L\) 表示截断阶数或模态数量，则严格极限写作
\[
  L\to\infty.
\]
若 \(L^\infty\) 表示范数意义，则要求一致收敛：
\[
  \left\lVert
  \mathcal C_{T,L}-\mathcal C_T
  \right\rVert_{L^\infty(\mathfrak T)}
  =
  \sup_{\Gamma\in\mathfrak T}
  \left|
  \mathcal C_{T,L}(\Gamma)-\mathcal C_T(\Gamma)
  \right|
  \to0.
\]
进一步要求
\[
  \left\lVert
  \nabla\mathcal C_{T,L}-\nabla\mathcal C_T
  \right\rVert_{L^\infty}
  \to0,
\]
以及
\[
  \left\lVert
  \mathfrak G_{k,L}-\mathfrak G_k
  \right\rVert_{L^\infty}
  \to0.
\]
这统一表达了有限阶高阶正交与有限维雅可比间隙逼近无限维对象的条件。
\end{definition}

\begin{theorem}[\(N\to\infty,L\to\infty\) 下的确定性极限收敛]
\label{thm:tex55-engdet-deterministic-limit}
设满足以下条件：
\[
  \left\lVert
  \widehat{\mathcal C}_{T,N,L}
  -
  \mathcal C_{T,L}
  \right\rVert_{L^\infty}
  \to0
  \quad(N\to\infty);
\]
\[
  \left\lVert
  \mathcal C_{T,L}
  -
  \mathcal C_T
  \right\rVert_{L^\infty}
  \to0
  \quad(L\to\infty);
\]
\[
  \nabla\widehat{\mathcal C}_{T,N,L}
  \to
  \nabla\mathcal C_T
  \quad\text{一致成立};
\]
\[
  J_{k,N,L}\to J_k
  \quad\text{一致成立};
\]
\[
  \Pi_{\mathcal W_{k-1,N,L}^{\perp}}
  \to
  \Pi_{\mathcal W_{k-1}^{\perp}}
  \quad\text{稳定成立};
\]
并且极限雅可比间隙问题有稳定极小点，或满足 PL、强凸、孤立稳定极小点等稳定性条件。
则有限截断系统产生的轨迹函数
\[
  \Gamma_{N,L}
\]
收敛到极限确定性轨迹
\[
  \Gamma_*.
\]
并且
\[
  \widehat{\mathcal J}_{T,N,L}(\Gamma_{N,L})
  \to
  \mathcal J_T(\Gamma_*).
\]
\end{theorem}

\begin{Certificate}[三角不等式与稳定极小点证书]
\label{cert:tex55-engdet-deterministic-limit}
总误差由两段组成：
\[
  \widehat{\mathcal C}_{T,N,L}-\mathcal C_T
  =
  \left(
  \widehat{\mathcal C}_{T,N,L}-\mathcal C_{T,L}
  \right)
  +
  \left(
  \mathcal C_{T,L}-\mathcal C_T
  \right).
\]
因此
\[
  \lVert
  \widehat{\mathcal C}_{T,N,L}-\mathcal C_T
  \rVert_{L^\infty}
  \le
  \lVert
  \widehat{\mathcal C}_{T,N,L}-\mathcal C_{T,L}
  \rVert_{L^\infty}
  +
  \lVert
  \mathcal C_{T,L}-\mathcal C_T
  \rVert_{L^\infty}.
\]
若两项都趋于零，则经验截断命中指数一致收敛到极限命中指数。
\end{Certificate}

\begin{proof}
由证书中的三角不等式，
\[
  \left\lVert
  \widehat{\mathcal C}_{T,N,L}
  -
  \mathcal C_T
  \right\rVert_{L^\infty}
  \to0.
\]
由梯度一致收敛、雅可比一致收敛与投影稳定收敛，可得
\[
  \widehat{\mathfrak G}_{k,N,L}
  \to
  \mathfrak G_k
\]
在相应参数域上一致成立。
由极小点稳定性条件，有限问题的最优点或稳定下降极限点收敛到极限问题的对应点：
\[
  \Gamma_{N,L}\to\Gamma_*.
\]
又因为
\[
  \mathcal J_T(\Gamma)=1-e^{-\mathcal C_T(\Gamma)}
\]
是 \(\mathcal C_T\) 的连续函数，所以
\[
  \widehat{\mathcal J}_{T,N,L}(\Gamma_{N,L})
  \to
  \mathcal J_T(\Gamma_*).
\]
证毕。
\end{proof}

\subsection{随机性到确定性的因果链与误差分解}
\label{subsec:tex55-engdet-chain-error}

\begin{proposition}[随机性到确定性的转换链]
\label{prop:tex55-engdet-random-to-deterministic-chain}
统计力学同伦扭曲命中的随机过程，可通过如下链条转化为确定性极限问题：
\[
\begin{aligned}
  &\text{统计力学同伦扭曲命中}\\
  &\Longrightarrow
  \text{随机变量簇 }\{X_t^\Gamma\}\\
  &\Longrightarrow
  \text{命中指示变量簇 }Y_t^\Gamma=\mathbf 1_S(X_t^\Gamma)\\
  &\Longrightarrow
  \text{命中泛函 }\mathcal J_T(\Gamma)\\
  &\Longrightarrow
  \text{经验命中泛函 }\widehat{\mathcal J}_{T,N,L}(\Gamma)\\
  &\Longrightarrow
  \text{有限维高阶正交命中残差 }\widehat b_{k,N,L}\\
  &\Longrightarrow
  \text{有限维雅可比间隙 }\widehat{\mathfrak G}_{k,N,L}\\
  &\Longrightarrow
  \text{雅可比间隙梯度下降}\\
  &\Longrightarrow
  \text{工程截断确定性轨迹 }\Gamma_{N,L}\\
  &\Longrightarrow
  N\to\infty,\ L\to\infty\\
  &\Longrightarrow
  \text{数学严格确定性极限轨迹 }\Gamma_*.
\end{aligned}
\]
因此可概括为
\[
  \boxed{
  \text{随机命中}
  \to
  \text{经验确定性}
  \to
  \text{截断确定性}
  \to
  \text{极限确定性}.
  }
\]
\end{proposition}

\begin{proof}
由定义 \ref{def:tex55-engdet-hit-process}，统计力学同伦扭曲命中首先表现为随机变量簇。
由命中指示变量得到命中泛函。
由定义 \ref{def:tex55-engdet-empirical-functional}，有限样本把命中泛函替换为经验命中泛函。
由定义 \ref{def:tex55-engdet-L-truncation}，有限 \(L\) 把完整轨迹空间替换为截断空间。
由定义 \ref{def:tex55-engdet-empirical-jacobian-gap}，每阶高阶正交命中方向被转化为经验雅可比间隙。
由命题 \ref{prop:tex55-engdet-finite-deterministic}，固定样本与截断后得到确定性下降系统。
由定理 \ref{thm:tex55-engdet-deterministic-limit}，在一致收敛与稳定性条件下，工程截断轨迹收敛到数学严格的确定性极限轨迹。
证毕。
\end{proof}

\begin{proposition}[总误差分解]
\label{prop:tex55-engdet-total-error}
有限工程截断与数学极限之间的误差可分解为
\[
  \boxed{
  \varepsilon_{\mathrm{total}}
  \le
  \varepsilon_N
  +
  \varepsilon_L
  +
  \varepsilon_{\mathrm{gap}}
  +
  \varepsilon_{\mathrm{orth}}
  +
  \varepsilon_{\mathrm{opt}}.
  }
\]
其中 \(\varepsilon_N\) 是采样误差，
\(\varepsilon_L\) 是截断误差，
\(\varepsilon_{\mathrm{gap}}\) 是雅可比间隙逼近误差，
\(\varepsilon_{\mathrm{orth}}\) 是高阶正交泄漏误差，
\(\varepsilon_{\mathrm{opt}}\) 是数值优化误差。
\end{proposition}

\begin{proof}
把工程系统到极限系统的偏差沿对象链拆开：
第一段是经验平均到理论泛函的偏差，记为 \(\varepsilon_N\)。
第二段是有限 \(L\) 截断到完整空间的偏差，记为 \(\varepsilon_L\)。
第三段是实际雅可比方向与理想命中残差方向的偏差，记为 \(\varepsilon_{\mathrm{gap}}\)。
第四段是实际方向泄漏到已用空间造成的正交破坏，记为 \(\varepsilon_{\mathrm{orth}}\)。
第五段是梯度下降未完全求解经验间隙问题造成的优化残差，记为 \(\varepsilon_{\mathrm{opt}}\)。
由三角不等式，总误差不超过这些误差分量之和。
证毕。
\end{proof}

\subsection{高阶正交与雅可比间隙的作用}
\label{subsec:tex55-engdet-role}

\begin{remark}[高阶正交的作用]
高阶正交负责不重复已用命中加速方向，并持续提取剩余命中增益。
第 \(k\) 阶方向为
\[
  b_k
  =
  \Pi_{\mathcal W_{k-1}^{\perp}}
  \nabla\mathcal C_T(\Gamma^{[k-1]}),
\]
其中
\[
  \mathcal W_{k-1}
  =
  \operatorname{span}\{u_0,\delta_1,\ldots,\delta_{k-1}\}.
\]
若 \(b_k\neq0\)，说明仍有剩余可用命中增益。
若 \(b_k=0\)，说明当前正交补空间中没有一阶命中加速方向。
\end{remark}

\begin{remark}[雅可比间隙的作用]
雅可比间隙负责把理想高阶命中方向 \(b_k\) 转化为实际参数可实现方向 \(J_kd_k\)。
即
\[
  J_kd_k\approx b_k.
\]
如果
\[
  \lVert J_kd_k-b_k\rVert
\]
足够小，则该阶同伦扭曲能提高命中指数。
如果该间隙太大，则虽然理想方向存在，实际参数系统仍可能无法实现它。
所以雅可比间隙是
\[
  \boxed{
  \text{工程可实现性与理论命中方向之间的桥。}
  }
\]
\end{remark}

\subsection{边界与最终严格表述}
\label{subsec:tex55-engdet-boundary-final}

\begin{remark}[必须保留的边界]
第一，单个样本不会变成确定性。
即使 \(N\to\infty\)，单个
\[
  X_t^\Gamma
\]
仍然是随机变量。
确定的是经验测度、经验命中泛函、经验命中指数、经验雅可比间隙与轨迹极限等对象。

第二，有限 \((N,L)\) 是工程确定性，不是数学精确性。
有限 \((N,L)\) 下通常只有
\[
  \widehat{\mathcal J}_{T,N,L}\approx\mathcal J_T,
\]
误差由采样误差与截断误差控制。

第三，非凸雅可比间隙只能无条件保证驻点收敛。
若 \(\widehat{\mathfrak G}_{k,N,L}\) 非凸，则梯度下降一般只能保证
\[
  \nabla\widehat{\mathfrak G}_{k,N,L}\to0
\]
意义下的驻点收敛。
要保证全局最小，需要额外结构。

第四，幂化加速仍需幂压缩条件。
命中指数提高
\[
  \mathcal C_T(\Gamma^{[k]})
  >
  \mathcal C_T(\Gamma^{[k-1]})
\]
只说明命中加速。
若要得到
\[
  Q_T(\Gamma^{[k]})
  \le
  Q_T(\Gamma^{[k-1]})^{A_k},
  \qquad A_k>1,
\]
还需要额外的幂压缩条件。
\end{remark}

\begin{remark}[最终严格表述]
用户原句可严格改写为：
\[
  \boxed{
  \begin{gathered}
  \text{上述框架等价于：对统计力学同伦扭曲命中的随机过程，}\\
  \text{构造由 }N\text{ 个样本与 }L\text{ 阶高阶正交截断形成的经验确定性系统；}\\
  \text{在该系统中，每阶通过雅可比间隙梯度下降逼近理想高阶正交命中方向，}\\
  \text{从而得到工程可执行的确定性收敛轨迹。}
  \end{gathered}
  }
\]
进一步，
\[
  \boxed{
  \begin{gathered}
  \text{当 }N\to\infty,\ L\to\infty\text{，并且经验泛函一致收敛、}\\
  \text{截断误差消失、投影稳定、雅可比稳定、下降问题稳定时，}\\
  \text{该工程确定性轨迹收敛到无限维统计力学同伦命中问题的数学严格确定性极限。}
  \end{gathered}
  }
\]
最简表达为
\[
  \boxed{
  \begin{gathered}
  \text{有限 }N,L\text{ 下是工程截断确定性收敛；}\\
  N\to\infty,\quad L\to\infty
  \text{ 且满足一致收敛与稳定性条件时，}\\
  \text{是数学严格确定性收敛。}
  \end{gathered}
  }
\]
对应极限公式为
\[
  \boxed{
  \Gamma_{N,L}\to\Gamma_*,
  \qquad
  \widehat{\mathcal J}_{T,N,L}(\Gamma_{N,L})
  \to
  \mathcal J_T(\Gamma_*).
  }
\]
\end{remark}

\clearpage

%%%%%%%%%%%%%%%%%%%%%%%%%%%%%%%%%%%%%%%%%%%%%%%%%%%%%%%%%%%%
\section{形式逻辑：统计力学同伦命中对高阶偏微分方程求解范式的替代}
\label{sec:tex55-pde-paradigm-replacement}
%%%%%%%%%%%%%%%%%%%%%%%%%%%%%%%%%%%%%%%%%%%%%%%%%%%%%%%%%%%%

\begin{remark}[严格替代口径]
本章把“统计力学同伦扭曲命中对高阶 PDE 具有决定意义”表述为条件命题。
严格替代口径是：
\[
  \boxed{
  \begin{gathered}
  \text{保留高阶 PDE 的残差语义，}\\
  \text{并把“解高阶 PDE”转化为}\\
  \text{“在离散统计力学同伦扭曲系统中命中低残差解集”。}
  \end{gathered}
  }
\]
所以所谓“替代”更准确地说是：
\[
  \boxed{
  \text{用“离散残差能量命中 + 同伦扭曲轨迹收敛”替代“直接求解高阶微分算子方程”。}
  }
\]
高阶 PDE 的传统困难来自高阶导数不稳定、复杂边界条件、非线性强、局部极小多、
离散刚性强与求解器初值敏感。
本章把这些困难转化为残差能量函数、统计力学分布律、命中随机变量簇、同伦扭曲轨迹、
高阶正交修正与雅可比间隙下降。
偏微分方程和弱解背景可参见 \cite{Evans2010PDE}，
离散残差、有限元近似与相容性背景可参见 \cite{BrennerScott2008FEM}。
\end{remark}

\subsection{高阶偏微分方程的替代对象}
\label{subsec:tex55-pde-replacement-object}

\begin{definition}[高阶 PDE 与边界条件]
\label{def:tex55-pde-high-order}
设高阶偏微分方程为
\[
  \mathcal P[u]=0,
  \qquad
  x\in D,
\]
并带边界条件
\[
  \mathcal B[u]=0,
  \qquad
  x\in\partial D.
\]
其中
\[
  \mathcal P[u]
  =
  \mathcal P(x,u,\nabla u,\nabla^2u,\ldots,\nabla^m u),
  \qquad
  m\ge2.
\]
当 \(m\) 很高或 \(\mathcal P\) 强非线性时，直接求解
\[
  \mathcal P[u]=0
\]
通常会遇到高阶导数离散不稳定、边界条件耦合复杂、非线性残差刚性等问题。
\end{definition}

\begin{definition}[离散状态、重构与残差能量]
\label{def:tex55-pde-residual-energy}
令 \(\Xi_L\) 为第 \(L\) 阶离散状态空间。
元素
\[
  z\in\Xi_L
\]
表示网格值、谱系数、有限元自由度或其他离散状态。
令
\[
  R_L:\Xi_L\to\mathcal U_L
\]
为重构算子，\(R_Lz\) 是由离散状态 \(z\) 重构出的函数近似。
令
\[
  \mathcal P_L[R_Lz]
\]
为 PDE 算子的离散残差，令
\[
  \mathcal B_L[R_Lz]
\]
为边界残差。
定义残差能量
\[
  E_L(z)
  =
  \lVert\mathcal P_L[R_Lz]\rVert^2
  +
  \alpha
  \lVert\mathcal B_L[R_Lz]\rVert^2
  +
  \lambda
  \mathcal R_L(z),
\]
其中 \(\alpha,\lambda\ge0\)，\(\mathcal R_L\) 是稳定化或正则化项。
\end{definition}

\begin{definition}[低残差命中集合]
\label{def:tex55-pde-low-residual-hit-set}
给定容忍度 \(\varepsilon>0\)，定义离散低残差命中集合为
\[
  S_{\varepsilon,L}
  =
  \{z\in\Xi_L:E_L(z)\le\varepsilon\}.
\]
于是直接解 PDE 的目标
\[
  \mathcal P[u]=0,
  \qquad
  \mathcal B[u]=0
\]
被替代为命中条件
\[
  z\in S_{\varepsilon,L}.
\]
这就是范式替代的第一层：
\[
  \boxed{
  \text{高阶 PDE 求解}
  \Longrightarrow
  \text{离散残差低能集合命中。}
  }
\]
\end{definition}

\begin{proposition}[低残差命中推出离散残差受控]
\label{prop:tex55-pde-low-residual-controlled}
若
\[
  z_L\in S_{\varepsilon,L},
\]
且 \(\alpha>0\)、\(\lambda\ge0\)、\(\mathcal R_L(z_L)\ge0\)，则
\[
  \lVert\mathcal P_L[R_Lz_L]\rVert
  \le
  \varepsilon^{1/2},
\]
并且
\[
  \lVert\mathcal B_L[R_Lz_L]\rVert
  \le
  \alpha^{-1/2}\varepsilon^{1/2}.
\]
\end{proposition}

\begin{Certificate}[符号推演证书]
\label{cert:tex55-pde-low-residual-controlled}
证明链为
\[
\begin{aligned}
  &z_L\in S_{\varepsilon,L}\\
  &\Longrightarrow
  E_L(z_L)\le\varepsilon\\
  &\Longrightarrow
  \lVert\mathcal P_L[R_Lz_L]\rVert^2
  +
  \alpha\lVert\mathcal B_L[R_Lz_L]\rVert^2
  \le\varepsilon\\
  &\Longrightarrow
  \lVert\mathcal P_L[R_Lz_L]\rVert\le\varepsilon^{1/2},
  \quad
  \lVert\mathcal B_L[R_Lz_L]\rVert\le\alpha^{-1/2}\varepsilon^{1/2}.
\end{aligned}
\]
\end{Certificate}

\begin{proof}
由 \(z_L\in S_{\varepsilon,L}\) 的定义，
\[
  E_L(z_L)\le\varepsilon.
\]
由残差能量定义，
\[
  \lVert\mathcal P_L[R_Lz_L]\rVert^2
  +
  \alpha
  \lVert\mathcal B_L[R_Lz_L]\rVert^2
  +
  \lambda\mathcal R_L(z_L)
  \le
  \varepsilon.
\]
由于 \(\lambda\mathcal R_L(z_L)\ge0\)，得到
\[
  \lVert\mathcal P_L[R_Lz_L]\rVert^2
  \le
  \varepsilon,
\]
以及
\[
  \alpha\lVert\mathcal B_L[R_Lz_L]\rVert^2
  \le
  \varepsilon.
\]
取平方根即可。
证毕。
\end{proof}

\subsection{统计力学同伦扭曲命中形式}
\label{subsec:tex55-pde-statistical-hit}

\begin{definition}[离散统计力学分布律]
\label{def:tex55-pde-statistical-law}
在离散状态空间 \(\Xi_L\) 上定义统计力学分布
\[
  \pi_{\lambda,L}(z)
  =
  \frac{
  w_L(z)
  \exp\left(
  -\beta E_L(z)
  +
  \sum_{j=1}^{K_L}\theta_j\phi_{j,L}(z)
  \right)
  }{
  Z_{\lambda,L}
  }.
\]
其中
\[
  \lambda=(\beta,w_L,K_L,\theta_1,\ldots,\theta_{K_L})
\]
是分布参数，\(\beta\) 控制低残差集中，
\(\phi_{j,L}\) 是频域模态或结构模态，\(\theta_j\) 是同伦扭曲参数。
\end{definition}

\begin{definition}[同伦扭曲命中随机过程]
\label{def:tex55-pde-homotopy-hit-process}
令同伦扭曲轨迹函数为 \(\Gamma\)，并令
\[
  \lambda_t=\Gamma_t(\mathcal F_{t-1}).
\]
定义随机变量
\[
  X_t^{\Gamma,L}\sim\pi_{\lambda_t,L}.
\]
命中指示变量为
\[
  Y_t^{\Gamma,L}
  =
  \mathbf 1_{S_{\varepsilon,L}}(X_t^{\Gamma,L}).
\]
首次命中时间为
\[
  \tau_{\varepsilon,L}^{\Gamma}
  =
  \inf\{t\ge1:Y_t^{\Gamma,L}=1\}.
\]
命中泛函、生存概率与命中指数分别为
\[
  \mathcal J_{T,L}(\Gamma)
  =
  \mathbb P_\Gamma(\tau_{\varepsilon,L}^{\Gamma}\le T),
\]
\[
  Q_{T,L}(\Gamma)
  =
  \mathbb P_\Gamma(\tau_{\varepsilon,L}^{\Gamma}>T),
\]
\[
  \mathcal C_{T,L}(\Gamma)
  =
  -\log Q_{T,L}(\Gamma).
\]
于是高阶 PDE 的离散求解目标被转化为
\[
  \boxed{
  \max_\Gamma\mathcal J_{T,L}(\Gamma)
  }
\]
或等价地
\[
  \boxed{
  \max_\Gamma\mathcal C_{T,L}(\Gamma).
  }
\]
\end{definition}

\begin{theorem}[低残差命中的幂化概率提升]
\label{thm:tex55-pde-power-hit}
若高阶正交同伦扭曲的每一阶满足
\[
  Q_{T,L}(\Gamma^{[k]})
  \le
  Q_{T,L}(\Gamma^{[k-1]})^{A_k},
  \qquad
  A_k>1,
\]
则经过 \(n\) 阶后
\[
  Q_{T,L}(\Gamma^{[n]})
  \le
  Q_{T,L}(\Gamma^{[0]})^{A_1A_2\cdots A_n},
\]
并且
\[
  \mathcal J_{T,L}(\Gamma^{[n]})
  \ge
  1-
  Q_{T,L}(\Gamma^{[0]})^{A_1A_2\cdots A_n}.
\]
\end{theorem}

\begin{Certificate}[数学归纳法证书]
\label{cert:tex55-pde-power-hit}
目标谓词为
\[
  P(n):
  Q_{T,L}(\Gamma^{[n]})
  \le
  Q_{T,L}(\Gamma^{[0]})^{\prod_{k=1}^{n}A_k}.
\]
基例 \(n=1\) 由假设直接成立。
归纳步：若 \(P(n)\) 成立，则
\[
  Q_{T,L}(\Gamma^{[n+1]})
  \le
  Q_{T,L}(\Gamma^{[n]})^{A_{n+1}}
  \le
  \left(
  Q_{T,L}(\Gamma^{[0]})^{\prod_{k=1}^{n}A_k}
  \right)^{A_{n+1}},
\]
即 \(P(n+1)\) 成立。
\end{Certificate}

\begin{proof}
由证书 \ref{cert:tex55-pde-power-hit} 得到
\[
  Q_{T,L}(\Gamma^{[n]})
  \le
  Q_{T,L}(\Gamma^{[0]})^{A_1A_2\cdots A_n}.
\]
又因为
\[
  \mathcal J_{T,L}(\Gamma)=1-Q_{T,L}(\Gamma),
\]
所以
\[
  \mathcal J_{T,L}(\Gamma^{[n]})
  \ge
  1-
  Q_{T,L}(\Gamma^{[0]})^{A_1A_2\cdots A_n}.
\]
证毕。
\end{proof}

\subsection{高阶正交与雅可比间隙的决定作用}
\label{subsec:tex55-pde-orthogonal-jacobian}

\begin{definition}[PDE 低残差命中的高阶正交残差]
\label{def:tex55-pde-orthogonal-residual}
第 \(k\) 阶命中梯度为
\[
  g_k
  =
  \nabla\mathcal C_{T,L}(\Gamma^{[k-1]}).
\]
已用方向空间为
\[
  \mathcal W_{k-1}
  =
  \operatorname{span}\{u_0,\delta_1,\ldots,\delta_{k-1}\}.
\]
第 \(k\) 阶高阶正交命中残差为
\[
  b_k
  =
  \Pi_{\mathcal W_{k-1}^{\perp}}g_k.
\]
若 \(b_k\neq0\)，则当前正交补空间中仍有可用的低残差命中增益方向。
\end{definition}

\begin{definition}[PDE 替代范式中的雅可比间隙]
\label{def:tex55-pde-jacobian-gap}
实际同伦扭曲由参数产生，其雅可比方向为
\[
  J_k(\theta_k)d_k.
\]
定义第 \(k\) 阶雅可比间隙为
\[
  \mathfrak G_k(\theta_k,d_k)
  =
  \frac12
  \left\lVert
  \Pi_{\mathcal W_{k-1}^{\perp}}
  J_k(\theta_k)d_k
  -
  b_k
  \right\rVert^2
  +
  \frac{\rho_k}{2}
  \left\lVert
  \Pi_{\mathcal W_{k-1}}
  J_k(\theta_k)d_k
  \right\rVert^2.
\]
它的意义是
\[
  \boxed{
  \mathfrak G_k
  =
  \text{实际同伦扭曲方向与理想高阶命中方向之间的工程误差。}
  }
\]
通过
\[
  (\theta_k,d_k)
  \leftarrow
  (\theta_k,d_k)
  -
  \alpha_k
  \nabla\mathfrak G_k(\theta_k,d_k)
\]
可以使实际扭曲方向逼近理想命中方向。
\end{definition}

\begin{proposition}[随机命中获得确定性收敛指引]
\label{prop:tex55-pde-deterministic-guide}
若雅可比间隙下降使
\[
  J_k(\theta_k)d_k\approx b_k,
\]
且逼近误差满足上述正增长判据，则第 \(k\) 阶同伦扭曲提高命中指数
\[
  \mathcal C_{T,L}(\Gamma^{[k]})
  >
  \mathcal C_{T,L}(\Gamma^{[k-1]}).
\]
因此统计力学随机命中不再只是概率启发式，而获得每阶同伦扭曲的确定性收敛指引。
\end{proposition}

\begin{proof}
由定义 \ref{def:tex55-pde-orthogonal-residual}，\(b_k\) 是命中指数梯度在剩余正交补空间中的投影。
若 \(J_k(\theta_k)d_k\) 足够接近 \(b_k\)，则由雅可比间隙增长判据可得
\[
  D\mathcal C_{T,L}(\Gamma^{[k-1]})[J_k(\theta_k)d_k]>0.
\]
取足够小正步长，即得
\[
  \mathcal C_{T,L}(\Gamma^{[k]})
  >
  \mathcal C_{T,L}(\Gamma^{[k-1]}).
\]
证毕。
\end{proof}

\subsection{核心替代定理：低残差命中极限恢复 PDE 解}
\label{subsec:tex55-pde-core-replacement-theorem}

\begin{theorem}[公理降级：低残差闭性传递原则]
\label{thm:tex55-pde-residual-closedness}
设存在离散状态序列 \(z_{N,L}\in\Xi_L\)，并令
\[
  u_{N,L}=R_Lz_{N,L}.
\]
若满足：
\[
  E_L(z_{N,L})\le\varepsilon_{N,L},
  \qquad
  \varepsilon_{N,L}\to0;
\]
\[
  u_{N,L}\to u_*
  \quad
  \text{在某个可传递残差的函数空间中成立};
\]
\[
  \mathcal P_L[u_{N,L}]\to0,
  \qquad
  \mathcal B_L[u_{N,L}]\to0;
\]
并且离散残差对连续 PDE 残差相容且闭合，即
\[
  u_{N,L}\to u_*,
  \quad
  \mathcal P_L[u_{N,L}]\to0
  \Longrightarrow
  \mathcal P[u_*]=0,
\]
\[
  u_{N,L}\to u_*,
  \quad
  \mathcal B_L[u_{N,L}]\to0
  \Longrightarrow
  \mathcal B[u_*]=0,
\]
则
\[
  \mathcal P[u_*]=0,
  \qquad
  \mathcal B[u_*]=0.
\]
因此 \(u_*\) 是原高阶 PDE 的弱解、黏性解、能量解或经典解，具体取决于所选残差范数与闭性框架。
\end{theorem}

\begin{Certificate}[谓词因果链证书]
\label{cert:tex55-pde-residual-closedness}
证明链为
\[
\begin{aligned}
  &z_{N,L}\in S_{\varepsilon_{N,L},L}\\
  &\Longrightarrow
  E_L(z_{N,L})\le\varepsilon_{N,L}\\
  &\Longrightarrow
  \lVert\mathcal P_L[R_Lz_{N,L}]\rVert\to0,
  \quad
  \lVert\mathcal B_L[R_Lz_{N,L}]\rVert\to0\\
  &\Longrightarrow
  u_{N,L}=R_Lz_{N,L}\to u_*\\
  &\Longrightarrow
  \text{残差相容与闭性}\\
  &\Longrightarrow
  \mathcal P[u_*]=0,\quad
  \mathcal B[u_*]=0.
\end{aligned}
\]
\end{Certificate}

\begin{proof}
由 \(E_L(z_{N,L})\le\varepsilon_{N,L}\) 与命题
\ref{prop:tex55-pde-low-residual-controlled}，得到
\[
  \lVert\mathcal P_L[R_Lz_{N,L}]\rVert\to0,
\]
并且在 \(\alpha>0\) 时
\[
  \lVert\mathcal B_L[R_Lz_{N,L}]\rVert\to0.
\]
记
\[
  u_{N,L}=R_Lz_{N,L}.
\]
由假设 \(u_{N,L}\to u_*\)，以及残差相容和闭性，
\[
  \mathcal P_L[u_{N,L}]\to0
  \Longrightarrow
  \mathcal P[u_*]=0,
\]
\[
  \mathcal B_L[u_{N,L}]\to0
  \Longrightarrow
  \mathcal B[u_*]=0.
\]
因此
\[
  \mathcal P[u_*]=0,
  \qquad
  \mathcal B[u_*]=0.
\]
证毕。
\end{proof}

\begin{theorem}[统计力学同伦命中替代高阶 PDE 直接求解]
\label{thm:tex55-pde-paradigm-replacement}
若满足以下条件：
\[
  \text{一、残差能量相容：}
  \quad
  E_L(z_L)\to0
  \Longrightarrow
  \mathcal P[R_Lz_L]\to0,\quad
  \mathcal B[R_Lz_L]\to0;
\]
\[
  \text{二、离散重构紧性：}
  \quad
  u_L=R_Lz_L
  \text{ 在目标函数空间中预紧};
\]
\[
  \text{三、PDE 残差闭性：}
  \quad
  u_L\to u_*,
  \ \mathcal P_L[u_L]\to0
  \Longrightarrow
  \mathcal P[u_*]=0;
\]
\[
  \text{四、命中泛函一致收敛：}
  \quad
  \lVert\widehat{\mathcal C}_{T,N,L}-\mathcal C_T\rVert_{L^\infty}\to0;
\]
\[
  \text{五、雅可比间隙稳定：}
  \quad
  \lVert\nabla\widehat{\mathfrak G}_{k,N,L}-\nabla\mathfrak G_k\rVert_{L^\infty}\to0;
\]
\[
  \text{六、高阶正交残差非零：}
  \quad
  b_k\neq0
  \text{ 在需要加速的阶上成立};
\]
并且同伦命中算法输出
\[
  z_{N,L}\in S_{\varepsilon_{N,L},L},
  \qquad
  \varepsilon_{N,L}\to0,
\]
则存在子列使
\[
  R_Lz_{N,L}\to u_*,
\]
且
\[
  \mathcal P[u_*]=0,
  \qquad
  \mathcal B[u_*]=0.
\]
若原高阶 PDE 解唯一，则整个重构序列收敛到该解。
\end{theorem}

\begin{Certificate}[替代定理证书]
\label{cert:tex55-pde-paradigm-replacement}
本定理的核心链条为
\[
\begin{aligned}
  &\text{同伦命中输出 }z_{N,L}\in S_{\varepsilon_{N,L},L}\\
  &\Longrightarrow
  E_L(z_{N,L})\le\varepsilon_{N,L}\to0\\
  &\Longrightarrow
  \mathcal P_L[R_Lz_{N,L}]\to0,\quad
  \mathcal B_L[R_Lz_{N,L}]\to0\\
  &\Longrightarrow
  \text{紧性给出 }R_Lz_{N,L}\to u_*\text{ 的子列}\\
  &\Longrightarrow
  \text{残差相容与闭性}\\
  &\Longrightarrow
  \mathcal P[u_*]=0,\quad
  \mathcal B[u_*]=0.
\end{aligned}
\]
\end{Certificate}

\begin{proof}
由命中输出
\[
  z_{N,L}\in S_{\varepsilon_{N,L},L}
\]
可得
\[
  E_L(z_{N,L})\le\varepsilon_{N,L}\to0.
\]
由残差能量相容与命题 \ref{prop:tex55-pde-low-residual-controlled}，
\[
  \mathcal P_L[R_Lz_{N,L}]\to0,
  \qquad
  \mathcal B_L[R_Lz_{N,L}]\to0.
\]
由离散重构紧性，存在子列仍记作 \(R_Lz_{N,L}\)，使
\[
  R_Lz_{N,L}\to u_*.
\]
由 PDE 残差闭性，得到
\[
  \mathcal P[u_*]=0,
  \qquad
  \mathcal B[u_*]=0.
\]
若解唯一，则任意收敛子列的极限相同，故整个序列收敛到唯一解。
证毕。
\end{proof}

\subsection{决定意义与传统范式对应}
\label{subsec:tex55-pde-significance}

\begin{remark}[替代直接处理高阶导数的中心地位]
传统高阶 PDE 解法通常围绕高阶导数
\[
  \nabla^m u
\]
构造离散算子、弱形式、有限元空间、谱方法或差分格式。
本范式把中心对象换成残差能量
\[
  E_L(z)=\text{PDE 残差能量}.
\]
也就是说，高阶导数不再直接作为求解动力系统的核心，而是被吸收到能量函数中：
\[
  \mathcal P[u]=0
  \quad
  \leadsto
  \quad
  E_L(z)\approx0.
\]
因此
\[
  \boxed{
  \text{高阶 PDE 的微分困难被转化为低残差集合命中困难。}
  }
\]
\end{remark}

\begin{remark}[从求解转化为命中]
传统范式是
\[
  \text{求 }u\text{ 使 }\mathcal P[u]=0.
\]
本范式是
\[
  \text{构造随机变量簇 }X_t^{\Gamma,L}
  \text{，使命中 }S_{\varepsilon,L}.
\]
也就是
\[
  X_t^{\Gamma,L}\in S_{\varepsilon,L}.
\]
这使高阶 PDE 求解从确定性点求解转为
\[
  \boxed{
  \text{命中残差零集或低残差集的随机过程优化。}
  }
\]
随后再通过 \(N\to\infty\)、\(L\to\infty\) 与 \(\varepsilon\to0\) 恢复确定性极限。
\end{remark}

\begin{remark}[从随机搜索升级为确定性极限收敛系统]
如果只有统计力学采样，结论最多是高概率命中。
引入
\[
  \text{高阶正交}
  +
  \text{雅可比间隙梯度下降}
\]
后，形成
\[
  \text{随机命中}
  \to
  \text{工程截断确定性}
  \to
  \text{数学极限确定性}.
\]
即
\[
  \widehat{\mathcal C}_{T,N,L}\to\mathcal C_T,
  \qquad
  \widehat{\mathfrak G}_{k,N,L}\to\mathfrak G_k,
  \qquad
  \Gamma_{N,L}\to\Gamma_*.
\]
这使 PDE 替代范式不只是概率启发式，而可以写成具有收敛证书的确定性极限算法。
\end{remark}

\begin{remark}[绕开部分局部线性化病态]
高阶非线性 PDE 的 Newton 型或局部线性化方法常面对
\[
  D\mathcal P[u]
\]
的病态、奇异、强非线性或边界耦合不稳定。
本范式使用
\[
  \nabla\mathcal C_T(\Gamma),
  \qquad
  b_k=\Pi_{\mathcal W_{k-1}^{\perp}}\nabla\mathcal C_T(\Gamma),
  \qquad
  \mathfrak G_k.
\]
核心下降不直接依赖
\[
  D\mathcal P[u]^{-1},
\]
而依赖命中泛函与同伦扭曲轨迹的雅可比间隙。
因此它把
\[
  \text{PDE 算子反演病态}
\]
转化为
\[
  \text{同伦轨迹可实现方向与理想命中方向的间隙控制}.
\]
\end{remark}

\begin{remark}[传统范式与同伦命中范式的对应]
二者的对应关系可写成：
\begin{center}
\renewcommand{\arraystretch}{1.35}
\begin{tabular}{p{0.40\linewidth}|p{0.50\linewidth}}
\textbf{传统高阶 PDE 范式} & \textbf{统计力学同伦命中范式}\\
\hline
\(\mathcal P[u]=0\) & \(E_L(z)\le\varepsilon\)\\
\(u\) 是解 & \(z\in S_{\varepsilon,L}\)\\
高阶导数稳定性 & 残差能量稳定性\\
Newton/Galerkin/差分求解 & 同伦扭曲命中\\
局部线性化方向 & 命中梯度正交投影方向\\
算子雅可比病态 & 雅可比间隙可控性\\
误差估计 &
\(\varepsilon_N+\varepsilon_L+\varepsilon_{\mathrm{gap}}
+\varepsilon_{\mathrm{orth}}\)\\
解收敛 & \(\Gamma_{N,L}\to\Gamma_*,\ R_Lz_{N,L}\to u_*\)
\end{tabular}
\end{center}
所以
\[
  \boxed{
  \text{它保留 PDE 理论的残差语义，并把 PDE 求解过程重构为命中泛函轨迹优化。}
  }
\]
\end{remark}

\subsection{严格条件、边界与最终表述}
\label{subsec:tex55-pde-boundary-final}

\begin{remark}[决定意义成立的条件]
该范式对高阶 PDE 具有决定意义，需要以下条件：

第一，残差能量相容：
\[
  E_L(z_L)\to0
  \Longrightarrow
  \mathcal P[R_Lz_L]\to0,
  \qquad
  \mathcal B[R_Lz_L]\to0.
\]

第二，离散重构紧性：
\[
  u_L=R_Lz_L
\]
在某个函数空间中预紧。

第三，PDE 残差闭性：
\[
  u_L\to u_*,
  \qquad
  \mathcal P_L[u_L]\to0
  \Longrightarrow
  \mathcal P[u_*]=0.
\]

第四，命中泛函一致收敛：
\[
  \lVert
  \widehat{\mathcal C}_{T,N,L}-\mathcal C_T
  \rVert_{L^\infty}\to0.
\]

第五，雅可比间隙稳定：
\[
  \lVert
  \nabla\widehat{\mathfrak G}_{k,N,L}
  -
  \nabla\mathfrak G_k
  \rVert_{L^\infty}\to0.
\]

第六，高阶正交残差非零：
\[
  b_k\neq0
\]
在仍需加速的阶上成立。
\end{remark}

\begin{remark}[最终严格表述]
用户命题可严格写成：
\[
  \boxed{
  \begin{gathered}
  \text{统计力学离散同伦扭曲命中范式，对高阶 PDE 的决定意义在于：}\\
  \text{它把直接求解高阶微分算子方程，}\\
  \text{替换为对离散残差能量低能集合的命中泛函轨迹优化。}
  \end{gathered}
  }
\]
进一步，
\[
  \boxed{
  \begin{gathered}
  \text{高阶正交提供命中方向的层级分解，}\\
  \text{雅可比间隙提供每阶同伦扭曲的工程收敛指引，}\\
  N\to\infty,\ L\to\infty
  \text{ 后给出确定性极限收敛。}
  \end{gathered}
  }
\]
最压缩的公式是：
\[
  \mathcal P[u]=0
\]
被替换为
\[
  E_L(z)\le\varepsilon,
  \qquad
  z\sim\pi_{\lambda,L},
  \qquad
  X_t^{\Gamma,L}\in S_{\varepsilon,L}.
\]
算法目标从
\[
  \text{求 }u
\]
变成
\[
  \Gamma^\star
  \in
  \arg\max_\Gamma
  \mathcal C_{T,L}(\Gamma).
\]
每阶收敛指引为
\[
  \min_{\theta_k,d_k}
  \mathfrak G_k(\theta_k,d_k).
\]
极限恢复为
\[
  N\to\infty,\quad L\to\infty,\quad \varepsilon\to0,
\]
\[
  R_Lz_{N,L}\to u_*,
  \qquad
  \mathcal P[u_*]=0,
  \qquad
  \mathcal B[u_*]=0.
\]
最终可表述为：
\[
  \boxed{
  \begin{gathered}
  \text{这是对高阶 PDE 解方程范式的决定性替代框架：}\\
  \text{以统计力学同伦命中取代直接微分算子求解，}\\
  \text{以高阶正交和雅可比间隙下降保证工程截断下的确定性收敛，}\\
  \text{并在 }N,L\to\infty\text{ 的极限中恢复严格 PDE 解。}
  \end{gathered}
  }
\]
但必须保留边界：
\[
  \boxed{
  \begin{gathered}
  \text{若没有残差相容、紧性、闭性、一致收敛和雅可比稳定性，}\\
  \text{则问题转化仍然成立；}\\
  \text{只是不能进一步推出该转化在极限中与原 PDE 求解完全等价。}
  \end{gathered}
  }
\]
\end{remark}

\clearpage

%%%%%%%%%%%%%%%%%%%%%%%%%%%%%%%%%%%%%%%%%%%%%%%%%%%%%%%%%%%%
\section{形式逻辑：高阶 PDE 同伦命中范式中的问题转化与极限保真性分离}
\label{sec:tex55-pde-transformation-fidelity-separation}
%%%%%%%%%%%%%%%%%%%%%%%%%%%%%%%%%%%%%%%%%%%%%%%%%%%%%%%%%%%%

\begin{remark}[层级分离口径]
上一节把高阶 PDE 的直接求解范式转化为统计力学离散同伦扭曲命中范式。
这里进一步区分构造层与保真层：
\[
  \boxed{
  \text{问题转化在构造层成立；极限解恢复属于保真层命题。}
  }
\]
在缺少保真条件时，不能推出转化后的命中解必然在极限中恢复原 PDE 解；
但这不影响构造层的问题转化成立。
因此必须区分：
\[
  \boxed{
  \text{问题转化成立}
  \neq
  \text{极限解恢复定理成立}.
  }
\]
前者是构造定义；后者是带条件的保真性定理。
\end{remark}

\subsection{公理降级为定义：原问题、命中问题与转化算子}
\label{subsec:tex55-transformation-definitions}

\begin{definition}[原高阶 PDE 问题]
\label{def:tex55-transform-original-pde}
令 \(D\subset\RR^d\) 为定义域，令 \(\partial D\) 为其边界。
给定高阶偏微分算子与边界算子
\[
  \mathcal P[u]=0,\qquad x\in D,
\]
\[
  \mathcal B[u]=0,\qquad x\in\partial D.
\]
记原高阶 PDE 问题为
\[
  \mathsf{PDE}
  \DefEq
  (\mathcal P,\mathcal B,D).
\]
这里的 \(\mathsf{PDE}\) 表示“直接寻找满足微分方程与边界条件的函数 \(u\)”这一问题对象。
\end{definition}

\begin{definition}[统计力学离散同伦扭曲命中问题]
\label{def:tex55-transform-hit-problem}
对每个离散尺度 \(L\)，取离散状态空间 \(\Xi_L\)，重构算子
\[
  R_L:\Xi_L\longrightarrow\mathcal U_L,
\]
离散 PDE 残差算子 \(\mathcal P_L\)，离散边界残差算子 \(\mathcal B_L\)，以及正则项 \(\mathcal R_L\)。
定义残差能量
\[
  E_L(z)
  \DefEq
  \lVert \mathcal P_L[R_Lz]\rVert^2
  +\alpha\lVert \mathcal B_L[R_Lz]\rVert^2
  +\lambda\mathcal R_L(z),
  \qquad z\in\Xi_L.
\]
给定容忍度 \(\varepsilon\ge0\)，定义低残差命中集合
\[
  S_{\varepsilon,L}
  \DefEq
  \{z\in\Xi_L:E_L(z)\le\varepsilon\}.
\]
给定分布参数
\[
  \lambda_L=(\beta,w_L,K_L,\theta_1,\dots,\theta_{K_L}),
\]
定义统计力学分布
\[
  \pi_{\lambda,L}(z)
  \DefEq
  \frac{
  w_L(z)
  \exp\!\left(
  -\beta E_L(z)
  +\sum_{j=1}^{K_L}\theta_j\phi_{j,L}(z)
  \right)
  }{
  Z_{\lambda,L}
  }.
\]
若同伦扭曲轨迹函数 \(\Gamma\) 在第 \(t\) 步选择参数 \(\lambda_t\)，则生成随机变量
\[
  X_t^{\Gamma,L}\sim\pi_{\lambda_t,L}.
\]
命中泛函定义为
\[
  \mathcal J_{T,L}(\Gamma)
  \DefEq
  \mathbb P_\Gamma
  \left(
  \exists t\le T:
  X_t^{\Gamma,L}\in S_{\varepsilon,L}
  \right).
\]
转化后的命中问题记为
\[
  \mathsf{Hit}_{L,\varepsilon}
  \DefEq
  (\Xi_L,E_L,S_{\varepsilon,L},\pi_{\lambda,L},\mathcal J_{T,L},\Gamma).
\]
\end{definition}

\begin{definition}[问题转化算子]
\label{def:tex55-transform-operator}
定义问题转化算子
\[
  \mathfrak T_L:
  \mathsf{PDE}
  \longmapsto
  \mathsf{Hit}_{L,\varepsilon}
\]
为如下构造：
\[
  (\mathcal P,\mathcal B,D)
  \longmapsto
  (\Xi_L,E_L,S_{\varepsilon,L},\pi_{\lambda,L},\mathcal J_{T,L},\Gamma).
\]
也就是说，
\[
  \boxed{
  \mathfrak T_L
  \text{ 的存在本身就是问题转化。}
  }
\]
这一步只要求转化对象被定义，不要求已经证明极限解恢复。
\end{definition}

\begin{Certificate}[构造层证书]
\label{cert:tex55-transform-constructive-certificate}
目标谓词为
\[
  P_{\mathrm{tr}}(L):
  \mathfrak T_L(\mathsf{PDE})=\mathsf{Hit}_{L,\varepsilon}
  \text{ 是一个良定义的问题转化}.
\]
证书分解为：
\[
  \mathcal P,\mathcal B
  \Rightarrow
  \mathcal P_L,\mathcal B_L
  \Rightarrow
  E_L
  \Rightarrow
  S_{\varepsilon,L}
  \Rightarrow
  \pi_{\lambda,L}
  \Rightarrow
  X_t^{\Gamma,L}
  \Rightarrow
  \mathcal J_{T,L}.
\]
只要每个箭头右端对象存在，谓词 \(P_{\mathrm{tr}}(L)\) 成立。
\end{Certificate}

\subsection{命题：问题转化不依赖极限保真条件}
\label{subsec:tex55-transformation-independent}

\begin{proposition}[问题转化的构造性成立]
\label{prop:tex55-transformation-independent-of-fidelity}
只要对象
\[
  \Xi_L,\quad
  R_L,\quad
  \mathcal P_L,\quad
  \mathcal B_L,\quad
  E_L,\quad
  S_{\varepsilon,L},\quad
  \pi_{\lambda,L},\quad
  \mathcal J_{T,L}
\]
均已定义，则
\[
  \mathsf{PDE}
  \longmapsto
  \mathsf{Hit}_{L,\varepsilon}
\]
作为问题转化成立。
\end{proposition}

\begin{Certificate}[谓词因果链证书]
\label{cert:tex55-transformation-independent-of-fidelity}
命题 \ref{prop:tex55-transformation-independent-of-fidelity} 的谓词因果链为
\[
  \mathsf{PDE}=(\mathcal P,\mathcal B,D)
\]
\[
  \Longrightarrow
\]
\[
  \text{选择离散状态空间与重构算子 }(\Xi_L,R_L)
\]
\[
  \Longrightarrow
\]
\[
  \text{构造离散残差能量 }E_L
\]
\[
  \Longrightarrow
\]
\[
  \text{定义低残差命中集合 }S_{\varepsilon,L}
\]
\[
  \Longrightarrow
\]
\[
  \text{定义统计力学分布 }\pi_{\lambda,L}
\]
\[
  \Longrightarrow
\]
\[
  \text{生成同伦扭曲随机变量簇 }X_t^{\Gamma,L}
\]
\[
  \Longrightarrow
\]
\[
  \text{定义命中泛函 }\mathcal J_{T,L}(\Gamma)
\]
\[
  \Longrightarrow
\]
\[
  \mathsf{Hit}_{L,\varepsilon}.
\]
该链条不包含 \(N\to\infty\)、\(L\to\infty\) 或闭性传递命题，
因此它证明的是构造性问题转化，而不是极限保真性。
\end{Certificate}

\begin{proof}
由定义 \ref{def:tex55-transform-original-pde}，原问题由
\[
  (\mathcal P,\mathcal B,D)
\]
给出。
由定义 \ref{def:tex55-transform-hit-problem}，一旦给定
\[
  \Xi_L,\ R_L,\ \mathcal P_L,\ \mathcal B_L,\ \mathcal R_L,
\]
就可定义残差能量
\[
  E_L(z)
  =
  \lVert \mathcal P_L[R_Lz]\rVert^2
  +\alpha\lVert \mathcal B_L[R_Lz]\rVert^2
  +\lambda\mathcal R_L(z).
\]
由 \(E_L\) 定义
\[
  S_{\varepsilon,L}
  =
  \{z\in\Xi_L:E_L(z)\le\varepsilon\}.
\]
由 \(E_L\) 与频域模态 \(\phi_{j,L}\) 定义
\[
  \pi_{\lambda,L}.
\]
由 \(\pi_{\lambda,L}\) 与轨迹函数 \(\Gamma\) 生成
\[
  X_t^{\Gamma,L}.
\]
由事件
\[
  X_t^{\Gamma,L}\in S_{\varepsilon,L}
\]
定义命中泛函
\[
  \mathcal J_{T,L}(\Gamma).
\]
于是原来的
\[
  \text{求 }u\text{ 使 }\mathcal P[u]=0,\ \mathcal B[u]=0
\]
被改写为
\[
  \max_{\Gamma}\mathcal J_{T,L}(\Gamma)
\]
或等价的命中指数优化
\[
  \max_{\Gamma}\mathcal C_{T,L}(\Gamma),
  \qquad
  \mathcal C_{T,L}=-\log Q_{T,L}.
\]
所以问题转化成立。
\end{proof}

\subsection{命题：需要条件的是极限保真性}
\label{subsec:tex55-fidelity-needs-conditions}

\begin{proposition}[极限保真性需要附加条件]
\label{prop:tex55-fidelity-needs-conditions}
若要从离散命中解
\[
  z_{N,L}\in S_{\varepsilon_{N,L},L}
\]
推出重构函数
\[
  u_{N,L}=R_Lz_{N,L}
\]
收敛到某个 \(u_*\)，并满足
\[
  \mathcal P[u_*]=0,
  \qquad
  \mathcal B[u_*]=0,
\]
则需要残差相容、重构紧性、PDE 残差闭性、命中泛函一致收敛、雅可比稳定性等保真条件。
\end{proposition}

\begin{Certificate}[保真性证书]
\label{cert:tex55-fidelity-needs-conditions}
保真性目标谓词为
\[
  P_{\mathrm{fid}}(N,L):
  z_{N,L}\in S_{\varepsilon_{N,L},L}
  \Longrightarrow
  R_Lz_{N,L}\to u_*,
  \quad
  \mathcal P[u_*]=0,
  \quad
  \mathcal B[u_*]=0.
\]
证书所需条件为
\[
  \varepsilon_{N,L}\to0,
\]
\[
  \mathcal P_L\to\mathcal P,
  \qquad
  \mathcal B_L\to\mathcal B,
\]
\[
  R_Lz_{N,L}\ \text{具有紧性},
\]
\[
  \text{PDE 残差图像在所用收敛拓扑下闭合},
\]
以及用于工程轨迹到理论轨迹传递的
\[
  \widehat{\mathcal C}_{T,N,L}\to\mathcal C_T,
  \qquad
  \nabla\widehat{\mathfrak G}_{k,N,L}\to\nabla\mathfrak G_k
\]
的一致收敛与稳定性。
\end{Certificate}

\begin{proof}
由
\[
  z_{N,L}\in S_{\varepsilon_{N,L},L}
\]
得
\[
  E_L(z_{N,L})\le\varepsilon_{N,L}.
\]
若
\[
  \varepsilon_{N,L}\to0,
\]
则由 \(E_L\) 的定义可推出离散残差趋零：
\[
  \lVert \mathcal P_L[R_Lz_{N,L}]\rVert\to0,
  \qquad
  \lVert \mathcal B_L[R_Lz_{N,L}]\rVert\to0.
\]
但是，要从离散残差趋零推出连续 PDE 残差趋零，需要相容性：
\[
  \mathcal P_L\to\mathcal P,
  \qquad
  \mathcal B_L\to\mathcal B.
\]
要从
\[
  u_{N,L}=R_Lz_{N,L}
\]
提取极限，需要紧性。
要把极限传递进算子，得到
\[
  \mathcal P[u_*]=0,
  \qquad
  \mathcal B[u_*]=0,
\]
需要残差闭性。
而若还要把有限样本、有限截断的工程同伦轨迹传递到理论极限轨迹，则需要
\[
  \widehat{\mathcal C}_{T,N,L}
  \to
  \mathcal C_T
\]
和
\[
  \nabla\widehat{\mathfrak G}_{k,N,L}
  \to
  \nabla\mathfrak G_k
\]
等一致收敛与雅可比稳定性条件。
因此这些条件证明的是
\[
  \mathsf{Hit}_{L,\varepsilon}
  \to
  \mathsf{PDE}
\]
的极限保真性，而不是证明问题转化本身。
\end{proof}

\subsection{定理：范式转化与保真定理的层级分离}
\label{subsec:tex55-transformation-fidelity-theorem}

\begin{theorem}[问题转化与极限保真性的层级分离]
\label{thm:tex55-transformation-fidelity-separation}
对高阶 PDE 问题 \(\mathsf{PDE}\)，存在两个逻辑层级：

第一层是构造层的问题转化：
\[
  \mathfrak T_L(\mathsf{PDE})
  =
  \mathsf{Hit}_{L,\varepsilon}.
\]
这一层只依赖定义 \ref{def:tex55-transform-operator} 中对象的可构造性。

第二层是极限层的保真性：
\[
  z_{N,L}\in S_{\varepsilon_{N,L},L},
  \quad
  N,L\to\infty,
  \quad
  \varepsilon_{N,L}\to0
  \Longrightarrow
  R_Lz_{N,L}\to u_*,
  \quad
  \mathcal P[u_*]=0,
  \quad
  \mathcal B[u_*]=0.
\]
这一层需要命题 \ref{prop:tex55-fidelity-needs-conditions} 中的保真条件。

因此
\[
  \boxed{
  \text{问题转化成立}
  }
\]
与
\[
  \boxed{
  \text{转化问题的命中解在极限中恢复原 PDE 解}
  }
\]
是两个不同的命题。
\end{theorem}

\begin{Certificate}[层级分离证书]
\label{cert:tex55-transformation-fidelity-separation}
定义两个谓词：
\[
  A_L:
  \mathfrak T_L(\mathsf{PDE})
  \text{ 良定义},
\]
\[
  B:
  \mathfrak T_L(\mathsf{PDE})
  \text{ 的低残差命中解在极限中恢复 PDE 解}.
\]
由命题 \ref{prop:tex55-transformation-independent-of-fidelity} 得
\[
  A_L
\]
由构造对象存在推出。
由命题 \ref{prop:tex55-fidelity-needs-conditions} 得
\[
  B
\]
需要相容性、紧性、闭性与稳定性。
故
\[
  A_L
  \not\Rightarrow
  B
  \quad
  \text{在无附加条件时不成立},
\]
但
\[
  A_L
\]
本身仍成立。
\end{Certificate}

\begin{proof}
由定义 \ref{def:tex55-transform-operator}，
\[
  \mathfrak T_L
\]
把
\[
  (\mathcal P,\mathcal B,D)
\]
映到
\[
  (\Xi_L,E_L,S_{\varepsilon,L},\pi_{\lambda,L},\mathcal J_{T,L},\Gamma).
\]
只要右侧对象被定义，则由命题
\ref{prop:tex55-transformation-independent-of-fidelity}，
\[
  \mathfrak T_L(\mathsf{PDE})
  =
  \mathsf{Hit}_{L,\varepsilon}
\]
作为问题转化成立。

另一方面，若要从
\[
  z_{N,L}\in S_{\varepsilon_{N,L},L}
\]
推出
\[
  R_Lz_{N,L}\to u_*,
  \qquad
  \mathcal P[u_*]=0,
  \qquad
  \mathcal B[u_*]=0,
\]
则必须把离散低残差、函数极限和连续算子方程连接起来。
由命题 \ref{prop:tex55-fidelity-needs-conditions}，
这一连接需要残差相容、紧性、闭性、一致收敛和雅可比稳定性。

所以，转化命题与保真命题不同：
\[
  \text{转化命题是构造命题，}
\]
\[
  \text{保真命题是极限恢复命题。}
\]
二者不能混写为同一个“是否已证明”的判断。
\end{proof}

\subsection{三层结构与最终口径}
\label{subsec:tex55-three-level-correction}

\begin{corollary}[三层正确表述]
\label{cor:tex55-three-level-correction}
高阶 PDE 同伦命中范式应分为三层：

第一层，范式转化：
\[
  \mathcal P[u]=0,\quad \mathcal B[u]=0
  \quad
  \leadsto
  \quad
  E_L(z)\le\varepsilon.
\]
也就是
\[
  \text{求 PDE 解}
  \quad
  \leadsto
  \quad
  \text{命中低残差集合}.
\]

第二层，工程截断确定性：
\[
  \text{随机命中}
  \quad
  \leadsto
  \quad
  \widehat{\mathcal J}_{T,N,L}(\Gamma),
\]
\[
  \text{同伦扭曲}
  \quad
  \leadsto
  \quad
  \min_{\theta_k,d_k}
  \widehat{\mathfrak G}_{k,N,L}(\theta_k,d_k).
\]

第三层，极限保真性：
\[
  N\to\infty,\qquad
  L\to\infty,\qquad
  \varepsilon\to0.
\]
在相容性、紧性、闭性与稳定性成立时，
\[
  R_Lz_{N,L}\to u_*,
  \qquad
  \mathcal P[u_*]=0,
  \qquad
  \mathcal B[u_*]=0.
\]
\end{corollary}

\begin{proof}
第一层直接由定义 \ref{def:tex55-transform-operator} 与命题
\ref{prop:tex55-transformation-independent-of-fidelity} 得出。
第二层来自经验命中泛函与雅可比间隙下降系统的有限样本、有限截断构造。
第三层由命题 \ref{prop:tex55-fidelity-needs-conditions} 给出的保真条件支撑。
三者依赖的前提不同，因此必须分层表述。
\end{proof}

\begin{remark}[层级分离后的最终表达]
严格结论应写为：
\[
  \boxed{
  \begin{gathered}
  \text{统计力学离散同伦扭曲命中范式对高阶 PDE 的意义，}\\
  \text{首先是问题转化：}\\
  \text{把直接求解高阶微分算子方程，}\\
  \text{转化为命中离散残差能量低能集合。}
  \end{gathered}
  }
\]
进一步：
\[
  \boxed{
  \begin{gathered}
  \text{高阶正交与雅可比间隙梯度下降}\\
  \text{提供的是转化后问题的工程截断确定性收敛机制。}
  \end{gathered}
  }
\]
再进一步：
\[
  \boxed{
  \begin{gathered}
  \text{残差相容、紧性、闭性、一致收敛、雅可比稳定性等条件，}\\
  \text{用于证明转化后命中解在 }N,L\to\infty\\
  \text{ 时恢复原 PDE 解的保真条件。}
  \end{gathered}
  }
\]
构造层的结论是：
\[
  \boxed{
  \text{问题转化成立。}
  }
\]
同时必须保留另一个独立层级：
\[
  \boxed{
  \text{从转化问题的命中解严格回到原 PDE 解，需要额外的保真性证书。}
  }
\]
\end{remark}

\begin{remark}[谓词因果链总表]
完整因果链为
\[
  \text{高阶 PDE 问题 }\mathsf{PDE}
\]
\[
  \Longrightarrow
\]
\[
  \text{问题转化算子 }\mathfrak T_L
\]
\[
  \Longrightarrow
\]
\[
  \text{离散残差能量 }E_L
\]
\[
  \Longrightarrow
\]
\[
  \text{低残差命中集合 }S_{\varepsilon,L}
\]
\[
  \Longrightarrow
\]
\[
  \text{统计力学分布 }\pi_{\lambda,L}
\]
\[
  \Longrightarrow
\]
\[
  \text{同伦扭曲随机变量簇 }X_t^{\Gamma,L}
\]
\[
  \Longrightarrow
\]
\[
  \text{命中泛函 }\mathcal J_{T,L}(\Gamma)
\]
\[
  \Longrightarrow
\]
\[
  \text{工程截断雅可比间隙下降}
\]
\[
  \Longrightarrow
\]
\[
  \text{若保真条件成立，则 }R_Lz_{N,L}\to u_*
\]
\[
  \Longrightarrow
\]
\[
  \mathcal P[u_*]=0,\qquad \mathcal B[u_*]=0.
\]
其中从 \(\mathsf{PDE}\) 到 \(\mathsf{Hit}_{L,\varepsilon}\) 是问题转化；
从 \(\mathsf{Hit}_{L,\varepsilon}\) 的命中解回到 \(u_*\) 是极限保真性。
\end{remark}

\clearpage

%%%%%%%%%%%%%%%%%%%%%%%%%%%%%%%%%%%%%%%%%%%%%%%%%%%%%%%%%%%%
\section{形式逻辑：问题转化后 PDE 显式求解效用的降级}
\label{sec:tex55-pde-explicit-solution-demotion}
%%%%%%%%%%%%%%%%%%%%%%%%%%%%%%%%%%%%%%%%%%%%%%%%%%%%%%%%%%%%

\begin{remark}[求解层口径]
只要问题转化成立，PDE 的显式表达就不再保留为\textbf{求解目标}。
严格说，PDE 不再作为直接求解对象，而是降级为
\[
  \boxed{
  \text{残差能量的定义来源}
  +
  \text{极限解的隐式存在性解释}
  +
  \text{保真性证书的语义背景}.
  }
\]
核心求解对象变成
\[
  \boxed{
  \text{统计力学离散同伦扭曲命中问题}.
  }
\]
这并不否定 PDE 的语义作用；它只是把 PDE 从求解层中心移到定义层和保真层。
\end{remark}

\subsection{公理降级为定义：定义层角色与求解层角色}
\label{subsec:tex55-explicit-demotion-definitions}

\begin{definition}[PDE 的定义层角色]
\label{def:tex55-pde-definition-layer-role}
原高阶 PDE 问题为
\[
  \mathcal P[u]=0,
  \qquad
  \mathcal B[u]=0.
\]
在问题转化中，它用于构造离散残差能量
\[
  E_L(z)
  \DefEq
  \lVert\mathcal P_L[R_Lz]\rVert^2
  +\alpha\lVert\mathcal B_L[R_Lz]\rVert^2
  +\lambda\mathcal R_L(z).
\]
因此 PDE 在定义层的角色是给出
\[
  \text{什么叫残差小}.
\]
也就是说，它定义命中集合
\[
  S_{\varepsilon,L}
  \DefEq
  \{z\in\Xi_L:E_L(z)\le\varepsilon\}.
\]
\end{definition}

\begin{definition}[求解层对象]
\label{def:tex55-pde-solving-layer-object}
一旦
\[
  E_L,
  \qquad
  S_{\varepsilon,L},
  \qquad
  \pi_{\lambda,L},
  \qquad
  \mathcal J_{T,L}
\]
被定义，求解层不再直接求
\[
  \mathcal P[u]=0,
  \qquad
  \mathcal B[u]=0.
\]
求解层的对象变为轨迹函数优化：
\[
  \Gamma^\star
  \in
  \arg\max_\Gamma
  \mathcal J_{T,L}(\Gamma),
\]
其中
\[
  \mathcal J_{T,L}(\Gamma)
  \DefEq
  \mathbb P_\Gamma
  \left(
  \exists t\le T:
  X_t^{\Gamma,L}\in S_{\varepsilon,L}
  \right).
\]
于是
\[
  \boxed{
  \text{求 PDE 解}
  \quad
  \leadsto
  \quad
  \text{命中低残差集合}.
  }
\]
\end{definition}

\begin{definition}[显式 PDE 求解效用降级]
\label{def:tex55-explicit-pde-demotion}
称 PDE 显式求解效用被降级，意思是：
\[
  \mathcal P[u]=0,\ \mathcal B[u]=0
\]
不再是算法直接优化的目标式。
它只作为
\[
  E_L,\quad S_{\varepsilon,L},\quad \text{极限保真性}
\]
的语义来源。
因此算法不追求显式闭式解、显式积分公式或显式高阶算子反演，而追求
\[
  X_t^{\Gamma,L}\in S_{\varepsilon,L}
\]
和
\[
  \Gamma^\star\in\arg\max_\Gamma\mathcal J_{T,L}(\Gamma).
\]
\end{definition}

\begin{Certificate}[层级角色证书]
\label{cert:tex55-pde-role-layer-certificate}
目标谓词为
\[
  P_{\mathrm{role}}:
  \text{PDE 在定义层保留语义，在求解层被命中问题替代}.
\]
谓词链为
\[
  \mathcal P,\mathcal B
  \Longrightarrow
  E_L
  \Longrightarrow
  S_{\varepsilon,L}
  \Longrightarrow
  \pi_{\lambda,L}
  \Longrightarrow
  X_t^{\Gamma,L}
  \Longrightarrow
  \mathcal J_{T,L}
  \Longrightarrow
  \Gamma^\star.
\]
其中只有第一箭头使用 PDE 显式表达；后续求解链以 \(E_L\)、\(S_{\varepsilon,L}\) 和
\(\mathcal J_{T,L}\) 为中心。
\end{Certificate}

\subsection{定理：PDE 求解问题到离散统计力学同伦命中问题的转化}
\label{subsec:tex55-pde-to-hit-theorem}

\begin{theorem}[问题转化定理]
\label{thm:tex55-pde-to-hit-problem-transformation}
设原问题为
\[
  \mathsf P:
  \quad
  \text{求 }u\in\mathcal U
  \text{ 使 }
  \mathcal P[u]=0,
  \quad
  \mathcal B[u]=0.
\]
构造离散状态空间 \(\Xi_L\)、重构映射
\[
  R_L:\Xi_L\to\mathcal U_L,
\]
残差能量
\[
  E_L:\Xi_L\to\RR_{\ge0},
\]
命中集合
\[
  S_{\varepsilon,L}
  =
  \{z\in\Xi_L:E_L(z)\le\varepsilon\},
\]
统计力学分布
\[
  \pi_{\lambda,L}(z)
  =
  \frac{
  w_L(z)
  \exp\!\left(
  -\beta E_L(z)
  +\sum_{j=1}^{K_L}\theta_j\phi_{j,L}(z)
  \right)
  }{
  Z_{\lambda,L}
  }.
\]
令
\[
  X_t^{\Gamma,L}\sim\pi_{\lambda_t,L},
  \qquad
  \lambda_t=\Gamma_t(\mathcal F_{t-1}),
\]
并定义
\[
  \mathcal J_{T,L}(\Gamma)
  =
  \mathbb P_\Gamma
  \left(
  \exists t\le T:
  X_t^{\Gamma,L}\in S_{\varepsilon,L}
  \right).
\]
则原问题被转化为
\[
  \boxed{
  \mathsf P
  \quad
  \longmapsto
  \quad
  \mathsf H_{L,\varepsilon}:
  \max_\Gamma\mathcal J_{T,L}(\Gamma).
  }
\]
\end{theorem}

\begin{Certificate}[转化证书]
\label{cert:tex55-pde-to-hit-transformation}
目标谓词为
\[
  P_{\mathrm{hit}}:
  \mathsf P
  \longmapsto
  \mathsf H_{L,\varepsilon}.
\]
证书链为
\[
  \mathcal P,\mathcal B
  \Longrightarrow
  E_L
  \Longrightarrow
  S_{\varepsilon,L}
  \Longrightarrow
  \pi_{\lambda,L}
  \Longrightarrow
  X_t^{\Gamma,L}
  \Longrightarrow
  Y_t^{\Gamma,L}
  \Longrightarrow
  \mathcal J_{T,L}(\Gamma)
  \Longrightarrow
  \max_\Gamma\mathcal J_{T,L}(\Gamma),
\]
其中
\[
  Y_t^{\Gamma,L}
  \DefEq
  \mathbf 1_{S_{\varepsilon,L}}(X_t^{\Gamma,L}).
\]
\end{Certificate}

\begin{proof}
由 \(\mathcal P,\mathcal B\) 构造残差能量 \(E_L\)。
由 \(E_L\) 构造低残差集合
\[
  S_{\varepsilon,L}
  =
  \{z:E_L(z)\le\varepsilon\}.
\]
由 \(E_L\)、权重 \(w_L\)、模态 \(\phi_{j,L}\)、逆温度 \(\beta\) 与同伦参数
\(\theta_j\) 构造统计力学分布 \(\pi_{\lambda,L}\)。
由轨迹函数 \(\Gamma\) 选择 \(\lambda_t\)，得到随机变量
\[
  X_t^{\Gamma,L}\sim\pi_{\lambda_t,L}.
\]
定义命中指示变量
\[
  Y_t^{\Gamma,L}
  =
  \mathbf 1_{S_{\varepsilon,L}}(X_t^{\Gamma,L}).
\]
于是
\[
  \mathcal J_{T,L}(\Gamma)
  =
  \mathbb P_\Gamma
  \left(
  \sum_{t=1}^{T}Y_t^{\Gamma,L}\ge1
  \right).
\]
因此，原来的“求 \(u\)”问题已经被改写为
\[
  \text{选择轨迹函数 }\Gamma
  \text{，使命中低残差集合的概率最大}.
\]
所以问题转化成立。
\end{proof}

\subsection{命题：残差能量确定后 PDE 显式表达不再参与求解}
\label{subsec:tex55-explicit-expression-irrelevance}

\begin{proposition}[显式表达求解效用降级]
\label{prop:tex55-explicit-expression-demoted}
若两个 PDE 表达
\[
  (\mathcal P_1,\mathcal B_1),
  \qquad
  (\mathcal P_2,\mathcal B_2)
\]
诱导同一个离散残差能量
\[
  E_L^{(1)}(z)=E_L^{(2)}(z),
  \qquad
  \forall z\in\Xi_L,
\]
且采用同一 \(w_L,\phi_{j,L},\beta,\theta_j\)，则它们诱导同一个命中集合、同一个统计力学分布和同一个命中泛函：
\[
  S_{\varepsilon,L}^{(1)}
  =
  S_{\varepsilon,L}^{(2)},
\]
\[
  \pi_{\lambda,L}^{(1)}
  =
  \pi_{\lambda,L}^{(2)},
\]
\[
  \mathcal J_{T,L}^{(1)}(\Gamma)
  =
  \mathcal J_{T,L}^{(2)}(\Gamma).
\]
因此在求解层，PDE 显式表达不再具有直接区分效用。
\end{proposition}

\begin{Certificate}[等残差能量证书]
\label{cert:tex55-equal-energy-same-hit-problem}
目标谓词为
\[
  P_{\mathrm{same}}:
  E_L^{(1)}=E_L^{(2)}
  \Longrightarrow
  \mathsf H_{L,\varepsilon}^{(1)}
  =
  \mathsf H_{L,\varepsilon}^{(2)}.
\]
因果链为
\[
  E_L^{(1)}=E_L^{(2)}
  \Longrightarrow
  S_{\varepsilon,L}^{(1)}=S_{\varepsilon,L}^{(2)}
  \Longrightarrow
  \pi_{\lambda,L}^{(1)}=\pi_{\lambda,L}^{(2)}
  \Longrightarrow
  X_t^{(1),\Gamma,L}\stackrel{\mathcal L}{=}X_t^{(2),\Gamma,L}
  \Longrightarrow
  \mathcal J_{T,L}^{(1)}=\mathcal J_{T,L}^{(2)}.
\]
\end{Certificate}

\begin{proof}
若
\[
  E_L^{(1)}(z)=E_L^{(2)}(z),
  \qquad
  \forall z\in\Xi_L,
\]
则对任意 \(\varepsilon\ge0\)，
\[
  \{z:E_L^{(1)}(z)\le\varepsilon\}
  =
  \{z:E_L^{(2)}(z)\le\varepsilon\}.
\]
所以
\[
  S_{\varepsilon,L}^{(1)}
  =
  S_{\varepsilon,L}^{(2)}.
\]
又因为统计力学分布只依赖
\[
  E_L,
  \qquad
  w_L,
  \qquad
  \phi_{j,L},
  \qquad
  \beta,
  \qquad
  \theta_j,
\]
所以在这些对象相同的条件下
\[
  \pi_{\lambda,L}^{(1)}
  =
  \pi_{\lambda,L}^{(2)}.
\]
于是对应随机变量簇的分布律相同：
\[
  X_t^{(1),\Gamma,L}
  \stackrel{\mathcal L}{=}
  X_t^{(2),\Gamma,L}.
\]
命中事件相同，故命中泛函相同：
\[
  \mathcal J_{T,L}^{(1)}(\Gamma)
  =
  \mathcal J_{T,L}^{(2)}(\Gamma).
\]
因此求解层无法区分这两个 PDE 显式表达。
\end{proof}

\subsection{引理：PDE 角色降级为隐式解释}
\label{subsec:tex55-pde-implicit-role}

\begin{lemma}[PDE 作为残差语义来源]
\label{lem:tex55-pde-residual-semantics}
问题转化后，PDE 不消失。
它的角色由显式求解对象降级为
\[
  \boxed{
  \text{隐式存在性解释}
  +
  \text{残差保真性解释}
  +
  \text{极限解语义解释}.
  }
\]
也就是说，\(\mathcal P[u]=0\) 不再作为算法直接求解的方程，
而是解释为什么
\[
  E_L(z)\approx0
\]
表示 \(R_Lz\) 接近 PDE 解。
\end{lemma}

\begin{Certificate}[语义保留证书]
\label{cert:tex55-pde-implicit-role}
语义保留的谓词链为
\[
  \mathcal P[u]=0,\ \mathcal B[u]=0
\]
\[
  \Longrightarrow
\]
\[
  E_L(z)
\]
\[
  \Longrightarrow
\]
\[
  E_L(z)\approx0
\]
\[
  \Longrightarrow
\]
\[
  R_Lz\text{ 近似满足 PDE 残差}
\]
\[
  \Longrightarrow
\]
\[
  \text{极限保真性解释}.
\]
该链条说明 PDE 在定义层与保真层保留语义，但不再是求解层的直接优化目标。
\end{Certificate}

\begin{proof}
由定义 \ref{def:tex55-pde-definition-layer-role}，PDE 表达被用于构造 \(E_L\)。
若
\[
  E_L(z)\approx0,
\]
则在残差能量所用范数下，
\[
  \mathcal P_L[R_Lz]\approx0,
  \qquad
  \mathcal B_L[R_Lz]\approx0.
\]
这说明 \(R_Lz\) 是低残差重构对象。
若进一步满足残差相容、紧性和闭性，则低残差重构序列在极限中解释为 PDE 解。
因此 PDE 的作用是提供残差语义与极限解释，而不是在求解层被直接反演或显式求解。
\end{proof}

\subsection{推论：工程等价与数学严格等价}
\label{subsec:tex55-engineering-and-mathematical-equivalence}

\begin{corollary}[工程意义等价]
\label{cor:tex55-engineering-equivalent-pde-solution}
若算法输出
\[
  z_{N,L}\in S_{\varepsilon,L},
\]
则
\[
  E_L(z_{N,L})\le\varepsilon.
\]
令
\[
  u_{N,L}=R_Lz_{N,L}.
\]
若工程容差满足
\[
  \varepsilon\le\varepsilon_{\mathrm{app}},
\]
并且应用验证指标满足
\[
  \operatorname{Err}_{\mathrm{app}}(u_{N,L})
  \le
  \delta_{\mathrm{app}},
\]
则 \(u_{N,L}\) 可称为工程等价 PDE 解。
\end{corollary}

\begin{proof}
由
\[
  z_{N,L}\in S_{\varepsilon,L}
\]
和 \(S_{\varepsilon,L}\) 的定义，直接得到
\[
  E_L(z_{N,L})\le\varepsilon.
\]
若
\[
  \varepsilon\le\varepsilon_{\mathrm{app}},
\]
则残差能量达到工程容差。
再由应用验证指标
\[
  \operatorname{Err}_{\mathrm{app}}(u_{N,L})
  \le
  \delta_{\mathrm{app}},
\]
可判定重构函数 \(u_{N,L}=R_Lz_{N,L}\) 满足工程应用要求。
所以它是工程意义下的等价 PDE 解。
\end{proof}

\begin{corollary}[数学严格等价]
\label{cor:tex55-mathematical-equivalent-pde-solution}
若存在序列 \(z_{N,L}\) 满足
\[
  E_L(z_{N,L})\to0,
\]
且
\[
  u_{N,L}=R_Lz_{N,L}\to u_*
\]
在目标函数空间中成立，并且残差极限传递成立，则
\[
  \mathcal P[u_*]=0,
  \qquad
  \mathcal B[u_*]=0.
\]
此时 \(u_*\) 是数学严格意义下的 PDE 解。
\end{corollary}

\begin{proof}
由
\[
  E_L(z_{N,L})\to0
\]
得到
\[
  \mathcal P_L[R_Lz_{N,L}]\to0,
  \qquad
  \mathcal B_L[R_Lz_{N,L}]\to0.
\]
又由
\[
  u_{N,L}=R_Lz_{N,L}\to u_*
\]
和残差极限传递条件，可推出
\[
  \mathcal P[u_*]=0,
  \qquad
  \mathcal B[u_*]=0.
\]
所以 \(u_*\) 是严格 PDE 解。
\end{proof}

\subsection{总定理与谓词因果链}
\label{subsec:tex55-explicit-demotion-final-theorem}

\begin{theorem}[显式 PDE 求解效用降级定理]
\label{thm:tex55-explicit-pde-solution-demotion}
若高阶 PDE 问题
\[
  \mathcal P[u]=0,
  \qquad
  \mathcal B[u]=0
\]
被转化为统计力学离散同伦扭曲命中问题
\[
  \max_\Gamma\mathcal J_{T,L}(\Gamma),
\]
其中
\[
  \mathcal J_{T,L}(\Gamma)
  =
  \mathbb P_\Gamma
  \left(
  \exists t\le T:
  X_t^{\Gamma,L}\in S_{\varepsilon,L}
  \right),
\]
\[
  S_{\varepsilon,L}
  =
  \{z:E_L(z)\le\varepsilon\},
\]
则在求解层，PDE 的显式表达不再是必要求解对象。
若命中过程在工程意义下收敛，则得到工程等价 PDE 解。
若命中过程在数学极限下收敛，则得到严格 PDE 解。
因此 PDE 只作为
\[
  \boxed{
  \text{隐式存在性解释与残差语义来源}
  }
\]
存在，而不再发挥显式表达与显式求解效用。
\end{theorem}

\begin{Certificate}[总定理证书]
\label{cert:tex55-explicit-pde-solution-demotion}
总命题的谓词因果链为
\[
  \mathcal P[u]=0,\ \mathcal B[u]=0
\]
\[
  \Longrightarrow
\]
\[
  E_L(z)
\]
\[
  \Longrightarrow
\]
\[
  S_{\varepsilon,L}=\{z:E_L(z)\le\varepsilon\}
\]
\[
  \Longrightarrow
\]
\[
  X_t^{\Gamma,L}\in S_{\varepsilon,L}
\]
\[
  \Longrightarrow
\]
\[
  \mathcal J_{T,L}(\Gamma)
\]
\[
  \Longrightarrow
\]
\[
  \Gamma^\star\in\arg\max_\Gamma\mathcal J_{T,L}(\Gamma)
\]
\[
  \Longrightarrow
\]
\[
  \text{求解层中心对象为命中泛函轨迹优化}.
\]
若附加工程容差，则得到工程等价解。
若附加极限保真性，则得到数学严格解。
\end{Certificate}

\begin{proof}
由定理 \ref{thm:tex55-pde-to-hit-problem-transformation}，原 PDE 求解问题被转化为
\[
  \max_\Gamma\mathcal J_{T,L}(\Gamma).
\]
由命题 \ref{prop:tex55-explicit-expression-demoted}，一旦 \(E_L\) 与命中结构确定，
求解层只依赖
\[
  E_L,\quad S_{\varepsilon,L},\quad \pi_{\lambda,L},\quad \mathcal J_{T,L}.
\]
因此 PDE 的显式表达不再作为直接求解对象。
由引理 \ref{lem:tex55-pde-residual-semantics}，PDE 的角色转为残差语义来源和极限解释背景。
由推论 \ref{cor:tex55-engineering-equivalent-pde-solution}，工程容差收敛给出工程等价 PDE 解。
由推论 \ref{cor:tex55-mathematical-equivalent-pde-solution}，数学极限收敛和残差传递给出严格 PDE 解。
故定理成立。
\end{proof}

\begin{remark}[最终压缩表述]
最终可压缩为两句话：
\[
  \boxed{
  \text{证明问题转化成立后，求解目标就从“显式求 PDE”转为“命中低残差集合”。}
  }
\]
\[
  \boxed{
  \begin{gathered}
  \text{若命中系统在工程意义或数学极限意义下收敛，则得到等价 PDE 解；}\\
  \text{此时 PDE 只提供隐式存在性与残差语义解释，}\\
  \text{不再承担显式表达和显式求解功能。}
  \end{gathered}
  }
\]
\end{remark}

\clearpage

%%%%%%%%%%%%%%%%%%%%%%%%%%%%%%%%%%%%%%%%%%%%%%%%%%%%%%%%%%%%
\section{形式逻辑：障碍度量化、变分完整化与确定性隧穿控制}
\label{sec:tex55-metric-variational-deterministic-tunneling}
%%%%%%%%%%%%%%%%%%%%%%%%%%%%%%%%%%%%%%%%%%%%%%%%%%%%%%%%%%%%

\begin{remark}[本章定位]
本章把前述各章的结论收束为一个统一形式系统：
\[
  \boxed{
  \begin{gathered}
  \text{度量化墙，能量化孔洞，命中化穿透，正交化方向，}\\
  \text{雅可比化控制，最终得到确定性隧穿。}
  \end{gathered}
  }
\]
第 \ref{sec:tex55-main} 节已经证明正厚度孔洞通道与穿墙芯线的同伦等价；
该结论的拓扑背景可参考代数拓扑中的同伦等价与强变形收缩理论 \cite{Hatcher2002AT}。
统计力学命中、马尔可夫采样、退火与经验过程背景分别参考
\cite{LevinPeresWilmer2017Markov,Norris1998MarkovChains,Kirkpatrick1983Optimization,Hajek1988Cooling,VanDerVaartWellner1996EmpiricalProcesses}。
高阶 PDE、有限元残差与数值优化背景参考
\cite{Evans2010PDE,BrennerScott2008FEM,NocedalWright2006NumericalOptimization}；
变分算子、凸分析、梯度流和 PL 型确定性下降背景参考
\cite{GelfandFomin2000CalculusVariations,Rockafellar1970ConvexAnalysis,AmbrosioGigliSavare2008GradientFlows,KarimiNutiniSchmidt2016PL}。
本文的锚点式形式系统写法继承仓内语义锚点与 \(\GFramework\) 相关稿件
\cite{RepoTex37SemanticGauge,RepoTex51TruthAnchor,GaoZheng2025GFramework}；
其“雅可比间隙--同伦扭曲--微观隧穿命中--PDE 策略簇”的仓内前史分别见
\cite{RepoTex28JacGap,RepoTex45ObstructionGlobalSection,RepoTex53StrategyCluster}。
\end{remark}

\subsection{定义：墙、孔洞与确定性隧穿的度量对象}
\label{subsec:tex55-metric-tunneling-definitions}

\begin{definition}[墙与孔洞的度量结构]
\label{def:tex55-wall-hole-metric-structure}
令 \(M\) 为环境空间，\(W\subset M\) 为墙或障碍。
第 \ref{sec:tex55-main} 节中的障碍修复给出正厚度通道
\[
  T_k=D^k\times I
\]
及其穿墙芯线
\[
  C_k=(D^k\times\{0\})\cup(\{0\}\times I)\cup(D^k\times\{1\}).
\]
为了把“能否穿透”的拓扑问题提升为“如何可控穿透”的确定性优化问题，
在 \(W,T_k,C_k\) 上附加度量数据
\[
  d_W,\qquad
  E_W,\qquad
  \operatorname{gap}(W),\qquad
  \operatorname{cost}_\rho(W),\qquad
  \operatorname{cap}(T_k).
\]
其中 \(d_W\) 度量墙附近距离结构，\(E_W\) 度量障碍能量，
\(\operatorname{gap}(W)\) 度量孔洞开口或可穿透间隙，
\(\operatorname{cost}_\rho(W)\) 度量修复代价，
\(\operatorname{cap}(T_k)\) 度量通道容量。
\end{definition}

\begin{definition}[穿透能量与命中集合]
\label{def:tex55-through-energy-target}
令 \(S_{\mathrm{through}}\) 表示已经穿过墙、进入障碍另一侧、或命中低残差通道的目标集合。
定义墙或孔洞的穿透能量为
\[
  E_{\mathrm{wall}}(x)
  =
  \operatorname{dist}(x,C_k)^2
  +
  \alpha\,\operatorname{barrier}_W(x)
  +
  \lambda\,\operatorname{cost}_\rho(x),
\]
其中 \(\alpha,\lambda\ge0\)。
给定容差 \(\varepsilon\ge0\)，定义可控穿透命中集合
\[
  S_{\mathrm{through},\varepsilon}
  =
  \{x:E_{\mathrm{wall}}(x)\le\varepsilon\}.
\]
当本文处理高阶 PDE 时，\(E_{\mathrm{wall}}\) 的角色由离散残差能量
\[
  E_L(z)
  =
  |\mathcal P_L[R_Lz]|^2
  +
  \alpha|\mathcal B_L[R_Lz]|^2
  +
  \lambda\mathcal R_L(z)
\]
承担，而命中集合变为
\[
  S_{\varepsilon,L}
  =
  \{z:E_L(z)\le\varepsilon\}.
\]
\end{definition}

\begin{definition}[统计力学同伦隧穿分布]
\label{def:tex55-statistical-tunnel-law}
给定分布律索引
\[
  \lambda_t=(\beta_t,w_t,K_t,\theta_t),
\]
定义同伦隧穿分布
\[
  \pi_{\lambda_t}(x)
  =
  \frac{
  w_t(x)
  \exp\left(
  -\beta_tE_{\mathrm{wall}}(x)
  +
  \sum_{j=1}^{K_t}\theta_{t,j}\phi_j(x)
  \right)
  }{
  Z_{\lambda_t}
  }.
\]
由同伦轨迹函数 \(\Gamma\) 生成参数轨迹
\[
  \lambda_t=\Gamma_t(\mathcal F_{t-1}).
\]
相应的状态随机变量、命中指示变量和首次命中时间为
\[
  X_t^\Gamma\sim\pi_{\lambda_t},
  \qquad
  Y_t^\Gamma=\mathbf 1_{S_{\mathrm{through},\varepsilon}}(X_t^\Gamma),
  \qquad
  \tau_{\mathrm{through}}^\Gamma=\inf\{t\ge1:Y_t^\Gamma=1\}.
\]
\end{definition}

\begin{definition}[命中泛函与命中指数]
\label{def:tex55-through-hit-functional-index}
定义有限时间穿透命中泛函
\[
  \mathcal J_T(\Gamma)
  =
  \mathbb P_\Gamma(\tau_{\mathrm{through}}^\Gamma\le T).
\]
定义未命中生存概率
\[
  Q_T(\Gamma)
  =
  \mathbb P_\Gamma(\tau_{\mathrm{through}}^\Gamma>T).
\]
定义命中指数
\[
  \mathcal C_T(\Gamma)
  =
  -\log Q_T(\Gamma).
\]
于是
\[
  \mathcal J_T(\Gamma)=1-e^{-\mathcal C_T(\Gamma)}.
\]
所以最大化穿透命中概率等价于最大化命中指数。
\end{definition}

\begin{definition}[变分算子的三层确定性扩展]
\label{def:tex55-three-layer-variational-completion}
传统变分法中的核心对象为
\[
  \mathcal E[u],
  \qquad
  \delta\mathcal E[u]=0,
  \qquad
  \nabla\mathcal E[u],
  \qquad
  D^2\mathcal E[u].
\]
本文把它扩展为三层确定性变分结构：
\[
  z\mapsto E_L(z),
  \qquad
  \Gamma\mapsto\mathcal C_T(\Gamma),
  \qquad
  (\theta_k,d_k)\mapsto\mathfrak G_k(\theta_k,d_k).
\]
第一层是残差能量变分；第二层是命中指数变分；第三层是雅可比间隙变分。
其中第 \(k\) 阶理想正交命中方向为
\[
  b_k
  =
  \Pi_{\mathcal W_{k-1}^{\perp}}
  \nabla\mathcal C_T(\Gamma^{[k-1]}),
\]
实际同伦扭曲雅可比方向为
\[
  J_k(\theta_k)d_k,
\]
雅可比间隙为
\[
  \mathfrak G_k(\theta_k,d_k)
  =
  \frac12
  \left|
  \Pi_{\mathcal W_{k-1}^{\perp}}
  J_k(\theta_k)d_k
  -
  b_k
  \right|^2
  +
  \frac{\rho_k}{2}
  \left|
  \Pi_{\mathcal W_{k-1}}
  J_k(\theta_k)d_k
  \right|^2.
\]
\end{definition}

\subsection{公理降级为定理：障碍度量化推出可控穿透问题}
\label{subsec:tex55-metric-tunnel-axiom-theorem}

\begin{theorem}[障碍度量化可控穿透定理（公理降级）]
\label{thm:tex55-metric-controllable-tunneling}
设 \(W\subset M\) 是墙或障碍，且存在第 \ref{sec:tex55-main} 节意义下的修复通道 \(T_k\) 与穿墙芯线 \(C_k\)。
若在 \(W,T_k,C_k\) 上定义度量数据
\[
  d_W,\quad E_W,\quad \operatorname{gap}(W),
  \quad
  \operatorname{cost}_\rho(W),
  \quad
  \operatorname{cap}(T_k),
\]
并由此定义穿透能量 \(E_{\mathrm{wall}}\) 与命中集合 \(S_{\mathrm{through},\varepsilon}\)，
则障碍穿透问题可表述为确定性可控的同伦命中优化问题
\[
  \Gamma^\star
  \in
  \arg\max_\Gamma \mathcal J_T(\Gamma)
  =
  \arg\max_\Gamma
  \mathbb P_\Gamma
  \left(
  \exists t\le T:
  X_t^\Gamma\in S_{\mathrm{through},\varepsilon}
  \right).
\]
\end{theorem}

\begin{Certificate}[障碍度量化证书]
\label{cert:tex55-metric-controllable-tunneling}
该定理的谓词因果链为
\[
  W\text{ 是墙}
\]
\[
  \Longrightarrow
\]
\[
  T_k=D^k\times I\text{ 是正厚度修复通道}
\]
\[
  \Longrightarrow
\]
\[
  C_k\subset T_k\text{ 是穿墙芯线}
\]
\[
  \Longrightarrow
\]
\[
  T_k\searrow C_k
\]
\[
  \Longrightarrow
\]
\[
  X_\mu\simeq X_\times
\]
\[
  \Longrightarrow
\]
\[
  \text{穿透存在拓扑通路}
\]
\[
  \Longrightarrow
\]
\[
  d_W,E_W,\operatorname{gap}(W),\operatorname{cost}_\rho(W),\operatorname{cap}(T_k)
\]
\[
  \Longrightarrow
\]
\[
  E_{\mathrm{wall}},\quad S_{\mathrm{through},\varepsilon}
\]
\[
  \Longrightarrow
\]
\[
  X_t^\Gamma\in S_{\mathrm{through},\varepsilon}
\]
\[
  \Longrightarrow
\]
\[
  \mathcal J_T(\Gamma)
\]
\[
  \Longrightarrow
\]
\[
  \Gamma^\star\in\arg\max_\Gamma\mathcal J_T(\Gamma).
\]
\end{Certificate}

\begin{proof}
由第 \ref{sec:tex55-main} 节可知，正厚度通道 \(T_k\) 强变形收缩到穿墙芯线 \(C_k\)，
因此修复后微观隧穿空间与穿墙芯线空间同伦等价：
\[
  X_\mu\simeq X_\times.
\]
该结论给出穿透的拓扑可达通路。
在此基础上，定义墙与孔洞的度量数据后，得到穿透能量
\[
  E_{\mathrm{wall}}(x)
  =
  \operatorname{dist}(x,C_k)^2
  +
  \alpha\,\operatorname{barrier}_W(x)
  +
  \lambda\,\operatorname{cost}_\rho(x).
\]
由穿透能量定义命中集合
\[
  S_{\mathrm{through},\varepsilon}
  =
  \{x:E_{\mathrm{wall}}(x)\le\varepsilon\}.
\]
由同伦轨迹 \(\Gamma\) 定义分布律轨迹 \(\pi_{\lambda_t}\)，并由该分布律轨迹生成随机变量簇 \(X_t^\Gamma\)。
命中事件为
\[
  X_t^\Gamma\in S_{\mathrm{through},\varepsilon}.
\]
所以穿透问题被表述为最大化
\[
  \mathcal J_T(\Gamma)
  =
  \mathbb P_\Gamma
  \left(
  \exists t\le T:X_t^\Gamma\in S_{\mathrm{through},\varepsilon}
  \right).
\]
故障碍穿透被转化为确定性可控的同伦命中优化问题。
\end{proof}

\subsection{引理：穿墙芯线是确定性隧穿的几何锚点}
\label{subsec:tex55-spine-anchor-lemma}

\begin{lemma}[穿墙芯线锚点引理]
\label{lem:tex55-crosswall-spine-anchor}
第 \ref{sec:tex55-main} 节中的同伦等价
\[
  \MicroTunnel_\rho(W)\simeq\CrossWall_\rho(W)
\]
为确定性隧穿提供几何锚点。
更具体地，正厚度通道的全部横截面自由度可以收缩到保端口的穿墙芯线 \(C_k\)，
因此可控穿透只需要控制芯线级通路，而不必保留全部正厚度几何复杂度。
\end{lemma}

\begin{Certificate}[几何锚点证书]
\label{cert:tex55-crosswall-spine-anchor}
证明链为
\[
  T_k=D^k\times I
\]
\[
  \Longrightarrow
\]
\[
  C_k=(D^k\times\{0\})\cup(\{0\}\times I)\cup(D^k\times\{1\})
\]
\[
  \Longrightarrow
\]
\[
  T_k\searrow C_k
\]
\[
  \Longrightarrow
\]
\[
  X_0\cup T_k\searrow X_0\cup C_k
\]
\[
  \Longrightarrow
\]
\[
  \MicroTunnel_\rho(W)\simeq\CrossWall_\rho(W)
\]
\[
  \Longrightarrow
\]
\[
  \text{穿墙芯线是可控穿透的几何锚点}.
\]
\end{Certificate}

\begin{proof}
由第 \ref{sec:tex55-main} 节的强变形收缩构造，\(T_k\) 保端口强变形收缩到 \(C_k\)。
粘合到环境空间 \(X_0\) 后，该收缩下降为商空间上的强变形收缩
\[
  X_0\cup T_k\searrow X_0\cup C_k.
\]
因此
\[
  \MicroTunnel_\rho(W)\simeq\CrossWall_\rho(W).
\]
这说明正厚度隧道的横截面自由度在同伦意义下不是控制穿透的必要复杂度；
保端口芯线已经足以表示穿透通路。
所以穿墙芯线是确定性隧穿的几何锚点。
\end{proof}

\subsection{命题：新求解范式的对象层级价值}
\label{subsec:tex55-paradigm-value-object-hierarchy}

\begin{proposition}[对象层级完整性命题]
\label{prop:tex55-object-hierarchy-value}
统计力学同伦命中范式的求解对象不是单层分布函数，而是由以下对象组成的层级系统：
\[
  \pi_\lambda
  \Rightarrow
  X_t^\lambda
  \Rightarrow
  Y_t^\lambda
  \Rightarrow
  \tau_S^\lambda
  \Rightarrow
  \mathcal J_T(\lambda)
  \Rightarrow
  \mathcal C_T(\lambda).
\]
因此，命中严格表示为随机变量取值落入命中集合，而不是分布律本身“命中”。
\end{proposition}

\begin{Certificate}[对象层级证书]
\label{cert:tex55-object-hierarchy-value}
对象层级谓词链为
\[
  \lambda\in\Lambda
  \Longrightarrow
  \pi_\lambda
  \Longrightarrow
  X_t^\lambda\sim\pi_\lambda
  \Longrightarrow
  Y_t^\lambda=\mathbf 1_S(X_t^\lambda)
  \Longrightarrow
  \tau_S^\lambda=\inf\{t:Y_t^\lambda=1\}
  \Longrightarrow
  \mathcal J_T(\lambda)
  \Longrightarrow
  \mathcal C_T(\lambda).
\]
\end{Certificate}

\begin{proof}
由分布律定义，\(\pi_\lambda\) 只是随机变量 \(X_t^\lambda\) 的法律。
命中事件定义为
\[
  X_t^\lambda\in S.
\]
命中指示变量定义为
\[
  Y_t^\lambda=\mathbf 1_S(X_t^\lambda).
\]
首次命中时间定义为
\[
  \tau_S^\lambda=\inf\{t\ge1:Y_t^\lambda=1\}.
\]
有限时间命中泛函定义为
\[
  \mathcal J_T(\lambda)=\mathbb P(\tau_S^\lambda\le T).
\]
命中指数定义为
\[
  \mathcal C_T(\lambda)=-\log\mathbb P(\tau_S^\lambda>T).
\]
所以求解对象确实是由分布律、随机变量、命中事件、命中泛函和命中指数组成的层级系统。
\end{proof}

\subsection{定理：变分算子的确定性完整化}
\label{subsec:tex55-variational-completion-theorem}

\begin{theorem}[变分算子确定性完整化定理]
\label{thm:tex55-variational-operator-completion}
设传统变分问题由能量泛函
\[
  u\mapsto\mathcal E[u]
\]
及其驻值条件
\[
  \delta\mathcal E[u]=0
\]
给出。
若引入残差能量 \(E_L\)、命中指数 \(\mathcal C_T\)、高阶正交残差方向 \(b_k\)
以及雅可比间隙 \(\mathfrak G_k\)，则变分算子被扩展为确定性完整链条
\[
  \boxed{
  \text{残差能量变分}
  \Rightarrow
  \text{命中指数变分}
  \Rightarrow
  \text{正交残差变分}
  \Rightarrow
  \text{雅可比间隙变分}
  \Rightarrow
  \text{确定性收敛轨迹}.
  }
\]
\end{theorem}

\begin{Certificate}[变分完整化证书]
\label{cert:tex55-variational-operator-completion}
证明链为
\[
  \mathcal E[u],\delta\mathcal E[u]
\]
\[
  \Longrightarrow
\]
\[
  E_L(z)
  =
  |\mathcal P_L[R_Lz]|^2
  +
  \alpha|\mathcal B_L[R_Lz]|^2
  +
  \lambda\mathcal R_L(z)
\]
\[
  \Longrightarrow
\]
\[
  S_{\varepsilon,L}=\{z:E_L(z)\le\varepsilon\}
\]
\[
  \Longrightarrow
\]
\[
  \mathcal C_T(\Gamma)=-\log Q_T(\Gamma)
\]
\[
  \Longrightarrow
\]
\[
  b_k=
  \Pi_{\mathcal W_{k-1}^{\perp}}
  \nabla\mathcal C_T(\Gamma^{[k-1]})
\]
\[
  \Longrightarrow
\]
\[
  \mathfrak G_k(\theta_k,d_k)
\]
\[
  \Longrightarrow
\]
\[
  \text{确定性同伦轨迹下降}.
\]
\end{Certificate}

\begin{proof}
传统变分算子给出局部驻值条件
\[
  \delta\mathcal E[u]=0.
\]
对于高阶 PDE 或强非线性问题，仅使用该条件容易停留在局部极值或病态线性化层面。
引入离散残差能量 \(E_L\) 后，变分对象从函数 \(u\) 的局部极值转向低残差集合
\[
  S_{\varepsilon,L}=\{z:E_L(z)\le\varepsilon\}.
\]
引入命中泛函后，目标从“能量下降”提升为“命中低残差集合”：
\[
  \mathcal J_T(\Gamma)=\mathbb P_\Gamma(\tau_{S_{\varepsilon,L}}^\Gamma\le T).
\]
等价地，最大化命中指数
\[
  \mathcal C_T(\Gamma)=-\log Q_T(\Gamma)
\]
给出降低未命中生存概率的确定性方向
\[
  \nabla\mathcal C_T(\Gamma).
\]
为了避免重复已用方向，将该梯度投影到
\[
  \mathcal W_{k-1}^{\perp}
\]
得到高阶正交残差方向 \(b_k\)。
实际同伦扭曲只能产生雅可比方向 \(J_k(\theta_k)d_k\)，因此必须通过 \(\mathfrak G_k\) 使其逼近 \(b_k\) 并抑制正交泄漏。
于是得到从残差能量到命中指数、从命中梯度到正交残差、从理想方向到雅可比实现方向的完整确定性链条。
\end{proof}

\subsection{定理：微观隧穿的确定化与可控化}
\label{subsec:tex55-deterministic-tunneling-control-theorem}

\begin{theorem}[微观隧穿确定化可控化定理]
\label{thm:tex55-deterministic-controllable-microtunneling}
设 \(S\) 为穿透命中集合，\(\Gamma\) 为同伦轨迹函数。
令
\[
  g_k=\nabla\mathcal C_T(\Gamma^{[k-1]}),
  \qquad
  b_k=
  \Pi_{\mathcal W_{k-1}^{\perp}}g_k.
\]
若 \(b_k\neq0\)，且实际同伦扭曲雅可比方向
\[
  \delta_k=J_k(\theta_k)d_k
\]
满足间隙条件
\[
  |\delta_k-b_k|
  <
  \frac{|b_k|^2}{|g_k|},
\]
则存在足够小的步长 \(\eta_k>0\)，使第 \(k\) 阶扭曲后
\[
  \mathcal C_T(\Gamma^{[k]})
  >
  \mathcal C_T(\Gamma^{[k-1]}),
\]
从而
\[
  Q_T(\Gamma^{[k]})
  <
  Q_T(\Gamma^{[k-1]}).
\]
因此微观隧穿从概率性试探提升为由高阶正交方向与雅可比间隙控制的确定性同伦命中过程。
\end{theorem}

\begin{Certificate}[确定性隧穿证书]
\label{cert:tex55-deterministic-controllable-microtunneling}
证书链为
\[
  \Gamma^{[k-1]}
  \Longrightarrow
  \mathcal C_T(\Gamma^{[k-1]})
  \Longrightarrow
  g_k=\nabla\mathcal C_T(\Gamma^{[k-1]})
\]
\[
  \Longrightarrow
\]
\[
  b_k=\Pi_{\mathcal W_{k-1}^{\perp}}g_k
\]
\[
  \Longrightarrow
\]
\[
  \delta_k=J_k(\theta_k)d_k
\]
\[
  \Longrightarrow
\]
\[
  |\delta_k-b_k|
  <
  \frac{|b_k|^2}{|g_k|}
\]
\[
  \Longrightarrow
\]
\[
  D\mathcal C_T(\Gamma^{[k-1]})[\delta_k]>0
\]
\[
  \Longrightarrow
\]
\[
  \mathcal C_T(\Gamma^{[k]})>\mathcal C_T(\Gamma^{[k-1]})
\]
\[
  \Longrightarrow
\]
\[
  Q_T(\Gamma^{[k]})<Q_T(\Gamma^{[k-1]}).
\]
\end{Certificate}

\begin{proof}
由 \(b_k=\Pi_{\mathcal W_{k-1}^{\perp}}g_k\)，有
\[
  \langle g_k,b_k\rangle=|b_k|^2.
\]
令
\[
  e_k=\delta_k-b_k.
\]
则
\[
  \langle g_k,\delta_k\rangle
  =
  \langle g_k,b_k+e_k\rangle
  =
  |b_k|^2+\langle g_k,e_k\rangle.
\]
由 Cauchy 不等式，
\[
  \langle g_k,\delta_k\rangle
  \ge
  |b_k|^2-|g_k||e_k|.
\]
若
\[
  |e_k|<
  \frac{|b_k|^2}{|g_k|},
\]
则
\[
  \langle g_k,\delta_k\rangle>0.
\]
因此
\[
  D\mathcal C_T(\Gamma^{[k-1]})[\delta_k]>0.
\]
由 \(\mathcal C_T\) 的可微性，存在足够小的 \(\eta_k>0\)，使得沿 \(\delta_k\) 更新后
\[
  \mathcal C_T(\Gamma^{[k]})
  >
  \mathcal C_T(\Gamma^{[k-1]}).
\]
又因为
\[
  Q_T(\Gamma)=e^{-\mathcal C_T(\Gamma)},
\]
故
\[
  Q_T(\Gamma^{[k]})
  <
  Q_T(\Gamma^{[k-1]}).
\]
这说明第 \(k\) 阶扭曲提高了命中指数并降低了未命中生存概率。
\end{proof}

\subsection{数学归纳法证书：高阶确定性隧穿的幂化控制}
\label{subsec:tex55-deterministic-tunnel-induction}

\begin{theorem}[高阶确定性隧穿幂化控制定理]
\label{thm:tex55-deterministic-tunnel-power-control}
若对每一阶 \(k=1,\ldots,n\) 均存在幂化因子 \(A_k>1\)，使得
\[
  Q_T(\Gamma^{[k]})
  \le
  Q_T(\Gamma^{[k-1]})^{A_k},
\]
则
\[
  \boxed{
  Q_T(\Gamma^{[n]})
  \le
  Q_T(\Gamma^{[0]})^{A_1A_2\cdots A_n}.
  }
\]
因此
\[
  \boxed{
  \mathcal J_T(\Gamma^{[n]})
  \ge
  1-
  Q_T(\Gamma^{[0]})^{A_1A_2\cdots A_n}.
  }
\]
\end{theorem}

\begin{Certificate}[数学归纳法证书]
\label{cert:tex55-deterministic-tunnel-power-control}
定义目标谓词
\[
  P(n):
  \quad
  Q_T(\Gamma^{[n]})
  \le
  Q_T(\Gamma^{[0]})^{\prod_{k=1}^{n}A_k}.
\]

基例 \(n=1\)：
由假设
\[
  Q_T(\Gamma^{[1]})
  \le
  Q_T(\Gamma^{[0]})^{A_1},
\]
所以 \(P(1)\) 成立。

归纳步：
假设 \(P(n)\) 成立，即
\[
  Q_T(\Gamma^{[n]})
  \le
  Q_T(\Gamma^{[0]})^{\prod_{k=1}^{n}A_k}.
\]
又由第 \(n+1\) 阶幂化压缩条件
\[
  Q_T(\Gamma^{[n+1]})
  \le
  Q_T(\Gamma^{[n]})^{A_{n+1}}.
\]
代入归纳假设，得
\[
  Q_T(\Gamma^{[n+1]})
  \le
  \left(
  Q_T(\Gamma^{[0]})^{\prod_{k=1}^{n}A_k}
  \right)^{A_{n+1}}
  =
  Q_T(\Gamma^{[0]})^{\prod_{k=1}^{n+1}A_k}.
\]
故 \(P(n+1)\) 成立。

由数学归纳法，
\[
  \forall n\ge1,\quad P(n)
\]
成立。
\end{Certificate}

\begin{proof}
证书已经给出完整归纳链。
由 \(Q_T(\Gamma)=1-\mathcal J_T(\Gamma)\)，得到
\[
  \mathcal J_T(\Gamma^{[n]})
  =
  1-Q_T(\Gamma^{[n]})
  \ge
  1-
  Q_T(\Gamma^{[0]})^{\prod_{k=1}^{n}A_k}.
\]
定理得证。
\end{proof}

\subsection{总定理：范式价值、变分完整性与障碍穿透控制}
\label{subsec:tex55-value-variational-tunnel-final-theorem}

\begin{theorem}[同伦命中 PDE 范式替代的总价值定理]
\label{thm:tex55-total-value-variational-tunnel}
本文构造的统计力学同伦命中范式具有以下八项形式价值：
\[
\begin{aligned}
  &\text{(1) PDE 求解被转化为低残差集合命中；}\\
  &\text{(2) 分布律、随机变量、命中事件、命中泛函层级被严格区分；}\\
  &\text{(3) 加速被形式化为生存概率下降、命中指数增长与幂化压缩；}\\
  &\text{(4) 雅可比间隙连接理想正交命中方向与工程可实现方向；}\\
  &\text{(5) 随机采样经经验泛函、轨迹函数和间隙下降进入确定性极限；}\\
  &\text{(6) 高阶 PDE 的困难被迁移到残差能量、命中集合与同伦轨迹控制；}\\
  &\text{(7) 问题转化与极限保真性被分层处理；}\\
  &\text{(8) 局部算法被提升为分布参数轨迹函数，}\\
  &\qquad\text{并通过频域结构与高阶正交残差实现全局命中桥接。}
\end{aligned}
\]
因此该范式给出一个由
\[
  \boxed{
  \text{随机变量簇、命中泛函、高阶正交残差、雅可比间隙、确定性极限}
  }
\]
共同组成的新型 PDE 求解替代框架。
\end{theorem}

\begin{Certificate}[总价值证书]
\label{cert:tex55-total-value-variational-tunnel}
总谓词因果链为
\[
  \text{墙或 PDE 障碍}
\]
\[
  \Longrightarrow
\]
\[
  \text{残差能量或穿透能量}
\]
\[
  \Longrightarrow
\]
\[
  \text{低残差命中集合}
\]
\[
  \Longrightarrow
\]
\[
  \text{统计力学分布律轨迹}
\]
\[
  \Longrightarrow
\]
\[
  \text{随机变量簇}
\]
\[
  \Longrightarrow
\]
\[
  \text{命中指示变量与首次命中时间}
\]
\[
  \Longrightarrow
\]
\[
  \text{命中泛函与命中指数}
\]
\[
  \Longrightarrow
\]
\[
  \text{高阶正交残差方向}
\]
\[
  \Longrightarrow
\]
\[
  \text{雅可比间隙下降}
\]
\[
  \Longrightarrow
\]
\[
  \text{工程截断确定性}
\]
\[
  \Longrightarrow
\]
\[
  \text{极限保真性}
\]
\[
  \Longrightarrow
\]
\[
  \text{PDE 求解范式替代与确定性隧穿控制}.
\]
\end{Certificate}

\begin{proof}
第 \ref{sec:tex55-main} 节给出障碍穿透的拓扑锚点：
\[
  \MicroTunnel_\rho(W)\simeq\CrossWall_\rho(W).
\]
第 \ref{sec:tex55-time-frequency-global-hit}--\ref{sec:tex55-resonance-random-variable-cluster} 节把时域遍历、频域调制和分布律索引的随机变量簇写成命中过程。
第 \ref{sec:tex55-homotopy-twist-hitting-functional}--\ref{sec:tex55-jacobian-gap-power-hitting} 节把命中过程提升为同伦轨迹函数、高阶正交残差方向和雅可比间隙下降。
第 \ref{sec:tex55-engineering-truncation-deterministic-convergence} 节说明经验测度、经验命中泛函、经验命中指数和经验雅可比间隙在有限样本与有限截断下构成工程确定性系统。
第 \ref{sec:tex55-pde-paradigm-replacement}--\ref{sec:tex55-pde-explicit-solution-demotion} 节说明高阶 PDE 被转化为低残差集合命中问题，并区分构造层的问题转化与保真层的极限恢复。
本章进一步把墙或孔洞度量化、把穿透能量化、把微观隧穿命中化，并把传统变分算子扩展到残差能量、命中指数与雅可比间隙三层。
因此八项形式价值均由前述定义、定理和证书链共同推出。
\end{proof}

\begin{corollary}[第一章的几何意义]
\label{cor:tex55-first-section-geometric-meaning}
第 \ref{sec:tex55-main} 节不是孤立的拓扑开场，而是全篇确定性隧穿控制的几何锚点。
它给出的核心范式为
\[
  \boxed{
  \text{空洞修复}
  \Rightarrow
  \text{穿墙芯线}
  \Rightarrow
  \text{同伦等价}
  \Rightarrow
  \text{障碍穿透确定化}.
  }
\]
后续章节将这一范式升级为
\[
  \boxed{
  \text{统计力学命中}
  \Rightarrow
  \text{高阶正交轨迹}
  \Rightarrow
  \text{雅可比间隙下降}
  \Rightarrow
  \text{微观隧穿可控化}.
  }
\]
\end{corollary}

\begin{proof}
由引理 \ref{lem:tex55-crosswall-spine-anchor}，穿墙芯线 \(C_k\) 是可控穿透的几何锚点。
由定理 \ref{thm:tex55-metric-controllable-tunneling}，附加度量与能量结构后，障碍穿透被转化为命中泛函优化。
由定理 \ref{thm:tex55-deterministic-controllable-microtunneling}，同伦命中过程在正交方向和雅可比间隙控制下成为确定性可控过程。
故第一章的同伦等价正是后续确定性隧穿控制的几何原型。
\end{proof}

\begin{remark}[最终压缩表述]
本章最终可压缩为四句话：
\[
  \boxed{
  \begin{gathered}
  \text{对墙或孔洞进行度量，是为了把“能否穿透”的拓扑问题}\\
  \text{提升为“如何可控穿透”的确定性优化问题。}
  \end{gathered}
  }
\]
\[
  \boxed{
  \begin{gathered}
  \text{PDE 求解从显式方程求解}\\
  \text{提升为统计力学同伦命中低残差集合的形式系统。}
  \end{gathered}
  }
\]
\[
  \boxed{
  \begin{gathered}
  \text{变分算子从函数空间上的局部极值条件，}\\
  \text{扩展为命中泛函轨迹空间上的高阶正交确定性收敛算子。}
  \end{gathered}
  }
\]
\[
  \boxed{
  \begin{gathered}
  \text{微观隧穿从概率性现象，提升为可参数化、可投影、}\\
  \text{可下降、可收敛、可保真的确定性同伦命中机制。}
  \end{gathered}
  }
\]
\end{remark}

\clearpage
 
\begin{thebibliography}{99}
\raggedright

\bibitem{RepoTex37SemanticGauge}
Gao Zheng.
\newblock \emph{G 框架正则语义规范场到自然语言语义逻辑占位规范场的高阶同伦扭曲双射、LLM 运行时降级与逻辑引擎（工作笔记：形式系统版）}.
\newblock 仓内稿件：
\code{NoteBook/notebook2/tex37_pub/gframework_semantic_gauge_field_natural_language_logic_placeholder_runtime_formal_system.tex}, 2026.

\bibitem{RepoTex51TruthAnchor}
Gao Zheng.
\newblock \emph{绝对逻辑锚点、粗暴超级对齐与 LLM 日常类幻觉的语义孔洞及边缘算子修复（工作笔记：形式系统版）}.
\newblock 仓内稿件：
\code{NoteBook/notebook3/tex51_pub/gframework_truth_anchor_superalignment_hallucination_obstruction_formal_system.tex}, 2026.

\bibitem{RepoTex17VariationalTower}
Gao Zheng.
\newblock \emph{边缘算子封闭性、比安基基准与变分生成的同伦修复塔（工作笔记：形式系统版）}.
\newblock 仓内稿件：
\code{NoteBook/notebook1/tex17_pub/gframework_edge_operator_homotopy_closure_variational_tower_formal_system.tex}, 2026.

\bibitem{RepoTex28JacGap}
Gao Zheng.
\newblock \emph{基于雅可比间隙动力驱动的同伦扭曲与统计力学最小势能路径收敛}.
\newblock 仓内稿件：
\code{NoteBook/notebook1/tex28_pub/gframework_jacgap_homotopy_distortion_statmech_min_potential_path_convergence.tex}, 2026.

\bibitem{RepoTex45ObstructionGlobalSection}
Gao Zheng.
\newblock \emph{钱学森弹道、钱学森障碍与观测--推断--障碍修复：能力元组、边缘算子与全局截面障碍（工作笔记：形式系统版）}.
\newblock 仓内稿件：
\code{NoteBook/notebook2/tex45_pub/gframework_observation_inference_obstruction_edge_operator_global_section_formal_system.tex}, 2026.

\bibitem{RepoTex50LinftyRepair}
Gao Zheng.
\newblock \emph{G 框架中性变态射力迫属性、\(L_\infty\) 无故障修复与三层证明强度（工作笔记：形式系统版）}.
\newblock 仓内稿件：
\code{NoteBook/notebook3/tex50_pub/gframework_forcing_property_morphism_linfty_failure_free_repair_formal_system.tex}, 2026.

\bibitem{RepoTex53StrategyCluster}
Gao Zheng.
\newblock \emph{统计力学参数格点、同伦扭曲命中与 PDE 策略簇：从时域求解到频域簇命中的近似量子软件层算法（工作笔记：形式系统版）}.
\newblock 仓内稿件：
\code{NoteBook/notebook3/tex53_pub/gframework_statistical_lattice_homotopy_tunnel_pde_strategy_cluster_formal_system.tex}, 2026.

\bibitem{Hatcher2002AT}
Allen Hatcher.
\newblock \emph{Algebraic Topology}.
\newblock Cambridge University Press, 2002.

\bibitem{GaoZheng2025GFramework}
Gao Zheng.
\newblock \emph{Meta-Mathematical Theory based on Pan-Logic Analysis and Pan-Iterative Analysis}.
\newblock \emph{(GaoZheng G-Framework) and the Principal-Bundle-Based Generalized}.
\newblock \emph{Noncommutative Lie Algebra (GaoZheng G-Algebra)}.
\newblock Zenodo, 2025.
\newblock \url{https://doi.org/10.5281/zenodo.17651584}.

\bibitem{LevinPeresWilmer2017Markov}
David A. Levin, Yuval Peres, and Elizabeth L. Wilmer.
\newblock \emph{Markov Chains and Mixing Times}.
\newblock American Mathematical Society, second edition, 2017.

\bibitem{Norris1998MarkovChains}
J. R. Norris.
\newblock \emph{Markov Chains}.
\newblock Cambridge University Press, 1998.

\bibitem{Kirkpatrick1983Optimization}
Scott Kirkpatrick, C. Daniel Gelatt, and Mario P. Vecchi.
\newblock Optimization by simulated annealing.
\newblock \emph{Science}, 220(4598):671--680, 1983.

\bibitem{Hajek1988Cooling}
Bruce Hajek.
\newblock Cooling schedules for optimal annealing.
\newblock \emph{Mathematics of Operations Research}, 13(2):311--329, 1988.

\bibitem{Rubinstein1999CrossEntropy}
Reuven Y. Rubinstein.
\newblock The cross-entropy method for combinatorial and continuous optimization.
\newblock \emph{Methodology and Computing in Applied Probability}, 1:127--190, 1999.

\bibitem{NocedalWright2006NumericalOptimization}
Jorge Nocedal and Stephen J. Wright.
\newblock \emph{Numerical Optimization}.
\newblock Springer, second edition, 2006.

\bibitem{VanDerVaartWellner1996EmpiricalProcesses}
Aad W. van der Vaart and Jon A. Wellner.
\newblock \emph{Weak Convergence and Empirical Processes: With Applications to Statistics}.
\newblock Springer, 1996.

\bibitem{GelfandFomin2000CalculusVariations}
I. M. Gelfand and S. V. Fomin.
\newblock \emph{Calculus of Variations}.
\newblock Dover Publications, 2000.

\bibitem{Rockafellar1970ConvexAnalysis}
R. Tyrrell Rockafellar.
\newblock \emph{Convex Analysis}.
\newblock Princeton University Press, 1970.

\bibitem{AmbrosioGigliSavare2008GradientFlows}
Luigi Ambrosio, Nicola Gigli, and Giuseppe Savar\'e.
\newblock \emph{Gradient Flows in Metric Spaces and in the Space of Probability Measures}.
\newblock Birkh\"auser, second edition, 2008.

\bibitem{KarimiNutiniSchmidt2016PL}
Hamed Karimi, Julie Nutini, and Mark Schmidt.
\newblock Linear convergence of gradient and proximal-gradient methods under the Polyak-Lojasiewicz condition.
\newblock In \emph{ECML PKDD 2016}, pages 795--811, 2016.

\bibitem{Evans2010PDE}
Lawrence C. Evans.
\newblock \emph{Partial Differential Equations}.
\newblock American Mathematical Society, second edition, 2010.

\bibitem{BrennerScott2008FEM}
Susanne C. Brenner and L. Ridgway Scott.
\newblock \emph{The Mathematical Theory of Finite Element Methods}.
\newblock Springer, third edition, 2008.

\end{thebibliography}

\end{CJK*}
\end{document}
